\documentclass[
    12pt,
    oneside,            % 1. REMOVE PÁGINAS EM BRANCO (substitui twoside e ignora openright)
    a4paper,
    brazilian
    ]{abntex2}


% ---
% PACOTES
% ---

% ---
% Pacotes fundamentais
% ---
% \usepackage{lmodern}% Usa a fonte Latin Modern
%\usepackage{mathptmx}
\usepackage[T1]{fontenc}		% Selecao de codigos de fonte.
% \usepackage{mathptmx}
\usepackage[utf8]{inputenc}		% Codificacao do documento (conversão automática dos acentos)
\usepackage{mathptmx}
% \usepackage{fancyhdr}

% \fontfamily{ptm}\selectfont

\usepackage{indentfirst}		% Indenta o primeiro parágrafo de cada seção.
\usepackage{color}				% Controle das cores
\usepackage{graphicx}			% Inclusão de gráficos
\usepackage{microtype} 			% para melhorias de justificação
\usepackage{pdfpages}
% ---

% ---
% Pacotes adicionais, usados no anexo do modelo de folha de identificação
% ---
\usepackage{multicol}
\usepackage{multirow}
% ---

% ---
% Pacotes adicionais, usados apenas no âmbito do Modelo Canônico do abnteX2
% ---
\usepackage{lipsum}				% para geração de dummy text
% ---

% ---
% Pacotes necessários para o apêndice de requisitos (quadros e equações)
% ---
\usepackage{amsmath}				% equações com \dfrac, \text e cases
\usepackage{longtable}				% quadros longos que quebram entre páginas
\usepackage{xurl}					% permite quebrar URLs longas em qualquer ponto
\usepackage[htt]{hyphenat}		% permite hifenizar nomes de entidades em \texttt
% ---

% ---
% Imagens ainda não geradas
% ---
% \imgura[altura]{largura}{arquivo}: se o arquivo já existir, inclui a imagem
% com a largura informada; caso contrário, desenha uma moldura com o nome do
% arquivo, para que o texto compile enquanto o artefato não é produzido. Ao
% criar o arquivo em img/, a imagem passa a aparecer automaticamente.
% \detokenize protege o "_" do nome do arquivo, que o LaTeX leria como subscrito.
\newcommand{\imgura}[3][5cm]{%
	\IfFileExists{#3}{%
		\includegraphics[width=#2]{#3}%
	}{%
		\begingroup
		\setlength{\fboxsep}{0pt}%
		\fbox{%
			\parbox[c][#1][c]{\dimexpr#2-2\fboxrule\relax}{%
				\centering\footnotesize\texttt{\detokenize{#3}}%
			}%
		}%
		\endgroup
	}%
}
\addto\captionsbrazilian{\renewcommand{\fontename}{Fonte}}
\renewcommand{\fontename}{Fonte}
% ---

% ---
% Pacotes de citações
% ---
\usepackage[brazilian,hyperpageref]{backref}	 % Paginas com as citações na bibl
\usepackage[alf]{abntex2cite}	% Citações padrão ABNT

% ---
% CONFIGURAÇÕES DE PACOTES
% ---
% ---
% TAMANHO NORMAL PARA TÍTULOS DE APÊNDICES E ANEXOS
% ---
\makeatletter
% Garante que a fonte de capítulos (usada por apêndices) seja normal
\renewcommand{\ABNTEXchapterfontsize}{\normalsize}
\makeatother
% ---
% Configurações do pacote backref
% Usado sem a opção hyperpageref de backref
\renewcommand{\backrefpagesname}{Citado na(s) página(s):~}
% Texto padrão antes do número das páginas
\renewcommand{\backref}{}
% Define os textos da citação
\renewcommand*{\backrefalt}[4]{
	\ifcase #1 %
		Nenhuma citação no texto.%
	\or
		Citado na página #2.%
	\else
		Citado #1 vezes nas páginas #2.%
	\fi}%
% ---
\renewcommand{\ABNTEXchapterfont}{\fontfamily{ptm}\fontseries{b}\selectfont}
\renewcommand{\imprimircapa}{%
\begin{capa}%
\center
\ABNTEXchapterfont \large \bfseries \MakeUppercase{\imprimirinstituicao}

\vspace*{\fill/2}
{\ABNTEXchapterfont \normalfont \large \imprimirautor}
\vspace*{\fill/2}
\begin{center}
\ABNTEXchapterfont\bfseries\large\imprimirtitulo
\end{center}
\vfill

\large\normalfont\imprimirlocal

\large\imprimirdata
\vspace*{1cm}
\end{capa}
}
\renewcommand{\folhaderostocontent}{\begin{center}
    \center
    \ABNTEXchapterfont \large \bfseries \MakeUppercase{\imprimirinstituicao}

    \vspace*{\fill}\vspace*{\fill}
    \begin{center}
      \ABNTEXchapterfont\bfseries\large\imprimirtitulo
    \end{center}
    \vspace*{\fill/2}


      \normalfont\normalsize
      \hspace{.45\textwidth}
      \begin{minipage}{.5\textwidth}
      	\SingleSpacing
         \imprimirpreambulo
       \end{minipage}%
       \vspace*{\fill} \vspace*{\fill}



    {\large\imprimirlocal}
    \par
    {\large\imprimirdata}
    \vspace*{1cm}

  \end{center}
}
% ---
% ALTERAR TAMANHO DA FONTE DAS SEÇÕES
% ---
\renewcommand{\ABNTEXsectionfontsize}{\normalsize}
\renewcommand{\ABNTEXsubsectionfontsize}{\normalsize}
\renewcommand{\ABNTEXsubsectionfont}{\normalfont\normalsize}
% ---
% Informações de dados para CAPA e FOLHA DE ROSTO
% ---
\titulo{Sistema Web de Análise de Segurança Territorial por Meio de Dados Públicos}
\autor{\texorpdfstring{Daniela Muniz de Souza, 23215770 \\ Pedro Henrique Felix Gomes, 24217853}{Daniela Muniz de Souza; Pedro Henrique Felix Gomes}}
\local{São Paulo - SP}
\data{2026}
\instituicao{%
  Universidade Virtual do Estado de São Paulo
  }
\tipotrabalho{Relatório técnico}
% O preambulo deve conter o tipo do trabalho, o objetivo,
% o nome da instituição e a área de concentração
\preambulo{Relatório Técnico-Científico apresentado na disciplina de Projeto Integrador para o curso de Bacharel em Tecnologia da Informação da Universidade Virtual do Estado de São Paulo (UNIVESP).}
% ---

% ---
% Configurações de aparência do PDF final

% alterando o aspecto da cor azul
\definecolor{blue}{RGB}{41,5,195}

% informações do PDF
\makeatletter
\hypersetup{
     	%pagebackref=true,
		pdftitle={\@title},
		pdfauthor={\@author},
    	pdfsubject={\imprimirpreambulo},
	    pdfcreator={LaTeX with abnTeX},
		pdfkeywords={abnt}{latex}{abntex}{abntex2}{relatório técnico},
		colorlinks=true,       		% false: boxed links; true: colored links
    	linkcolor=black,          	% color of internal links
    	citecolor=black,        		% color of links to bibliography
    	filecolor=magenta,      		% color of file links
		urlcolor=black,
		bookmarksdepth=4
}
\makeatother
% ---

% ---
% Espaçamentos entre linhas e parágrafos
% ---

% O tamanho do parágrafo é dado por:
\setlength{\parindent}{1.3cm}

% Controle do espaçamento entre um parágrafo e outro:
\setlength{\parskip}{0.2cm}  % tente também \onelineskip

% ---
% compila o indice
% ---
\makeindex
% ---
\usepackage{chngcntr}
\counterwithout{section}{chapter}
\renewcommand{\thesection}{\arabic{section}}
% ---
% REMOVER LINHA E TEXTO DO CABEÇALHO (Padrão ABNT)
% ---

% ----
% Início do documento
% ----
\begin{document}

% Seleciona o idioma do documento (conforme pacotes do babel)
%\selectlanguage{english}
\selectlanguage{brazilian}

% Retira espaço extra obsoleto entre as frases.
\frenchspacing

% ----------------------------------------------------------
% ELEMENTOS PRÉ-TEXTUAIS
% ----------------------------------------------------------
% \pretextual

% ---
% Capa
% ---
\imprimircapa
% ---
% ---
% Folha de rosto
% (o * indica que haverá a ficha bibliográfica)
% ---
\imprimirfolhaderosto*
% ---
% ---
% Anverso da folha de rosto:
% ---

% {
% \ABNTEXchapterfont

% \vspace*{\fill}

% Conforme a ABNT NBR 10719:2015, seção 4.2.1.1.1, o anverso da folha de rosto
% deve conter:

% \begin{alineas}
%   \item nome do órgão ou entidade responsável que solicitou ou gerou o
%    relatório;
%   \item título do projeto, programa ou plano que o relatório está relacionado;
%   \item título do relatório;
%   \item subtítulo, se houver, deve ser precedido de dois pontos, evidenciando a
%    sua subordinação ao título. O relatório em vários volumes deve ter um título
%    geral. Além deste, cada volume pode ter um título específico;
%   \item número do volume, se houver mais de um, deve constar em cada folha de
%    rosto a especificação do respectivo volume, em algarismo arábico;
%   \item código de identificação, se houver, recomenda-se que seja formado
%    pela sigla da instituição, indicação da categoria do relatório, data,
%    indicação do assunto e número sequencial do relatório na série;
%   \item classificação de segurança. Todos os órgãos, privados ou públicos, que
%    desenvolvam pesquisa de interesse nacional de conteúdo sigiloso, devem
%     informar a classificação adequada, conforme a legislação em vigor;
%   \item nome do autor ou autor-entidade. O título e a qualificação ou a função
%    do autor podem ser incluídos, pois servem para indicar sua autoridade no
%    assunto. Caso a instituição que solicitou o relatório seja a mesma que o
%    gerou, suprime-se o nome da instituição no campo de autoria;
%   \item local (cidade) da instituição responsável e/ou solicitante; NOTA: No
%    caso de cidades homônimas, recomenda-se o acréscimo da sigla da unidade da
%    federação.
%   \item ano de publicação, de acordo com o calendário universal (gregoriano),
%   deve ser apresentado em algarismos arábicos.
% \end{alineas}

% \vspace*{\fill}
% }

% % ---
% % Agradecimentos
% % ---
% \begin{agradecimentos}
% O agradecimento principal é direcionado a Youssef Cherem, autor do
% \nameref{formulado-identificacao} (\autopageref{formulado-identificacao}).

% Os agradecimentos especiais são direcionados ao Centro de Pesquisa em
% Arquitetura da Informação\footnote{\url{http://www.cpai.unb.br/}} da Universidade de
% Brasília (CPAI), ao grupo de usuários
% \emph{latex-br}\footnote{\url{http://groups.google.com/group/latex-br}} e aos
% novos voluntários do grupo
% \emph{\abnTeX}\footnote{\url{http://groups.google.com/group/abntex2} e
% \url{http://www.abntex.net.br/}}~que contribuíram e que ainda
% contribuirão para a evolução do abn\TeX.

% \end{agradecimentos}
% ---

% ---
% RESUMO
% ---

\noindent
SOUZA, D. M.; GOMES, P. H. F. \textbf{\imprimirtitulo{}}. Relatório Técnico-Científico. Bacharel em Tecnologia da Informação, \textbf{Universidade Virtual do Estado de São Paulo.} Tutor: Maycol Clister Fernández Chillcce. Polo Butantã, CEU Azul da Cor do Mar, 2026.
\vspace*{2.5cm}

\noindent
\textbf{RESUMO}

\vspace{0.6cm}


{\linespread{1.0}\selectfont \noindent
O presente trabalho visa contribuir com dados para a tomada de decisões complexas aos cidadãos da cidade de São Paulo, marcadas pela assimetria informativa e pela instrumentalização do medo da violência urbana. Embora dados públicos relevantes estejam formalmente disponíveis em repositórios abertos, como o portal GeoSampa e os relatórios estatísticos da Secretaria da Segurança Pública (SSP-SP), tais informações encontram-se dispersas em formatos técnicos e de baixa usabilidade, inviabilizando a capacidade de inferência dos usuários. O objetivo é desenvolver uma plataforma web cívica e acessível voltada a munícipes, que integra indicadores oficiais de segurança pública à oferta de infraestrutura urbana essencial (saúde, transporte coletivo, áreas verdes e educação) mediante a consulta simplificada por Código de Endereçamento Postal (CEP). Neste trabalho adota-se a abordagem de Design Thinking, organizada nas etapas de escuta do contexto, criação e prototipagem, e implementação e teste, associada a práticas de engenharia de software e métodos ágeis inspirados no Scrum. A escuta presencial realizada com dez moradores confirmou o desconhecimento generalizado acerca das bases estatais e orientou a especificação de um Produto Mínimo Viável (MVP) com 25 requisitos funcionais (14 essenciais), 18 requisitos não funcionais e 16 regras de negócio, estruturado em arquitetura em camadas desacopladas com retaguarda em Python, banco de dados relacional PostgreSQL/PostGIS, cliente web acessível em conformidade com as diretrizes WCAG 2.1 e WAI-ARIA, e plano de testes automatizados. Os resultados parciais incluem o diagnóstico empírico, a especificação técnica detalhada, os diagramas de casos de uso, arquitetura e entidade-relacionamento, além de wireframes conceituais e do cronograma incremental em sprints para a codificação e validação com a comunidade.
	\par}


\vspace{2cm}
\noindent
\MakeUppercase{\textbf{Palavras-chaves:}} diagnóstico socioespacial; segurança pública; dados abertos; tecnologia cívica; acessibilidade web; engenharia de software.
% ---
\newpage
% ---
% inserir lista de ilustrações
% ---
\pdfbookmark[0]{\listfigurename}{lof}
\listoffigures*
\cleardoublepage
% ---

% ---
% inserir lista de tabelas
% ---
% \pdfbookmark[0]{\listtablename}{lot}
% \listoftables*
% \cleardoublepage
% ---

% ---
% inserir lista de abreviaturas e siglas
% ---
\begin{siglas}
  \item[ABNT] Associação Brasileira de Normas Técnicas
  \item[AGPL] \textit{Affero General Public License}
  \item[AJAX] \textit{Asynchronous JavaScript and XML}
  \item[API] \textit{Application Programming Interface} (Interface de Programação de Aplicação)
  \item[CEP] Código de Endereçamento Postal
  \item[CEU] Centro Educacional Unificado
  \item[CSV] \textit{Comma-Separated Values}
  \item[DOM] \textit{Document Object Model}
  \item[DP] Distrito Policial
  \item[e-MAG] Modelo de Acessibilidade em Governo Eletrônico
  \item[ETL] \textit{Extract, Transform, Load} (Extração, Transformação e Carga)
  \item[GNU] \textit{GNU's Not Unix}
  \item[HTTPS] \textit{Hypertext Transfer Protocol Secure}
  \item[IBGE] Instituto Brasileiro de Geografia e Estatística
  \item[IPL] Índice de Proximidade Local
  \item[IVEP] Índice de Vitalidade do Espaço Público
  \item[JSON] \textit{JavaScript Object Notation}
  \item[LAI] Lei de Acesso à Informação
  \item[LGPD] Lei Geral de Proteção de Dados Pessoais
  \item[MoSCoW] \textit{Must, Should, Could, Won't}
  \item[MVP] Produto Mínimo Viável (\textit{Minimum Viable Product})
  \item[ODbL] \textit{Open Database License}
  \item[OWASP] \textit{Open Web Application Security Project}
  \item[PMSP] Prefeitura do Município de São Paulo
  \item[REST] \textit{Representational State Transfer}
  \item[RIA] Aplicações Ricas de Internet (\textit{Rich Internet Applications})
  \item[SIRGAS] Sistema de Referência Geocêntrico para as Américas
  \item[SQL] \textit{Structured Query Language}
  \item[SSP-SP] Secretaria da Segurança Pública do Estado de São Paulo
  \item[SUS] \textit{System Usability Scale}
  \item[UBS] Unidade Básica de Saúde
  \item[UNIVESP] Universidade Virtual do Estado de São Paulo
  \item[UTM] \textit{Universal Transverse Mercator} (Universal Transversa de Mercator)
  \item[UX] \textit{User Experience} (Experiência do Usuário)
  \item[WAI-ARIA] \textit{Web Accessibility Initiative - Accessible Rich Internet Applications}
  \item[WCAG] \textit{Web Content Accessibility Guidelines}
  \item[WFS] \textit{Web Feature Service}
  \item[ZAP] \textit{Zed Attack Proxy}
\end{siglas}
% ---

% ---
% inserir lista de símbolos
% ---
% \begin{simbolos}
%   \item[$ \Gamma $] Letra grega Gama
%   \item[$ \Lambda $] Lambda
%   \item[$ \zeta $] Letra grega minúscula zeta
%   \item[$ \in $] Pertence
% \end{simbolos}
% ---

% ---
% inserir o sumario
% ---
\pdfbookmark[0]{\contentsname}{toc}
\tableofcontents*
\cleardoublepage
% ---


% ----------------------------------------------------------
% ELEMENTOS TEXTUAIS
% ----------------------------------------------------------
\textual

% ----------------------------------------------------------
% Introdução (exemplo de capítulo sem numeração, mas presente no Sumário)
% ----------------------------------------------------------
% ---
% REMOVER LINHA E TEXTO DO CABEÇALHO (Padrão ABNT)
% ---
\makepagestyle{abntclean}
\makeevenhead{abntclean}{\thepage}{}{}
\makeoddhead{abntclean}{}{}{\thepage}
\makeevenfoot{abntclean}{}{}{}
\makeoddfoot{abntclean}{}{}{}

\pagestyle{abntclean}

\section{Introdução}

A busca por moradia na cidade de São Paulo impõe às famílias e cidadãos um processo decisório complexo, marcado por profundas assimetrias na distribuição de infraestrutura urbana, serviços públicos e equipamentos essenciais \cite{caldeira2000}. Dentre as variáveis a serem ponderadas na escolha de um novo lar, uma que ressalta particularmente para a comunidade escolhida como foco deste projeto integrador é a preocupação com a segurança e a habitabilidade das futuras vizinhanças de pessoas em transição habitacional na capital paulista.

A dinâmica metropolitana revela que a proliferação da violência e a instrumentalização do medo transformaram profundamente os modos de habitar, impulsionando a formação de enclaves fortificados e a privatização da vigilância, o que fragmenta o tecido público e reforça a segregação socioespacial, conforme criticado por \citeonline{caldeira2000}

No âmbito da economia urbana, a criminalidade é frequentemente tratada como uma desamenidade hedônica que deprecia o valor do solo \cite{teixeira2016}\footnotemark. No entanto, converter a segurança em um indicador isolado ou em mapas alarmistas tende a agravar preconceitos territoriais contra bairros populares e alimentar a especulação imobiliária, ignorando que crimes patrimoniais de rua frequentemente se concentram em áreas de alta renda e fluxo comercial, enquanto a segurança real depende da vitalidade comunitária, dos ``olhos na rua'' e da presença de infraestrutura urbana acolhedora \cite{jacobs2011,nev2024}.

Em meio a esse cenário, os cidadãos enfrentam uma assimetria informativa acentuada: por um lado, plataformas imobiliárias convencionais utilizam o discurso da segurança privada para supervalorizar condomínios murados e omitem carências do entorno; por outro, dados públicos relevantes como as estatísticas da Secretaria da Segurança Pública (SSP-SP) e as camadas de saúde, áreas verdes e transporte mantidas pelo portal GeoSampa permanecem fragmentados em painéis administrativos e formatos especializados, inacessíveis à população sem domínio de ferramentas de geoprocessamento \cite{ssp2024,geosampa2024,cem2021}.

Diante desse contexto, o tema norteador do projeto integrador: ``Desenvolvimento de um software com framework web que utilize noções de banco de dados, script web (Javascript), nuvem, uso de API, acessibilidade, controle de versão e testes'', foi aplicada da metodologia do \emph{Design Thinking} (nas etapas de ouvir, criar e implementar) junto a cidadãos que vivenciam a busca por habitação. Foi concebido um modelo de solução centrado na transparência cívica e no direito à cidade, uma ferramenta capaz de converter o endereço pretendido, consultado de forma simplificada por Código de Endereçamento Postal (CEP), em uma avaliação equilibrada do território, contextualizando indicadores de segurança com a proximidade a serviços essenciais (unidades de saúde, praças, parques e transporte público coletivo), desmistificando estigmas e fortalecendo o poder de escolha do morador.

A solução almejada, nomeada \textbf{GeoHabita: Diagnóstico Socioespacial Integrado}, tem como objetivo democratizar o acesso a informações urbanas qualificadas, servindo como instrumento de apoio à decisão habitacional e de cidadania. Do ponto de vista técnico, este trabalho busca construir um Produto Mínimo Viável (em inglês \emph{MVP}) estruturado em arquitetura web em camadas, utilizando JavaScript/TypeScript, banco de dados com suporte a consultas geoespaciais, microsserviços integradores de APIs abertas (ViaCEP, BrasilAPI e servidores WFS do município), infraestrutura em nuvem e conformidade primária com os padrões e diretrizes de acessibilidade para Aplicações Ricas de Internet (RIAs) analisados por \citeonline{padroesAcessibilidade2019}. Adotam-se, assim, as recomendações da WCAG 2.1 e a especificação WAI-ARIA, com destaque para \emph{roles}, estados, propriedades e regiões vivas (\emph{live regions}), garantindo que o dinamismo assíncrono das consultas e a apresentação dos dados cartográficos disponham de equivalência textual plena para sintetizadores de voz e leitores de tela.

Foram, contudo, excluídas do escopo inicial rotinas de extração automatizada não estruturada (\emph{Web Scraping}) e modelagens preditivas complexas de inteligência artificial, priorizando-se a integridade das consultas espaciais, a usabilidade universal e o prazo da disciplina. Todo o projeto será desenvolvido sob controle de versão contínuo e distribuído sob a licença GNU AGPL v3 \cite{agplv3}.

\footnotetext{
	A Teoria dos Preços Hedônicos decompõe o valor de um bem habitacional em uma cesta de atributos construtivos e amenidades/desamenidades locacionais.
	No contexto técnico nacional, a modelagem de atributos por regressão para avaliação de imóveis urbanos é padronizada pelos procedimentos de inferência estatística da norma ABNT NBR 14.653-2 \cite{abnt2011}.
}

\section{Desenvolvimento}

\hfill

\subsection{Objetivos}

Desenvolver uma plataforma web acessível e cidadã para a realização de um diagnóstico socioespacial integrado na cidade de São Paulo, articulando dados oficiais de segurança pública à disponibilidade de infraestrutura urbana essencial (como transporte público, unidades de saúde e áreas verdes).

Identificar o medo da violência urbana, e quais elementos os moradores entendem como prioritários e essenciais para a habitabilidade de suas vizinhanças (tais como caminhabilidade, iluminação e transporte seguro).

Busca-se caracterizar as bases de dados oficiais da Secretaria da Segurança Pública (SSP-SP) e do portal GeoSampa, analisando suas limitações de agregação por distrito policial e os desafios de harmonização espacial via código postal (CEP).

Pretende-se definir por meio de técnicas de engenharia de software ágil, documentar os requisitos funcionais, não funcionais e casos de uso, construindo uma especificação técnica robusta que incorpore os padrões de acessibilidade para aplicações ricas (RIAs).

Desta forma, construir, com controle de versão, rotinas de testes e banco de dados relacional com extensão geoespacial, uma plataforma web acessível, resiliente e razoavelmente escalável que possa mitigar a opacidade mercadológica e a fragmentação dos dados públicos observadas na tomada de decisão de moradia.

\hfill

\subsection{Justificativa e delimitação do problema}

À luz dos dilemas de segurança e habitabilidade apontados pelos cidadãos como pontos de maior aflição no processo de busca por moradia, e diante da constatação de que as informações públicas existentes não alcançam quem delas mais necessita, elaborou-se a seguinte pergunta norteadora: Como podemos integrar, simplificar e democratizar o acesso a dados abertos de segurança pública e infraestrutura urbana para apoiar a decisão de moradia na cidade de São Paulo, superando a ineficácia da transparência governamental e a assimetria do mercado imobiliário?

Como destaca \citeonline{caldeira2000}, a proliferação da violência e a instrumentalização da ``fala do crime'' na capital paulista impulsionaram historicamente a multiplicação de enclaves fortificados e a privatização da vigilância, fragmentando o tecido público e aprofundando a segregação socioespacial. Embora a criminalidade atue como desamenidade direta sobre os valores habitacionais sob a ótica da economia urbana \cite{teixeira2016}, o mercado imobiliário frequentemente apropria-se do medo para justificar a supervalorização de condomínios murados e a depreciação de bairros populares consolidados, desconsiderando que a segurança real decorre da vitalidade comunitária, dos ``olhos na rua'' e da presença de serviços e infraestrutura pública acolhedora \cite[cap.~1]{jacobs2011}.

Sem alternativas transparentes de consulta, as famílias permanecem vulneráveis ao discurso comercial das incorporadoras e ao sensacionalismo da mídia, perpetuando decisões de moradia desinformadas e estigmas territoriais prejudiciais à cidade \cite{nev2024}.

Durante a etapa de imersão do \emph{Design Thinking}, as entrevistas conduzidas com pessoas que buscam moradia na capital evidenciaram de forma contundente essa assimetria: nenhum dos entrevistados tinha conhecimento prévio da existência do portal GeoSampa ou dos boletins estatísticos periódicos da Secretaria da Segurança Pública (SSP-SP) \cite{ssp2024,geosampa2024}. Ao serem expressamente indagados pelo grupo se sabiam que poderiam checar os índices criminais e a distribuição de equipamentos públicos de suas futuras vizinhanças por meio de dados abertos do governo, a resposta foi unanimemente negativa. Tal constatação desvela a ineficácia da transparência pública estritamente formal, que disponibiliza microdados brutos em formatos burocráticos e painéis especializados, mas falha em convertê-los em transparência real para o cidadão comum.

Convém destacar que os dados provenientes dos munícipes ouvidos passaram por processo de \textbf{pseudonimização} neste relatório, prerrogativa fundamentada pelo caráter sensível dos relatos, que incluíam expor rotinas de deslocamento, receios de vitimização e preocupações efetivas acerca da segurança patrimonial e familiar. Os canais de contato disponibilizados foram mantidos em base segregada e sob confidencialidade rigorosa por parte da equipe responsável pelo desenvolvimento, em consonância com as diretrizes de consentimento estabelecidas pela LGPD, visando exclusivamente possibilitar o recontato com os participantes nas fases subsequentes de validação dos wireframes e do MVP.

O projeto, portanto, se propõe a atuar justamente nessa lacuna entre a disponibilidade do dado bruto e a sua utilidade social, fornecendo uma ferramenta web pública e acessível de diagnóstico socioespacial capaz de traduzir indicadores de segurança e infraestrutura em avaliações compreensíveis por meio da simples consulta por CEP, especificamente no contexto territorial da cidade de São Paulo.

\hfill

\subsection{Fundamentação Teórica}
Este trabalho sustenta-se em três pilares principais complementares: a dinâmica socioespacial urbana e os determinantes de segurança e habitabilidade na escolha residencial; a transparência da gestão pública, dados abertos e acessibilidade digital; e, por fim, a engenharia de software (métodos ágeis, arquitetura web em camadas, banco de dados com suporte geoespacial, controle de versão e garantia de qualidade). Sob estes prismas, compreende-se em profundidade o problema diagnosticado e motivador do projeto, a opacidade informativa do mercado imobiliário e a ineficácia prática da transparência governamental acerca da segurança e da infraestrutura pública em São Paulo, a estrutura conceitual da solução proposta, assim como as escolhas para o seu desenvolvimento, tanto do ponto de vista da \emph{stack} tecnológica e da modelagem espacial quanto das estratégias de teste e acessibilidade universal.

\subsubsection{Dinâmica socioespacial, a retórica do medo e a vitalidade urbana na escolha habitacional}
A conformação do tecido urbano na cidade de São Paulo é historicamente marcada por profundas desigualdades na oferta de equipamentos públicos e por padrões persistentes de segregação socioespacial. Como analisa criticamente \citeonline{caldeira2000}, a partir das últimas décadas do século XX, a disseminação da violência e a instrumentalização generalizada da ``fala do crime'' alteraram estruturalmente o padrão de urbanização da metrópole, impulsionando a privatização da segurança cotidiana e a proliferação dos chamados enclaves fortificados. A autora sintetiza o impacto dessa transformação na morfologia urbana e na convivência democrática:

\begin{citacao}
	Os condomínios fechados são a versão residencial de uma categoria mais ampla
	de novos empreendimentos urbanos que chamo de enclaves fortificados. Eles estão
	mudando consideravelmente a maneira como as pessoas das classes média e alta
	vivem, consomem, trabalham e gastam seu tempo de lazer. Eles estão mudando o
	panorama da cidade, seu padrão de segregação espacial e o caráter do espaço público
	e das interações públicas entre as classes. Os enclaves fortificados incluema
	conjuntos de escritórios, shopping centers, e cada vez mais outros espaços que têm
	sido adaptados para se conformarem a esse modelo, como escolas, hospitais, centros
	de lazer e parques temáticos. Todos os tipos de enclaves fortificados partilham
	algumas características básicas. São propriedade privada para uso coletivo e enfatizam
	o valor do que é privado e restrito ao mesmo tempo que desvalorizam o
	que é público e aberto na cidade. São fisicamente demarcados e isolados por muros,
	grades, espaços vazios e detalhes arquitetônicos. São voltados para o interior e
	não em direção à rua, cuja vida pública rejeitam explicitamente. São controlados
	por guardas armados e sistemas de segurança, que impõem as regras de inclusão e
	exclusão.[...] % São flexíveis: devido ao seu tamanho, às novas tecnologias de comunicação,
	% organização do trabalho e aos sistemas de segurança, eles são espaços autônomos,
	% independentes do seu entorno, que podem ser situados praticamente em
	% qualquer lugar. Em outras palavras, em contraste com formas anteriores de empreendimentos
	% comerciais e residenciais, eles pertencem não ao seus arredores imediatos,
	% mas a redes invisíveis (Cenzatti e Crawford 1998).! Em consequência, embora tendam
	% a ser espaços para as classes altas, podem ser situados em áreas rurais ou na
	% periferia, ao lado de favelas ou casas autoconstruídas. Finalmente, os enclaves tendem
	% a ser ambientes socialmente homogêneos. Aqueles que escolhem habitar esses
	% espaços valorizam viver entre pessoas seletas (ou seja, do mesmo grupo social) e
	% longe das interações indesejadas, movimento, heterogeneidade, perigo e imprevisibilidade
	% das ruas. Os enclaves privados e fortificados cultivam um relacionamento
	% de negação e ruptura com o resto da cidade e com o que pode ser chamado de
	% um estilo moderno de espaço público aberto à livre circulação. Eles estão transformando
	% a natureza do espaço público e a qualidade das interações públicas na cidade,
	% que estão se tornando cada vez mais marcadas por suspeita e restrição.
	\cite[p.~258-259]{caldeira2000}
\end{citacao}

No campo da economia urbana, a relação entre criminalidade e habitação é formalizada pela Teoria dos Preços Hedônicos, na qual o valor de um imóvel resulta do somatório implícito de seus atributos intrínsecos e das amenidades locacionais do entorno \cite{teixeira2016}. Sob essa perspectiva, a incidência delituosa recorrente atua como desamenidade exógena, impondo descontos mensuráveis sobre os valores residenciais. Contudo, na prática do mercado paulistano, observa-se uma apropriação predatória desse fenômeno: plataformas comerciais utilizam o medo difuso para supervalorizar condomínios murados, retroalimentando o modelo dos enclaves segregadores, e desvalorizar bairros populares periféricos com base em estigmas, ocultando que delitos patrimoniais de rua concentram-se com frequência em áreas nobres e corredores de grande circulação econômica \cite{nev2024}.


Em contraponto à lógica das barreiras físicas e da vigilância privada, \citeonline[cap.~1]{jacobs2011} demonstra que a segurança efetiva do espaço urbano não decorre do isolamento ou do policiamento ostensivo, mas sim da vitalidade pública e do controle voluntário exercido pelos próprios moradores, mecanismo sintetizado no conceito clássico dos ``olhos na rua''. Para a autora, calçadas ativas, fachadas voltadas para o logradouro e a presença equilibrada de comércios, áreas verdes e transporte coletivo atraem fluxos contínuos de pedestres em horários alternados, inibindo oportunidades criminais e construindo redes espontâneas de acolhimento comunitário.

A articulação entre \citeonline{caldeira2000}, \citeonline{teixeira2016} e \citeonline{jacobs2011} fundamenta a necessidade de uma ferramenta que rompa com o binarismo estigmatizante dos ``mapas de perigo''. Para viabilizar uma tomada de decisão habitacional verdadeiramente emancipatória, a análise de segurança precisa ser contextualizada conjuntamente com a oferta de infraestrutura urbana atenuante, equilibrando índices de criminalidade e vitalidade comunitária.

\subsubsection{Transparência pública, dados abertos e acessibilidade em aplicações ricas}
A concepção de transparência governamental que alicerça o GeoHabita ultrapassa o mero depósito passivo de dados em páginas institucionais. \citeonline{zuccolotto2019}, ao sistematizarem o debate conceitual da transparência na administração pública brasileira, diferenciam a transparência ativa, caracterizada pela divulgação proativa e espontânea de informações pelo poder público, da transparência passiva, que depende de demandas e pedidos formais de acesso provocados pelo munícipe com base na Lei de Acesso à Informação (Lei n\textordmasculine{} 12.527/2011, LAI). Para além da iniciativa, os autores enfatizam que a transparência efetiva exige a conjugação de duas dimensões estruturais: a visibilidade (a disponibilidade efetiva da informação ao escrutínio público) e a capacidade de inferência, que consiste na aptidão do cidadão de processar, compreender e extrair conclusões práticas a partir dos dados disponibilizados.


Essa distinção é central para o diagnóstico deste projeto: os portais da Secretaria da Segurança Pública de São Paulo (SSP-SP) e o mapa digital GeoSampa atendem, no plano formal, ao requisito de visibilidade da transparência ativa \cite{ssp2024,geosampa2024}. Todavia, a distribuição desses registros em arquivos brutos agregados por Circunscrição Policial ou em camadas geoespaciais herméticas anula a capacidade de inferência do cidadão comum, que desconhece essas ferramentas ou é incapaz de cruzá-las para aferir a qualidade de sua vizinhança.

Essa disfunção alinha-se às constatações empíricas de \citeonline{michener2018}, que apontam que a transparência estatal brasileira frequentemente frustra o exercício da cidadania devido à fragmentação, ao vocabulário técnico burocrático e à baixa usabilidade das plataformas mantidas pelo Estado.

Quando a administração pública falha em assegurar inferência direta à sociedade, legitima-se a necessidade do que a literatura reconhece como transparência compensatória e tecnologia cívica, em que iniciativas autônomas sistematizam, limpam e devolvem dados públicos de forma inteligível e acionável. Exemplo dessa atuação no cenário nacional é a Operação Serenata de Amor, que utilizou ciência de dados aberta para viabilizar o controle social sobre gastos públicos parlamentares \cite{dalpozzo2020}. O GeoHabita assume função análoga no tecido urbano: converte microdados públicos dispersos em um diagnóstico socioespacial de fácil apreensão para o cidadão em processo de escolha habitacional.

Contudo, a transparência não pode ser plena se mantiver barreiras excludentes de acesso tecnológico. No contexto do desenvolvimento web moderno, \citeonline{padroesAcessibilidade2019} demonstram que as Aplicações Ricas de Internet (RIAs) caracterizadas pela intensa manipulação do DOM em JavaScript e pela troca de mensagens assíncronas via AJAX/Fetch representaram avanços visuais notáveis, mas frequentemente resultaram em um ``apagão'' informativo para pessoas com deficiência visual. Como sistematizado pelos autores, mapas dinâmicos e atualizações parciais de tela deixam de ser capturados por sintetizadores de voz e leitores de tela (como NVDA, JAWS e Orca) na ausência das diretrizes de acessibilidade da WCAG 2.1 e das marcações semânticas da especificação WAI-ARIA (com destaque para \emph{roles}, estados, propriedades e regiões vivas). A garantia da acessibilidade digital erige-se, portanto, como condição estrutural de conformidade à Lei Brasileira de Inclusão e como dever ético da tecnologia cívica.

\subsubsection{Engenharia de software, modelagem geoespacial e desenvolvimento web}
A viabilidade técnica e a sustentabilidade operacional da plataforma apoiam-se em referências consagradas da engenharia de software. \citeonline{pressman2021} conceituam a engenharia de software como uma abordagem sistemática, disciplinada e quantificável aplicada ao desenvolvimento, à operação e à manutenção de sistemas computacionais. Os autores organizam a disciplina em dimensões essenciais: processo, modelagem, qualidade e segurança, e projeto de experiência do usuário (UX), destacando a convergência entre métodos ágeis e práticas rigorosas de engenharia. Essa visão norteia o GeoHabita desde a especificação de requisitos e elaboração de casos de uso até a modelagem estrutural e as rotinas de testes automatizados, estruturando o desenvolvimento em ciclos incrementais adequados ao horizonte temporal de um Projeto Integrador. A ênfase na experiência do usuário é diretamente aplicada na construção de uma interface inclusiva, fluida e descomplicada para indivíduos em busca de habitação.

\citeonline{sommerville2018}, por sua vez, dedica atenção especial à engenharia de requisitos, às práticas do desenvolvimento ágil com Scrum e ao projeto arquitetural de sistemas web distribuídos. O autor enfatiza a separação de responsabilidades em camadas desacopladas, separando apresentação, serviços de aplicação e persistência, a concepção de componentes reutilizáveis e o planejamento para evolução do código após o ciclo inicial de implantação. Essas diretrizes são basilares para a concepção do Produto Mínimo Viável (MVP) do GeoHabita, cuja arquitetura articula um cliente web reativo em JavaScript, um serviço de retaguarda estruturado em Node.js como camada integradora de APIs externas (ViaCEP, BrasilAPI e servidores municipais) e um banco de dados relacional dotado de extensão geoespacial (PostGIS).

A adoção do Scrum, conforme delineada por Sommerville, mostra-se ideal para a dinâmica da equipe, permitindo que sprints curtas produzam incrementos funcionais testáveis em ritmo constante. A complementaridade entre as abordagens de \citeonline{pressman2021} (com foco em processo, qualidade e UX) e \citeonline{sommerville2018} (com ênfase em requisitos, arquitetura distribuída e modularidade) fornece o arcabouço metodológico necessário para construir uma plataforma web tecnicamente robusta, acessível e orientada ao interesse público.

\subsection{Metodologia}

O desenvolvimento do GeoHabita adota a abordagem do \emph{Design Thinking}, metodologia ancorada no ser humano que estrutura a solução de problemas complexos em iterações cíclicas de empatia, definição, ideação, prototipação e teste \cite{brown2010}.

No âmbito deste Projeto Integrador, o trabalho organiza-se em torno das três fases fundamentais sugeridas pela Univesp: escuta e interpretação do contexto, criação e prototipagem, e execução e validação. Por se tratar de um relatório parcial, o projeto está concentrado no diagnóstico e no planejamento conceitual: realizaram-se a imersão de campo e a formalização da modelagem do sistema (especificação de requisitos, arquitetura em camadas e diagramas conceituais), restando à etapa subsequente a codificação plena, a ingestão definitiva do banco e os testes de homologação.

\textbf{Ouvir e interpretar o contexto.} A coleta de dados empíricos desenvolveu-se durante setembro,
mediante abordagens presenciais que aplicaram roteiro semiestruturado (reproduzido no Apêndice \ref{apendice:questionario}). Para garantir representatividade da mobilidade metropolitana, as ações concentraram-se em nós estratégicos de transporte intermodal e de grande fluxo de pedestres da capital: Terminal Pinheiros, Terminal Palmeiras-Barra Funda e praças de grande fluxo. O formulário incluía uma pergunta-filtro eliminatória inicial, permitindo prosseguimento apenas com transeuntes que confirmaram interesse ativo em imóvel (compra ou locação) em São Paulo nos últimos doze meses. Considerando o caráter sensível dos temas tratados — incluindo retorno noturno, questões de segurança patrimonial e dinâmicas familiares —, o instrumento instituiu a **pseudonimização** para as contribuições: as respostas foram dissociadas de qualquer dado identificativo no texto deste trabalho (atribuindo códigos anônimos de Munícipe 1 a 10), mantidos sob posse limitada e temporária da equipe exclusivamente para permitir as sessões de validação e testes de usabilidade planejados na fase de implementação. Para análise das respostas, empregou-se a técnica de análise de conteúdo categorial \cite{bardin2016}, que possibilitou extrair núcleos de sentido e classificar sistematicamente as preocupações relatadas em eixos temáticos.

\textbf{Criar e prototipar.} O processo de ideação converteu as categorias obtidas da escuta em
requisitos de engenharia de software, alinhado às práticas documentadas por \citeonline{pressman2021} e \citeonline{sommerville2018}. A especificação técnica abrangeu o inventário de requisitos funcionais e não funcionais hierarquizados pela matriz MoSCoW, a modelagem de casos de uso e o projeto da estrutura relacional do banco de dados com extensões espaciais. Pautado pelo princípio de \emph{Accessibility by Design}, reduzindo as barreiras das interfaces dinâmicas apontadas por
\citeonline{padroesAcessibilidade2019}, garantiu-se a equivalência textual dos mapas por meio de
tabelas semânticas e markup WAI-ARIA desde o início da concepção dos wireframes.

\textbf{Implementar e testar.} A concretização da ferramenta no ciclo seguinte utilizará um
ecossistema híbrido voltado para desempenho e segregação de responsabilidades: serviço de retaguarda em Python (coordenando rotas, acesso a APIs e operações analíticas), persistência espacial em SQL (PostgreSQL/PostGIS) e scripts front-end em JavaScript dedicados à interatividade de formulários, consumo assíncrono via \texttt{fetch}, controle de regiões ativas (\texttt{aria-live}) e renderização cartográfica simplificada. A garantia de qualidade contemplará testes unitários de regras de negócio, auditorias automáticas de acessibilidade e validação de usabilidade em duas etapas (confirmação conceitual dos wireframes e, em seguida, homologação do protótipo funcional com usuários reais e tecnologias assistivas).

A condução das atividades ocorre pela equipe sob metodologia ágil inspirada no
\emph{Scrum} \cite{sommerville2018}, estruturada em sprints quinzenais. O controle de versões é
administrado via repositório centralizado no GitHub, disponível em:
\url{https://github.com/DanielaMunizSouza/pi2-indicadores-seguranca-publica}.

\subsection{Resultados}
Os resultados são expostos conforme as três fases da metodologia, acompanhando a construção da solução desde a escuta dos moradores até a especificação e prototipagem da plataforma. Por se tratar de relatório parcial, os resultados focam no diagnóstico empírico e na modelagem do sistema; os artefatos aqui descritos são protótipos conceituais que orientarão a etapa executiva de codificação.

\subsubsection{Ouvir: diagnóstico junto à comunidade}
A realização do questionário em formato presencial nos terminais e praças envolveu dez munícipes com perfis diversificados em relação à renda, faixa etária e interesse habitacional (Quadro sintético e transcrições no Apêndice \ref{apendice:questionario}). No que tange à sensibilidade inerente ao tema, 100\% dos participantes optaram explicitamente pela \textbf{preservação de sua identidade pública por meio de pseudonimização} (identificados como Munícipe 1 a Munícipe 10), conferindo autorização expressa para a manutenção de contato confidencial destinada à devolutiva e avaliação das interfaces desenhadas. A análise das respostas evidenciou cenário de forte assimetria informacional e decisões tomadas com base em informação inadequada: todos os dez participantes relataram fundamentar suas escolhas apenas em comentários informais de conhecidos ou no material publicitário das imobiliárias. Quando questionados diretamente sobre conhecimento do mapa digital GeoSampa ou das planilhas estatísticas mensais da SSP-SP, a resposta foi unânime (dez respostas ``Não''). O achado validou empiricamente a disfunção conceituada por \citeonline{zuccolotto2019}: o Estado paulista mantém formalmente a transparência ativa visível, mas não oferece a capacidade de inferência necessária ao cidadão comum. A aplicação da análise categorial \cite{bardin2016} agrupou as queixas dos entrevistados em três eixos:



\begin{enumerate}
  \item \textbf{Iluminação pública e fluxo de pedestres} nos percursos a pé;
  \item \textbf{Acessibilidade e proximidade física a eixos estruturados de transporte coletivo};
  \item \textbf{Falta de dados criminais auditáveis} em escala de vizinhança.
\end{enumerate}

Notavelmente, as duas primeiras preocupações referem-se à infraestrutura urbana cotidiana, não a estatísticas criminais isoladas. Interpretado sob a ótica da teoria de vitalidade urbana de \citeonline[cap.~2]{jacobs2011}, esse resultado motivou reorientação profunda no escopo do projeto: a segurança pública passou a ser modelada como reflexo direto da ocupação do espaço público e da rede comunitária gerada por praças, parques, feiras livres, escolas, unidades de saúde e nós de transporte (os ``olhos na rua''). O desafio do GeoHabita consolidou-se em transformar o CEP em síntese equilibrada entre oferta de serviços urbanos essenciais e dados oficiais de segurança, superando a lógica estigmatizante dos ``mapas de calor e perigo'' criticada por \citeonline{caldeira2000}.

\subsubsection{Criar: especificação e arquitetura da solução}

Na ideação, foram expressamente descartadas rotinas de extração não estruturada (\emph{web scraping}) e modelagens preditivas criminais, por inviabilidade temporal do projeto e pelo risco de reforçar estigmas socioespaciais \cite{caldeira2000}. Para responder aos três eixos do diagnóstico sem fundir os dados em uma nota única artificial, definiu-se abordagem analítica paralela:

\begin{itemize}
  \item \textbf{Índice de Proximidade Local (IPL):} mede a distância métrica e o tempo a pé até o equipamento público mais próximo de cada categoria (saúde, transporte, lazer);
  \item \textbf{Índice de Vitalidade do Espaço Público (IVEP):} compara a densidade global de infraestrutura urbana entre os 96 distritos paulistanos;
  \item \textbf{Painel de Ocorrências Oficiais:} exibe taxas padronizadas por 100 mil habitantes da circunscrição policial correspondente, acompanhadas da data de referência e notas técnicas sobre subnotificação.
\end{itemize}

A especificação do MVP, elaborada segundo \citeonline{pressman2021} e \citeonline{sommerville2018}, resultou em 25 requisitos funcionais, 18 requisitos não funcionais e 16 regras de negócio. Os requisitos funcionais foram hierarquizados pelo método MoSCoW, sendo 14 classificados como essenciais, organizando-se em cinco módulos: consulta de endereço (RF01 a RF06), diagnóstico socioespacial (RF07 a RF12 e RF23 a RF25), visualização acessível (RF13 a RF15), comparação e compartilhamento (RF16 a RF18) e administração dos dados (RF19 a RF22). Os requisitos não funcionais cobrem desempenho, conformidade WCAG 2.1 nível AA, e-MAG, segurança, LGPD e Lei Brasileira de Inclusão (detalhados no Apêndice~\ref{ap:requisitos}).

O sistema prevê três perfis de interação, modelados no diagrama de casos de uso (Figura~\ref{fig:casos-de-uso}): o Munícipe (consultas públicas e ajustes de preferência), o Usuário de tecnologia assistiva (navegação semântica orientada por leitor de tela ou teclado) e o Administrador de dados (gestão das rotinas ETL de atualização).

\begin{figure}[!h]
	\caption{Diagrama de casos de uso do MVP do GeoHabita}
	\label{fig:casos-de-uso}
	\centering
	\imgura{\textwidth}{img/casos_uso.png}
	\fonte{Elaborado pelo autor (2026).}
\end{figure}

\begin{sloppypar}
O banco de dados relacional foi modelado com oito entidades no diagrama de entidade-relacionamento (Figura~\ref{fig:mer}). A entidade \texttt{Distrito} armazena os 96 polígonos municipais do GeoSampa e a população censitária. \texttt{Categoria} e \texttt{Tipo de Equipamento} estruturam as amenidades urbanas e fixam os parâmetros de distância de caminhabilidade do IPL. \texttt{Equipamento} guarda a coordenada geográfica pontual de cada serviço público, e \texttt{Ocorrência} registra os dados oficiais da SSP-SP. As entidades \texttt{Indicador do Distrito} e \texttt{Carga ETL} asseguram rastreabilidade temporal e auditoria das séries históricas, enquanto \texttt{Cache de CEP} otimiza as buscas geográficas reduzindo requisições externas.
\end{sloppypar}

\begin{figure}[!h]
	\caption{Diagrama de entidade-relacionamento do banco de dados}
	\label{fig:mer}
	\centering
	\imgura{\textwidth}{img/mer.png}
	\fonte{Elaborado pelo autor (2026).}
\end{figure}

A arquitetura de software estrutura-se em camadas desacopladas (Figura~\ref{fig:arquitetura}), seguindo \citeonline{sommerville2018}. A camada de apresentação web (HTML5 semântico com folhas de estilo responsivas, WAI-ARIA e JavaScript para interatividade acessível) comunica-se via JSON com o \emph{backend} em Python. Na aplicação servidora, a camada de aplicação orquestra as regras de negócio e o cálculo espacial dos indicadores, enquanto a infraestrutura acessa as consultas PostGIS em SQL (\texttt{ST\_DWithin} e \texttt{ST\_Contains}) e gerencia os serviços de contingência entre ViaCEP e BrasilAPI.

\begin{figure}[!h]
	\caption{Arquitetura em camadas do GeoHabita}
	\label{fig:arquitetura}
	\centering
	\imgura{\textwidth}{img/arquitetura.png}
	\fonte{Elaborado pelo autor (2026).}
\end{figure}

O fluxo de uso do cidadão, abrangendo desde a inserção do CEP até a exibição dos índices e o tratamento de eventuais erros de digitação, está esquematizado na Figura~\ref{fig:fluxo}.

\begin{figure}[!h]
	\caption{Fluxo de consulta do munícipe na plataforma}
	\label{fig:fluxo}
	\centering
	\imgura[7cm]{0.6\textwidth}{img/fluxo_usuario.png}
	\fonte{Elaborado pelo autor (2026).}
\end{figure}

Por fim, consolidaram-se os \emph{wireframes} de baixa fidelidade para validar a interface. A tela inicial (Figura~\ref{fig:wf-inicial}) concentra a consulta por CEP e controles fixos de acessibilidade (alto contraste, atalhos do teclado e redimensionamento). A tela de diagnóstico (Figura~\ref{fig:wf-diagnostico}) organiza os dados em abas (Proximidade, Vitalidade e Ocorrências) e traz a aba obrigatória ``Ver como tabela'', assegurando a equivalência textual de mapas recomendada por \citeonline{padroesAcessibilidade2019}.

\begin{figure}[!h]
	\caption{\emph{Wireframe} da tela inicial de consulta}
	\label{fig:wf-inicial}
	\centering
	\imgura{0.85\textwidth}{img/wireframe_inicial.png}
	\fonte{Elaborado pelo autor (2026).}
\end{figure}

\begin{figure}[!h]
	\caption{\emph{Wireframe} da tela de diagnóstico do endereço}
	\label{fig:wf-diagnostico}
	\centering
	\imgura{\textwidth}{img/wireframe_diagnostico.png}
	\fonte{Elaborado pelo autor (2026). Dados ilustrativos.}
\end{figure}
\subsubsection{Implementar: planejamento da construção e da validação}
A codificação do MVP ainda não foi iniciada. Nesta etapa preliminar, estruturaram-se o repositório de código com controle de versão e o planejamento incremental do desenvolvimento, dividido em uma sprint inicial de ambiente e quatro sprints quinzenais de entrega funcional: (1) Consulta de CEP com resolução de coordenadas e contingência; (2) Ingestão espacial do GeoSampa e cálculo do IPL; (3) Carga dos dados da SSP-SP, cálculo do IVEP e geração das tabelas semânticas; e (4) Ajustes de preferência, painel cívico e modo de alto contraste. Os casos de teste foram derivados previamente dos requisitos e regras de negócio, viabilizando sua escrita concomitante à codificação. A validação em campo ocorrerá em duas etapas sucessivas:
\begin{enumerate}
	\item \textbf{Validação Conceitual dos Wireframes:} apresentação dos protótipos de baixa fidelidade aos dez entrevistados, verificando se a disposição das informações atende às expectativas e se os termos utilizados são compreensíveis;
	\item \textbf{Testes de Usabilidade e Acessibilidade do MVP Funcional:} execução de roteiros de tarefas com o sistema publicado em nuvem, incluindo testes práticos de navegação por teclado e por leitores de tela (NVDA/Orca) para certificar a usabilidade universal antes da entrega final.
\end{enumerate}


\phantompart

% ---
% Conclusão
% ---
% \chapter{Conclusão}
% ---

% \lipsum[31-33]

% ----------------------------------------------------------
% ELEMENTOS PÓS-TEXTUAIS
% ----------------------------------------------------------
\postextual

% ----------------------------------------------------------
% Referências bibliográficas
% ----------------------------------------------------------
\bibliography{references}

% ----------------------------------------------------------
% Glossário
% ----------------------------------------------------------
%
% Consulte o manual da classe abntex2 para orientações sobre o glossário.
%
%\glossary

% ----------------------------------------------------------
% Apêndices
% ----------------------------------------------------------

% ---
% Inicia os apêndices
% ---
\begin{apendicesenv}

	% Imprime uma página indicando o início dos apêndices
	% \partapendices

	% ----------------------------------------------------------
	\chapter{Especificação de requisitos de software}
	\label{ap:requisitos}
	% ----------------------------------------------------------
	% =====================================================================
% GeoHabita --- Especificação de Requisitos de Software
% Trecho para inclusão no relatório (abnTeX2). Sugestão de uso:
%   % =====================================================================
% GeoHabita --- Especificação de Requisitos de Software
% Trecho para inclusão no relatório (abnTeX2). Sugestão de uso:
%   % =====================================================================
% GeoHabita --- Especificação de Requisitos de Software
% Trecho para inclusão no relatório (abnTeX2). Sugestão de uso:
%   \input{requisitos_geohabita}
% posicionado após a subseção "Metodologia" (ou dentro de
% "Resultados preliminares: solução inicial").
%
% Pacotes necessários no preâmbulo do relatório:
%   \usepackage{longtable}
%   \usepackage{tabularx}
%   \usepackage{booktabs}
%   \usepackage{array}
%   \usepackage{amsmath}
% Referências novas: ver requisitos_geohabita.bib
% =====================================================================

\counterwithin{section}{chapter}
\renewcommand{\thesection}{\thechapter.\arabic{section}}
\counterwithin{table}{chapter}
\renewcommand{\thetable}{\thechapter.\arabic{table}}
\counterwithin{equation}{chapter}
\renewcommand{\theequation}{\thechapter.\arabic{equation}}

\newcolumntype{L}[1]{>{\raggedright\arraybackslash}p{#1}}
\newcolumntype{C}[1]{>{\centering\arraybackslash}p{#1}}


\section{Especificação de requisitos}

A especificação de requisitos do GeoHabita consolida, em termos verificáveis, aquilo que a plataforma deve oferecer ao cidadão e as condições de qualidade sob as quais deve operar. Em relação à concepção inicial, o escopo foi reorientado: a segurança deixa de ser tratada prioritariamente a partir da estatística criminal e passa a ser lida também pela infraestrutura pública de uso do espaço, isto é, pelos equipamentos e lugares que atraem pessoas às ruas em horários variados e sustentam o que \citeonline[cap.~1]{jacobs2011} denomina vigilância natural dos ``olhos na rua''. Praças, parques, equipamentos culturais e esportivos, feiras livres, pontos de transporte coletivo, unidades de saúde e escolas passam a compor um indicador próprio de vitalidade do espaço público, apresentado lado a lado com os registros oficiais da SSP-SP \cite{ssp2024,geosampa2024}.

A estrutura adotada segue a norma ISO/IEC/IEEE 29148 \cite{iso29148} e a classificação proposta por \citeonline{sommerville2018}, que distingue requisitos funcionais, requisitos não funcionais (de produto, organizacionais e externos) e requisitos de domínio. Os atributos de qualidade são organizados segundo o modelo da ISO/IEC 25010 \cite{iso25010}, e o desdobramento em histórias de usuário e critérios de aceitação acompanha as práticas ágeis descritas por \citeonline{cohn2004} e \citeonline{pressman2021}. Cada requisito recebe identificador único, prioridade segundo a técnica MoSCoW (\emph{Must}, \emph{Should}, \emph{Could}, \emph{Won't}) e rastreabilidade até os casos de teste que o verificam.

\subsection{Escopo do produto e atores}

O Produto Mínimo Viável (MVP) delimita-se ao município de São Paulo e aos seus 96 distritos administrativos, unidade territorial escolhida por dispor de polígonos oficiais no GeoSampa e de população recenseada pelo IBGE, o que permite o cálculo de taxas comparáveis entre regiões. A consulta parte de um CEP ou de um endereço informado pelo usuário, que é convertido em coordenadas geográficas e associado ao distrito correspondente. A partir desse ponto, o sistema produz um diagnóstico socioespacial composto por três blocos: a proximidade do endereço consultado aos equipamentos públicos de uso do espaço, medida pela distância até o equipamento mais próximo de cada tipo; o indicador de vitalidade do espaço público do distrito; e os indicadores oficiais de ocorrências criminais, sempre contextualizados.

Permanecem fora do escopo do MVP a intermediação ou divulgação de anúncios imobiliários, a coleta colaborativa de relatos de crimes, a predição de ocorrências por modelos de aprendizado de máquina, a extração automatizada não estruturada de páginas (\emph{web scraping}) e aplicativos nativos para dispositivos móveis. A exclusão dos relatos colaborativos decorre, em particular, do risco de reforço de estigmas territoriais e de reidentificação de vítimas já discutido na fundamentação teórica.

Identificam-se três atores principais, descritos no Quadro~\ref{qua:atores}.

\begin{table}[htb]
	\caption{Atores do sistema}
	\label{qua:atores}
	\centering
	\small
	\begin{tabularx}{\textwidth}{L{3.6cm}X}
		\toprule
		\textbf{Ator}                   & \textbf{Descrição e interesse}                                                                                                                                                            \\
		\midrule
		Munícipe consulente             & Pessoa ou família em transição habitacional, ou morador interessado em conhecer o próprio território. Consulta endereços, compara regiões e não precisa de cadastro para o uso essencial. \\
		Usuário de tecnologia assistiva & Munícipe que acessa a plataforma por leitor de tela, navegação exclusiva por teclado ou ampliação. Exige equivalência textual integral para mapas e atualizações dinâmicas.               \\
		Administrador de dados          & Integrante da equipe responsável por executar as rotinas de carga (ETL), validar as bases importadas e publicar novas versões dos indicadores.                                            \\
		Sistemas externos               & Serviços de CEP e geocodificação, servidores WFS do GeoSampa, portais de dados abertos da SSP-SP e do IBGE. Impõem limites de uso, licenças e obrigações de atribuição.                   \\
		\bottomrule
	\end{tabularx}
	\fonte{Elaborado pelo autor (2026).}
\end{table}

\subsection{Modelo de indicadores}

A plataforma calcula três indicadores em duas escalas espaciais. Na escala do endereço, o Índice de Proximidade Local mede a distância do ponto consultado aos equipamentos públicos mais próximos. Na escala do distrito, o Índice de Vitalidade do Espaço Público mede a oferta de infraestrutura da região, e os indicadores oficiais de ocorrências descrevem os registros criminais. Por decisão de projeto, os indicadores de infraestrutura e os de ocorrências são exibidos de forma paralela e nunca fundidos em uma nota única de ``segurança'' ou ``periculosidade''. Essa separação atende à crítica de \citeonline{caldeira2000} à ``fala do crime'' e evita que a ferramenta reproduza os mapas de perigo que alimentam a segregação.

Os três indicadores compartilham a mesma classificação de equipamentos, agrupados em seis categorias: (i) lazer e convivência (praças, parques e áreas verdes); (ii) cultura e esporte (bibliotecas, centros culturais, CEUs, clubes-escola e equipamentos esportivos); (iii) mobilidade (pontos e terminais de ônibus, estações de metrô e trem, ciclovias); (iv) comércio de rua de uso público (feiras livres e mercados municipais); (v) serviços essenciais (unidades de saúde e escolas públicas); e (vi) conectividade e permanência (pontos de Wi-Fi público e, quando disponível, iluminação pública). Cada categoria reúne um ou mais tipos de equipamento, como unidade básica de saúde, escola pública, praça, parque, ponto de ônibus ou estação de metrô e trem.

O \textbf{Índice de Proximidade Local} (IPL) é o indicador específico do endereço consultado. Para cada tipo de equipamento $j$, o sistema calcula a distância $d_j$, em metros, entre o ponto obtido pelo CEP e o equipamento mais próximo daquele tipo, por meio de busca de vizinho mais próximo no banco geoespacial (operador de distância indexado do PostGIS sobre o tipo \emph{geography}). A distância é convertida em uma nota de 0 a 100 por uma função linear decrescente, definida por dois limiares próprios de cada tipo: a distância ideal $a_j$, até a qual a nota é máxima, e a distância máxima $b_j$, a partir da qual a nota é nula (Equação~\ref{eq:nota}):

\begin{equation}
	\label{eq:nota}
	s_j =
	\begin{cases}
		100,                                    & \text{se } d_j \leq a_j    \\[4pt]
		100 \cdot \dfrac{b_j - d_j}{b_j - a_j}, & \text{se } a_j < d_j < b_j \\[8pt]
		0,                                      & \text{se } d_j \geq b_j
	\end{cases}
\end{equation}

A nota da categoria $k$, $S_k$, é a média das notas dos tipos que a compõem, e o IPL do endereço $c$ é a média ponderada das categorias, com os mesmos pesos $w_k$ ajustáveis pelo usuário (Equação~\ref{eq:ipl}):

\begin{equation}
	\label{eq:ipl}
	\text{IPL}_c = \frac{\sum_{k=1}^{6} w_k \, S_k}{\sum_{k=1}^{6} w_k}
\end{equation}

Os limiares iniciais, apresentados no Quadro~\ref{qua:limiares}, tomam como referência a distância de 400 metros, percorrida em cerca de cinco minutos de caminhada, e são tratados como parâmetros de configuração, a serem revistos com base nos testes com os moradores. Além da nota, o sistema informa para cada tipo a distância em linha reta, o tempo estimado de caminhada e a quantidade de equipamentos a até 500 metros e a até 1 quilômetro.

\begin{table}[htb]
	\caption{Limiares iniciais de distância por tipo de equipamento}
	\label{qua:limiares}
	\centering
	\small
	\begin{tabularx}{\textwidth}{L{3.6cm}XC{2.2cm}C{2.2cm}}
		\toprule
		\textbf{Categoria}  & \textbf{Tipos de equipamento}                               & \textbf{Ideal $a_j$ (m)} & \textbf{Máxima $b_j$ (m)} \\
		\midrule
		Lazer e convivência & Praça; parque ou área verde                                 & 400                      & 1\,200                    \\
		Cultura e esporte   & Biblioteca ou centro cultural; CEU ou equipamento esportivo & 800                      & 2\,000                    \\
		Mobilidade          & Ponto de ônibus                                             & 300                      & 800                       \\
		Mobilidade          & Estação de metrô ou trem; ciclovia                          & 800                      & 2\,000                    \\
		Comércio de rua     & Feira livre ou mercado municipal                            & 800                      & 1\,500                    \\
		Serviços essenciais & Unidade básica de saúde; escola pública                     & 800                      & 2\,000                    \\
		Conectividade       & Ponto de Wi-Fi público                                      & 400                      & 1\,200                    \\
		\bottomrule
	\end{tabularx}
	\fonte{Elaborado pelo autor (2026).}
\end{table}

O \textbf{Índice de Vitalidade do Espaço Público} (IVEP) complementa o IPL na escala do distrito, indicando se a região como um todo oferece infraestrutura que atrai uso cotidiano do espaço. Para cada categoria $k$ e distrito $d$, calcula-se a densidade de equipamentos por quilômetro quadrado e por dez mil habitantes, convertida em posição percentil entre os 96 distritos ($P_{k,d}$, de 0 a 100). O índice é a média ponderada dos percentis, conforme a Equação~\ref{eq:ivep}:

\begin{equation}
	\label{eq:ivep}
	\text{IVEP}_d = \frac{\sum_{k=1}^{6} w_k \, P_{k,d}}{\sum_{k=1}^{6} w_k}
\end{equation}

\noindent em que $w_k$ é o peso da categoria $k$, igual a 1 por padrão e ajustável pelo usuário entre 0 e 5, de modo que cada pessoa possa refletir suas prioridades de rotina.

O terceiro indicador compõe os \textbf{indicadores oficiais de ocorrências}, estimados com base nos microdados georreferenciados provenientes de boletins de ocorrência divulgados pela SSP-SP \cite{ssp2025}, correlacionados aos distritos mediante operação espacial do tipo ponto-em-polígono. Em conformidade com as orientações do manual de interpretação elaborado pela própria secretaria \cite{ssp2005}, que ressalta a subnotificação e advertência quanto a comparações entre unidades de dimensões distintas, os valores são expressos como taxa anualizada por cem mil habitantes, discriminados por natureza da ocorrência (roubo, furto, roubo de celular, roubo de veículo e crimes contra a vida), calculados em janela móvel trienal, conforme a Equação~\ref{eq:taxa}:

\begin{equation}
\label{eq:taxa}
T_{n,d} = \frac{O_{n,d}}{\text{Pop}_d \times a} \times 100\,000
\end{equation}

\noindent em que $O_{n,d}$ denota o quantitativo de ocorrências da natureza $n$ localizadas no distrito $d$ ao longo do período considerado, $\text{Pop}_d$ representa a população do distrito conforme dados do Censo Demográfico 2022 \cite{ibge2022} e $a$ corresponde ao número de anos abrangidos pela janela. Cada taxa é disponibilizada acompanhada da média municipal, da tendência observada no período e de notificação permanente sobre a questão da subnotificação.

\subsection{Fontes de dados e integrações}

O GeoHabita consome exclusivamente fontes abertas ou gratuitas, organizadas em dois fluxos. As consultas de endereço ocorrem em tempo real e passam por camada de cache; as bases estatísticas e cartográficas são importadas por rotinas periódicas de extração, transformação e carga (ETL) e versionadas com data de referência. O Quadro~\ref{qua:fontes} sintetiza as integrações previstas.

\begin{longtable}{L{3.1cm}L{4.3cm}L{2.3cm}L{4.4cm}}
	\caption{Fontes de dados e integrações}
	\label{qua:fontes}                                                                                                                                                                                                                                \\
	\toprule
	\textbf{Fonte}                        & \textbf{Uso no sistema}                                                                              & \textbf{Acesso}                    & \textbf{Restrições}                                           \\
	\midrule
	\endfirsthead
	\toprule
	\textbf{Fonte}                        & \textbf{Uso no sistema}                                                                              & \textbf{Acesso}                    & \textbf{Restrições}                                           \\
	\midrule
	\endhead
	\bottomrule
	\endfoot
	ViaCEP                                & Validação do CEP e obtenção de logradouro, bairro e código IBGE                                      & API REST, tempo real               & Não retorna coordenadas; bloqueio em uso massivo              \\
	BrasilAPI (CEP v2) e AwesomeAPI       & Contingência do ViaCEP e obtenção de latitude e longitude                                            & API REST, tempo real               & Coordenadas aproximadas e eventualmente ausentes              \\
	Nominatim (OpenStreetMap)             & Geocodificação de endereço completo quando o CEP não traz coordenadas                                & API REST, tempo real               & Até 1 requisição por segundo; cache e atribuição obrigatórios \\
	GeoSampa                              & Polígonos dos 96 distritos e camadas de equipamentos públicos, áreas verdes, transporte e mobilidade & WFS, carga periódica               & Sistema de referência SIRGAS 2000/UTM 23S; exige reprojeção   \\
	Portal de dados abertos da Prefeitura & Feiras livres, pontos de Wi-Fi público e demais cadastros municipais complementares                  & Arquivos CSV/JSON, carga periódica & Coordenadas a validar; atualização irregular                  \\
	SSP-SP                                & Microdados de boletins de ocorrência com coordenadas                                                 & Planilhas anuais, carga periódica  & Colunas variam entre anos; coordenadas nulas; subnotificação  \\
	IBGE (Censo 2022)                     & População por distrito para cálculo de taxas e densidades                                            & Arquivos e API, carga periódica    & Atualização decenal                                           \\
	OpenStreetMap e Leaflet               & Base cartográfica do mapa interativo                                                                 & Biblioteca JavaScript              & Atribuição obrigatória (licença ODbL)                         \\
\end{longtable}
\vspace{-0.4cm}
\fonte{Elaborado pelo autor (2026).}

\subsection{Requisitos funcionais}

Os requisitos funcionais descrevem os comportamentos observáveis do sistema \cite{sommerville2018}. Foram identificados 25 requisitos, agrupados em cinco módulos, dos quais 14 são classificados como essenciais (\emph{Must}) e compõem o MVP (Quadro~\ref{qua:rf}).

\begin{longtable}{C{1.2cm}L{9.2cm}C{1.5cm}C{2.2cm}}
	\caption{Requisitos funcionais}
	\label{qua:rf}                                                                                                                                                                                                                                \\
	\toprule
	\textbf{ID} & \textbf{Descrição}                                                                                                                                                                          & \textbf{Prior.} & \textbf{Módulo} \\
	\midrule
	\endfirsthead
	\toprule
	\textbf{ID} & \textbf{Descrição}                                                                                                                                                                          & \textbf{Prior.} & \textbf{Módulo} \\
	\midrule
	\endhead
	\bottomrule
	\endfoot
	RF01        & O sistema deve permitir a consulta por CEP, validando o formato de oito dígitos antes de qualquer chamada externa.                                                                          & Must            & Consulta        \\
	RF02        & O sistema deve preencher automaticamente logradouro, bairro e município a partir do CEP informado.                                                                                          & Must            & Consulta        \\
	RF03        & O sistema deve permitir a consulta por endereço textual quando o usuário não souber o CEP.                                                                                                  & Should          & Consulta        \\
	RF04        & O sistema deve converter o endereço em coordenadas geográficas, utilizando serviço de contingência em caso de falha do serviço principal.                                                   & Must            & Consulta        \\
	RF05        & O sistema deve identificar o distrito administrativo correspondente às coordenadas por operação de ponto-em-polígono.                                                                       & Must            & Consulta        \\
	RF06        & O sistema deve informar ao usuário quando o endereço estiver fora do município de São Paulo.                                                                                                & Must            & Consulta        \\
	RF07        & O sistema deve listar os equipamentos públicos de uso do espaço em raio de 500 metros e de 1 quilômetro do ponto consultado, agrupados pelas seis categorias do IVEP, com nome e distância. & Must            & Diagnóstico     \\
	RF23        & O sistema deve calcular, para o endereço consultado, a distância em linha reta até o equipamento mais próximo de cada tipo, com o tempo estimado de caminhada.                              & Must            & Diagnóstico     \\
	RF24        & O sistema deve calcular e exibir o Índice de Proximidade Local do endereço e a nota de cada categoria, destacando os tipos de equipamento ausentes no raio máximo.                          & Must            & Diagnóstico     \\
	RF25        & O sistema deve calcular a distância e o tempo pela rede de calçadas e vias até o equipamento mais próximo de cada tipo, quando o usuário solicitar.                                         & Could           & Diagnóstico     \\
	RF08        & O sistema deve exibir o IVEP do distrito, o percentil de cada categoria e a comparação com a mediana do município.                                                                          & Must            & Diagnóstico     \\
	RF09        & O sistema deve permitir que o usuário ajuste os pesos das categorias do IPL e do IVEP e recalcule os índices sem recarregar a página.                                                       & Should          & Diagnóstico     \\
	RF10        & O sistema deve exibir as taxas de ocorrências por cem mil habitantes do distrito, por natureza, com média municipal e tendência dos últimos três anos.                                      & Must            & Diagnóstico     \\
	RF11        & O sistema deve apresentar, junto às taxas criminais, nota explicativa sobre subnotificação e sobre a metodologia de cálculo.                                                                & Must            & Diagnóstico     \\
	RF12        & O sistema deve exibir a fonte e a data de referência de cada indicador apresentado.                                                                                                         & Must            & Diagnóstico     \\
	RF13        & O sistema deve exibir mapa interativo com o ponto consultado, o limite do distrito e os equipamentos do entorno, filtráveis por categoria.                                                  & Must            & Visualização    \\
	RF14        & O sistema deve oferecer, para todo conteúdo cartográfico, tabela textual equivalente navegável por teclado e leitor de tela.                                                                & Must            & Visualização    \\
	RF15        & O sistema deve anunciar por região viva (\emph{live region}) WAI-ARIA a conclusão de cada consulta assíncrona e eventuais erros.                                                            & Should          & Visualização    \\
	RF16        & O sistema deve permitir a comparação lado a lado de até três endereços ou distritos.                                                                                                        & Should          & Comparação      \\
	RF17        & O sistema deve permitir compartilhar um diagnóstico por link, sem armazenar dados pessoais do consulente.                                                                                   & Could           & Comparação      \\
	RF18        & O sistema deve permitir exportar o diagnóstico em formato PDF acessível ou CSV.                                                                                                             & Could           & Comparação      \\
	RF19        & O sistema deve permitir ao administrador executar as rotinas de carga das bases externas e registrar o resultado de cada execução.                                                          & Should          & Administração   \\
	RF20        & O sistema deve validar as bases importadas, descartando registros sem coordenadas ou fora dos limites do município e reportando a quantidade descartada.                                    & Should          & Administração   \\
	RF21        & O sistema deve manter o histórico de versões dos indicadores com data de referência de cada carga.                                                                                          & Could           & Administração   \\
	RF22        & O sistema deve disponibilizar uma API pública de leitura dos indicadores agregados por distrito, documentada em OpenAPI.                                                                    & Could           & Administração   \\
\end{longtable}
\vspace{-0.4cm}
\fonte{Elaborado pelo autor (2026).}

\subsection{Requisitos não funcionais}

Os requisitos não funcionais estabelecem restrições sobre os serviços oferecidos e sobre o processo de desenvolvimento \cite{sommerville2018}. Para garantir sua verificabilidade, cada requisito de produto está associado a uma métrica objetiva, conforme recomendam \citeonline{pressman2021}. Os requisitos de produto estão classificados segundo as características de qualidade da ISO/IEC 25010 (Quadro~\ref{qua:rnf}).

\begin{longtable}{C{1.3cm}L{2.6cm}L{6.1cm}L{4.1cm}}
	\caption{Requisitos não funcionais de produto}
	\label{qua:rnf}                                                                                                                                                                                                                                                                                                                                             \\
	\toprule
	\textbf{ID} & \textbf{Característica}              & \textbf{Requisito}                                                                                                                                                                                    & \textbf{Métrica de verificação}                                                                \\
	\midrule
	\endfirsthead
	\toprule
	\textbf{ID} & \textbf{Característica}              & \textbf{Requisito}                                                                                                                                                                                    & \textbf{Métrica de verificação}                                                                \\
	\midrule
	\endhead
	\bottomrule
	\endfoot
	RNF01       & Eficiência de desempenho             & O diagnóstico de um CEP já armazenado em cache deve ser exibido em até 2 segundos.                                                                                                                    & Percentil 95 $\leq$ 2 s com 30 usuários simultâneos                                            \\
	RNF02       & Eficiência de desempenho             & O diagnóstico de um CEP inédito deve ser exibido em até 6 segundos.                                                                                                                                   & Percentil 95 $\leq$ 6 s em teste de carga                                                      \\
	RNF03       & Usabilidade (acessibilidade)         & A interface deve atender às diretrizes WCAG 2.1 no nível AA e empregar marcações WAI-ARIA em todos os componentes dinâmicos.                                                                          & Zero erros críticos em Lighthouse e axe; contraste $\geq$ 4,5:1                                \\
	RNF04       & Usabilidade (operabilidade)          & Todas as funções devem ser operáveis exclusivamente por teclado, com foco visível e ordem lógica.                                                                                                     & Roteiro de teste manual com NVDA e Orca                                                        \\
	RNF05       & Usabilidade (inteligibilidade)       & Os textos explicativos devem usar linguagem simples, sem jargão técnico ou geográfico sem explicação.                                                                                                 & Avaliação em teste com usuários; SUS $\geq$ 70                                                 \\
	RNF06       & Confiabilidade (tolerância a falhas) & A falha de uma fonte externa não deve impedir a exibição dos demais blocos do diagnóstico.                                                                                                            & Teste com serviço externo simulado fora do ar                                                  \\
	RNF07       & Confiabilidade (disponibilidade)     & A aplicação deve manter disponibilidade mensal mínima de 98\%.                                                                                                                                        & Monitoramento externo de disponibilidade                                                       \\
	RNF08       & Segurança                            & Todo o tráfego deve ocorrer por HTTPS, e as rotas públicas devem limitar-se a 60 requisições por minuto por endereço IP.                                                                              & Teste automatizado; resposta HTTP 429 acima do limite                                          \\
	RNF09       & Segurança                            & A aplicação não deve apresentar vulnerabilidades de alta severidade do OWASP Top 10.                                                                                                                  & Varredura com OWASP ZAP e auditoria de dependências                                            \\
	RNF10       & Manutenibilidade                     & O código deve adotar arquitetura em camadas (apresentação, aplicação, domínio e infraestrutura), com adaptadores isolando cada serviço externo.                                                       & Revisão em \emph{pull request}; ausência de chamadas externas fora da camada de infraestrutura \\
	RNF11       & Manutenibilidade (testabilidade)     & As regras de negócio devem ter cobertura mínima de 70\% por testes unitários.                                                                                                                         & Relatório de cobertura na integração contínua                                                  \\
	RNF12       & Compatibilidade e portabilidade      & A interface deve funcionar nas duas versões mais recentes dos navegadores Chrome, Firefox, Edge e Safari, em telas a partir de 360 pixels de largura, e a aplicação deve ser empacotada em contêiner. & Roteiro multiplataforma; implantação em nuvem com um único comando                             \\
\end{longtable}
\vspace{-0.4cm}
\fonte{Elaborado pelo autor (2026).}

Os requisitos organizacionais e externos, que não se referem diretamente ao comportamento do produto, estão reunidos no Quadro~\ref{qua:rnf-org}.

\begin{table}[htb]
	\caption{Requisitos não funcionais organizacionais e externos}
	\label{qua:rnf-org}
	\centering
	\small
	\begin{tabularx}{\textwidth}{C{1.3cm}C{2.4cm}X}
		\toprule
		\textbf{ID} & \textbf{Tipo}      & \textbf{Requisito}                                                                                                                                                                                                         \\
		\midrule
		RNF13       & Organizacional     & O desenvolvimento deve seguir práticas do Scrum, com \emph{sprints} de duas semanas, controle de versão em Git e integração contínua executando os testes a cada envio ao repositório.                                     \\
		RNF14       & Organizacional     & A pilha tecnológica deve compreender Python com micro-framework web no serviço de retaguarda, JavaScript no cliente, PostgreSQL com extensão PostGIS na persistência e hospedagem em provedor de nuvem com plano gratuito. \\
		RNF15       & Externo (legal)    & O tratamento de dados deve observar a Lei Geral de Proteção de Dados \cite{lgpd2018}: a consulta não exige cadastro e nenhum CEP consultado é associado a pessoa identificável.                                            \\
		RNF16       & Externo (legal)    & A plataforma deve atender ao art. 63 da Lei Brasileira de Inclusão \cite{lbi2015}, que torna obrigatória a acessibilidade em sítios da internet.                                                                           \\
		RNF17       & Externo (licenças) & Todas as fontes devem ser atribuídas conforme suas licenças, e o código deve ser distribuído sob a GNU AGPL v3 \cite{agplv3}.                                                                                              \\
		RNF18       & Externo (ético)    & Os indicadores não devem empregar rótulos estigmatizantes, rankings de ``bairros perigosos'' nem escalas cromáticas de alarme para dados criminais.                                                                        \\
		\bottomrule
	\end{tabularx}
	\fonte{Elaborado pelo autor (2026).}
\end{table}

\subsection{Regras de negócio e requisitos de domínio}

As regras de negócio condensam o conhecimento do domínio urbano e estatístico que o sistema deve respeitar independentemente da interface. Elas são implementadas na camada de domínio e testadas de forma isolada (Quadro~\ref{qua:rn}).

\begin{longtable}{C{1.2cm}L{10.6cm}L{2.6cm}}
	\caption{Regras de negócio}
	\label{qua:rn}                                                                                                                                                                                                                                                                                               \\
	\toprule
	\textbf{ID} & \textbf{Regra}                                                                                                                                                                                                                                                        & \textbf{Relacionada a} \\
	\midrule
	\endfirsthead
	\toprule
	\textbf{ID} & \textbf{Regra}                                                                                                                                                                                                                                                        & \textbf{Relacionada a} \\
	\midrule
	\endhead
	\bottomrule
	\endfoot
	RN01        & Um CEP válido possui exatamente oito dígitos numéricos; hífens, pontos e espaços são removidos antes da validação.                                                                                                                                                    & RF01                   \\
	RN02        & Um CEP inexistente gera a mensagem ``CEP não encontrado'' e libera a consulta por endereço textual.                                                                                                                                                                   & RF02, RF03             \\
	RN03        & O serviço de contingência é acionado quando o serviço principal retorna erro ou excede 3 segundos de resposta.                                                                                                                                                        & RF04, RNF06            \\
	RN04        & Resultados de CEP e geocodificação permanecem em cache por 30 dias; indicadores agregados, até a próxima carga.                                                                                                                                                       & RF04, RNF01            \\
	RN05        & A coordenada obtida pelo CEP é tratada como aproximada; a interface informa que ela identifica o distrito, mas não a posição exata do imóvel.                                                                                                                         & RF05, RF07             \\
	RN06        & O IVEP é calculado por percentis entre os 96 distritos e varia de 0 a 100; os pesos padrão são iguais a 1 e aceitam valores inteiros de 0 a 5.                                                                                                                        & RF08, RF09             \\
	RN07        & Os indicadores de vitalidade e os indicadores de ocorrências nunca são combinados em uma pontuação única.                                                                                                                                                             & RF08, RF10, RNF18      \\
	RN08        & As taxas de ocorrências usam janela de três anos completos e são expressas por cem mil habitantes; o ano corrente incompleto não entra no cálculo.                                                                                                                    & RF10                   \\
	RN09        & Naturezas com menos de 30 ocorrências no distrito durante a janela são exibidas como ``número insuficiente para comparação'', sem taxa nem tendência.                                                                                                                 & RF10, RF11             \\
	RN10        & Registros da SSP-SP com coordenadas nulas, iguais a zero ou fora do município são descartados na carga e contabilizados no relatório de validação.                                                                                                                    & RF20                   \\
	RN11        & Indicadores com data de referência superior a 24 meses recebem o aviso ``dado desatualizado''.                                                                                                                                                                        & RF12                   \\
	RN12        & A distância até o equipamento mais próximo é a distância geodésica em linha reta, em metros, entre o ponto do CEP e o ponto do equipamento; para áreas como parques, usa-se o ponto mais próximo do contorno. A interface identifica a medida como ``em linha reta''. & RF23                   \\
	RN13        & A nota de cada tipo segue a Equação~\ref{eq:nota}, com os limiares do Quadro~\ref{qua:limiares}; os limiares são parâmetros de configuração alteráveis pelo administrador sem mudança de código.                                                                      & RF24                   \\
	RN14        & O tempo estimado de caminhada é calculado à velocidade de 80 metros por minuto (4,8 km/h) e arredondado para o minuto inteiro superior.                                                                                                                               & RF23                   \\
	RN15        & Tipo sem nenhum equipamento a menos da distância máxima recebe nota zero, e a interface informa ``nenhum equipamento deste tipo a menos de $b_j$ metros''.                                                                                                            & RF24                   \\
	RN16        & O IPL é calculado sempre sobre a coordenada do endereço consultado, e nunca sobre o centroide do distrito ou do bairro.                                                                                                                                               & RF24                   \\
\end{longtable}
\vspace{-0.4cm}
\fonte{Elaborado pelo autor (2026).}

A regra RN09 transpõe para o sistema a advertência do manual de interpretação da SSP-SP sobre comparações apoiadas em poucos casos, e a regra RN07 materializa a opção, discutida na fundamentação, de não reduzir a qualidade de um território a um único número.

\subsection{Casos de uso}

Os requisitos funcionais foram agrupados em oito casos de uso principais, que orientam o diagrama de casos de uso e o desenho das telas:

\begin{alineas}
	\item \textbf{UC01 -- Consultar endereço}: o munícipe informa CEP ou endereço, e o sistema valida, geocodifica e identifica o distrito (RF01 a RF06);
	\item \textbf{UC02 -- Visualizar diagnóstico socioespacial}: o sistema apresenta equipamentos do entorno, IVEP e indicadores de ocorrências, com fontes e datas (RF07, RF08, RF10 a RF12);
	\item \textbf{UC02a -- Avaliar proximidade do endereço}: o sistema calcula a distância até o equipamento mais próximo de cada tipo e o IPL do endereço (RF23 a RF25);
	\item \textbf{UC03 -- Personalizar prioridades}: o munícipe ajusta os pesos das categorias e acompanha o recálculo do IVEP (RF09);
	\item \textbf{UC04 -- Explorar mapa ou tabela equivalente}: o munícipe navega pelo mapa ou pela tabela textual e filtra categorias (RF13 a RF15);
	\item \textbf{UC05 -- Comparar localidades}: o munícipe compara até três endereços ou distritos (RF16);
	\item \textbf{UC06 -- Compartilhar ou exportar diagnóstico}: o munícipe gera link ou arquivo do diagnóstico (RF17, RF18);
	\item \textbf{UC07 -- Atualizar bases de dados}: o administrador executa as cargas, valida os registros e publica nova versão dos indicadores (RF19 a RF21).
\end{alineas}

\subsection{Histórias de usuário e critérios de aceitação}

Para o planejamento das \emph{sprints}, os casos de uso foram desdobrados em histórias de usuário no formato ``Como \emph{papel}, quero \emph{ação}, para \emph{benefício}'' \cite{cohn2004}, estimadas em pontos de complexidade relativa (Quadro~\ref{qua:hu}).

\begin{longtable}{C{1.2cm}L{9.0cm}C{2.1cm}C{1.3cm}}
	\caption{Histórias de usuário}
	\label{qua:hu}                                                                                                                                                                                                                             \\
	\toprule
	\textbf{ID} & \textbf{História}                                                                                                                                                                    & \textbf{Requisitos} & \textbf{Pontos} \\
	\midrule
	\endfirsthead
	\toprule
	\textbf{ID} & \textbf{História}                                                                                                                                                                    & \textbf{Requisitos} & \textbf{Pontos} \\
	\midrule
	\endhead
	\bottomrule
	\endfoot
	HU01        & Como munícipe, quero digitar um CEP e ver o endereço reconhecido, para iniciar a consulta sem preencher formulários longos.                                                          & RF01, RF02          & 3               \\
	HU02        & Como munícipe, quero saber em que distrito fica o endereço, para entender a região a que ele pertence.                                                                               & RF04--RF06          & 5               \\
	HU03        & Como munícipe, quero ver praças, parques, feiras, transporte e serviços públicos a até 1 km, para saber se a vizinhança tem vida nas ruas.                                           & RF07, RF13          & 8               \\
	HU04        & Como munícipe, quero ver o índice de vitalidade do distrito em comparação com a cidade, para avaliar a região de forma objetiva.                                                     & RF08                & 5               \\
	HU05        & Como munícipe, quero ver as taxas oficiais de ocorrências com explicação sobre seus limites, para decidir com informação confiável e sem alarmismo.                                  & RF10--RF12          & 8               \\
	HU06        & Como usuário de leitor de tela, quero acessar todo o conteúdo do mapa em forma de tabela e ser avisado quando a consulta terminar, para usar a plataforma com autonomia.             & RF14, RF15          & 5               \\
	HU07        & Como munícipe, quero dar mais peso ao transporte do que ao lazer, para que o índice reflita minha rotina.                                                                            & RF09                & 5               \\
	HU08        & Como munícipe, quero comparar três endereços lado a lado, para escolher entre as opções finais.                                                                                      & RF16                & 8               \\
	HU09        & Como administrador, quero carregar e validar novas bases da SSP-SP e do GeoSampa, para manter os indicadores atualizados.                                                            & RF19, RF20          & 8               \\
	HU10        & Como munícipe, quero saber a que distância do endereço ficam a UBS, a escola, a praça, o ponto de ônibus e a estação mais próximos, para avaliar se consigo fazer minha rotina a pé. & RF23, RF24          & 8               \\
\end{longtable}
\vspace{-0.4cm}
\fonte{Elaborado pelo autor (2026).}

Os critérios de aceitação são redigidos na sintaxe Gherkin (Dado, Quando, Então), o que permite convertê-los diretamente em testes automatizados de aceitação. A título de exemplo, a história HU01 é aceita quando todos os cenários a seguir forem atendidos:

\begin{alineas}
	\item \textbf{CEP válido}: \emph{dado} que o munícipe está na página inicial, \emph{quando} informa ``01310-100'', \emph{então} o sistema exibe ``Avenida Paulista, Bela Vista, São Paulo -- SP'';
	\item \textbf{Formato inválido}: \emph{quando} informa ``1234'', \emph{então} o sistema exibe ``O CEP deve ter 8 dígitos'' e nenhuma API externa é chamada;
	\item \textbf{CEP inexistente}: \emph{quando} informa ``99999-999'', \emph{então} o sistema exibe ``CEP não encontrado'' e oferece a consulta por endereço;
	\item \textbf{Serviço principal indisponível}: \emph{dado} que o ViaCEP não responde em 3 segundos, \emph{quando} o munícipe informa um CEP válido, \emph{então} o endereço é obtido pelo serviço de contingência.
\end{alineas}

Considera-se uma história pronta para desenvolvimento (\emph{Definition of Ready}) quando possui critérios de aceitação escritos, estimativa acordada e dependências externas identificadas. Considera-se concluída (\emph{Definition of Done}) quando o código foi revisado, os testes unitários e de aceitação passam na integração contínua, a verificação de acessibilidade não aponta erros críticos e o incremento está implantado no ambiente em nuvem.

\subsection{Priorização e planejamento incremental}

As dez histórias totalizam 63 pontos, distribuídos em quatro \emph{sprints} de duas semanas, precedidas de uma \emph{sprint} de preparação do ambiente (Quadro~\ref{qua:sprints}). O MVP encerra-se ao fim da terceira \emph{sprint}, quando o fluxo completo de consulta e diagnóstico estará disponível para validação com os moradores.

\begin{table}[htb]
	\caption{Planejamento das \emph{sprints}}
	\label{qua:sprints}
	\centering
	\small
	\begin{tabularx}{\textwidth}{C{1.4cm}XC{3.2cm}C{1.4cm}}
		\toprule
		\textbf{Sprint} & \textbf{Meta}                                                                                  & \textbf{Histórias} & \textbf{Pontos} \\
		\midrule
		0               & Repositório, integração contínua, banco com PostGIS e primeira carga dos distritos do GeoSampa & --                 & --              \\
		1               & Consultar um CEP e localizar o distrito                                                        & HU01, HU02, HU09   & 16              \\
		2               & Exibir a infraestrutura do entorno e a distância ao equipamento mais próximo de cada tipo      & HU03, HU10         & 16              \\
		3               & Calcular IVEP e indicadores de ocorrências, com mapa e tabela acessível (MVP)                  & HU04, HU05, HU06   & 18              \\
		4               & Personalização de pesos, comparação e ajustes após testes com usuários                         & HU07, HU08         & 13              \\
		\bottomrule
	\end{tabularx}
	\fonte{Elaborado pelo autor (2026).}
\end{table}

\subsection{Plano de testes}

A estratégia de verificação segue a pirâmide de testes: grande volume de testes unitários sobre as regras de negócio, testes de integração sobre os adaptadores de dados e o banco geoespacial, e um conjunto reduzido de testes ponta a ponta sobre os fluxos críticos \cite{pressman2021}. Nos testes automatizados, os serviços externos são substituídos por simulações (\emph{mocks}), o que torna os resultados reprodutíveis; testes de contrato agendados verificam, separadamente, se o formato das respostas reais das APIs permanece inalterado. O Quadro~\ref{qua:niveis} apresenta os níveis de teste previstos.

\begin{longtable}{L{2.6cm}L{5.0cm}L{3.5cm}L{3.0cm}}
	\caption{Níveis e tipos de teste}
	\label{qua:niveis}                                                                                                                                                                       \\
	\toprule
	\textbf{Nível} & \textbf{Objeto de verificação}                                  & \textbf{Ferramentas}                           & \textbf{Critério de saída}                           \\
	\midrule
	\endfirsthead
	\toprule
	\textbf{Nível} & \textbf{Objeto de verificação}                                  & \textbf{Ferramentas}                           & \textbf{Critério de saída}                           \\
	\midrule
	\endhead
	\bottomrule
	\endfoot
	Unitário       & Regras RN01 a RN16, cálculo do IPL, do IVEP e das taxas         & pytest                                         & 100\% aprovados; cobertura $\geq$ 70\%               \\
	Integração     & Adaptadores de CEP e geocodificação, consultas espaciais, cache & pytest com banco PostGIS em contêiner          & 100\% aprovados                                      \\
	Contrato       & Formato das respostas de ViaCEP, BrasilAPI e Nominatim          & Execução diária agendada                       & Alerta em caso de divergência                        \\
	Aceitação      & Cenários Gherkin das histórias de usuário                       & pytest-bdd ou Behave                           & Todas as histórias \emph{Must} aprovadas             \\
	Ponta a ponta  & Fluxo consulta, diagnóstico, mapa e tabela                      & Playwright                                     & Aprovado em Chrome e Firefox                         \\
	Desempenho     & RNF01 e RNF02                                                   & k6                                             & Metas de percentil 95 atendidas                      \\
	Segurança      & RNF08 e RNF09                                                   & OWASP ZAP; pip-audit                           & Nenhum alerta alto                                   \\
	Acessibilidade & RNF03 e RNF04                                                   & Lighthouse, axe, NVDA, Orca                    & Zero erros críticos                                  \\
	Usabilidade    & RNF05 e aderência às necessidades levantadas                    & Teste moderado com moradores; questionário SUS & Tarefas concluídas por $\geq$ 80\% dos participantes \\
\end{longtable}
\vspace{-0.4cm}
\fonte{Elaborado pelo autor (2026).}

Os casos de teste derivam diretamente dos requisitos e das regras de negócio. O Quadro~\ref{qua:ct} relaciona os principais.

\begin{longtable}{C{1.1cm}C{1.8cm}L{4.6cm}L{4.4cm}C{1.8cm}}
	\caption{Casos de teste principais}
	\label{qua:ct}                                                                                                                                                                                \\
	\toprule
	\textbf{ID} & \textbf{Requisito} & \textbf{Entrada ou condição}                                 & \textbf{Resultado esperado}                                                & \textbf{Tipo}  \\
	\midrule
	\endfirsthead
	\toprule
	\textbf{ID} & \textbf{Requisito} & \textbf{Entrada ou condição}                                 & \textbf{Resultado esperado}                                                & \textbf{Tipo}  \\
	\midrule
	\endhead
	\bottomrule
	\endfoot
	CT01        & RF01, RN01         & ``1234'', ``abcdefgh'' e ``123456789''                       & Mensagem de formato inválido e nenhuma chamada externa                     & Unitário       \\
	CT02        & RF02               & CEP 01310-100                                                & Logradouro, bairro e município preenchidos corretamente                    & Integração     \\
	CT03        & RN02               & CEP 99999-999                                                & ``CEP não encontrado'' e consulta textual liberada                         & Integração     \\
	CT04        & RN03, RNF06        & Serviço principal simulado com atraso de 5 s                 & Resposta obtida pelo serviço de contingência                               & Integração     \\
	CT05        & RF05               & Coordenadas conhecidas no interior de três distritos         & Distrito correto em todos os casos                                         & Integração     \\
	CT06        & RF06               & CEP de município vizinho                                     & Aviso de endereço fora de São Paulo                                        & Integração     \\
	CT07        & RF07               & Ponto com equipamentos cadastrados a 300 m, 800 m e 1,5 km   & Os dois primeiros listados com distância; o terceiro ausente               & Integração     \\
	CT08        & RF08, RN06         & Base fictícia com percentis conhecidos e pesos iguais        & IVEP igual à média aritmética dos percentis                                & Unitário       \\
	CT09        & RF09               & Peso de mobilidade igual a 5 e demais iguais a 1             & IVEP recalculado conforme a Equação~\ref{eq:ivep}, sem recarregar a página & Ponta a ponta  \\
	CT10        & RN08               & Distrito com 300 ocorrências em 3 anos e 100\,000 habitantes & Taxa de 100 por cem mil habitantes ao ano                                  & Unitário       \\
	CT11        & RN09               & Natureza com 12 ocorrências na janela                        & ``Número insuficiente para comparação'', sem taxa                          & Unitário       \\
	CT12        & RN07               & Resposta da API interna do diagnóstico                       & Ausência de qualquer campo que combine os dois conjuntos de indicadores    & Unitário       \\
	CT13        & RN10, RF20         & Planilha com 5 registros de coordenadas nulas ou zeradas     & Registros descartados e contabilizados no relatório de carga               & Integração     \\
	CT14        & RF14, RNF04        & Navegação apenas por teclado com leitor de tela              & Todos os equipamentos do mapa acessíveis pela tabela equivalente           & Acessibilidade \\
	CT15        & RF15               & Consulta concluída e consulta com erro                       & Ambas anunciadas pela região viva                                          & Acessibilidade \\
	CT16        & RNF08              & 61 requisições em um minuto a partir do mesmo IP             & A 61ª requisição recebe HTTP 429                                           & Segurança      \\
	CT17        & RF23, RN12         & Ponto com duas UBS cadastradas a 250 m e a 600 m             & UBS mais próxima informada a 250 m, ``em linha reta''                      & Integração     \\
	CT18        & RN13               & Ponto de ônibus a 550 m ($a_j = 300$, $b_j = 800$)           & Nota do tipo igual a 50                                                    & Unitário       \\
	CT19        & RN14               & Distância de 650 m                                           & Tempo estimado de 9 minutos                                                & Unitário       \\
	CT20        & RN15               & Nenhuma estação a menos de 2\,000 m                          & Nota zero e mensagem de ausência                                           & Unitário       \\
	CT21        & RF24               & Notas de categoria 100, 50, 0, 80, 60 e 40, pesos iguais a 1 & IPL igual a 55                                                             & Unitário       \\
	CT22        & RN16               & Dois CEPs distintos no mesmo distrito                        & IPL diferentes e mesmo IVEP                                                & Integração     \\
\end{longtable}
\vspace{-0.4cm}
\fonte{Elaborado pelo autor (2026).}

Os defeitos encontrados são registrados como \emph{issues} no repositório, classificados por severidade; defeitos críticos e altos impedem o encerramento da \emph{sprint}.

\subsection{Matriz de rastreabilidade}

A matriz de rastreabilidade (Quadro~\ref{qua:rastreab}) liga as necessidades identificadas junto à comunidade aos requisitos que as atendem e aos testes que comprovam seu cumprimento, permitindo avaliar o impacto de qualquer alteração de escopo.

\begin{longtable}{L{4.8cm}L{4.2cm}C{2.0cm}C{2.6cm}}
	\caption{Matriz de rastreabilidade}
	\label{qua:rastreab}                                                                                                          \\
	\toprule
	\textbf{Necessidade}                                    & \textbf{Requisitos e regras}  & \textbf{História} & \textbf{Testes} \\
	\midrule
	\endfirsthead
	\toprule
	\textbf{Necessidade}                                    & \textbf{Requisitos e regras}  & \textbf{História} & \textbf{Testes} \\
	\midrule
	\endhead
	\bottomrule
	\endfoot
	Consultar uma região de forma simples                   & RF01--RF06, RN01--RN05        & HU01, HU02        & CT01--CT06      \\
	Saber a que distância ficam os serviços do dia a dia    & RF23--RF25, RN12--RN16        & HU10              & CT17--CT22      \\
	Conhecer a infraestrutura pública do entorno            & RF07, RF13                    & HU03              & CT07            \\
	Avaliar a vitalidade do espaço público                  & RF08, RF09, RN06              & HU04, HU07        & CT08, CT09      \\
	Dispor de dados criminais auditáveis e contextualizados & RF10--RF12, RN07--RN09, RNF18 & HU05              & CT10--CT12      \\
	Usar a plataforma com tecnologia assistiva              & RF14, RF15, RNF03, RNF04      & HU06              & CT14, CT15      \\
	Manter os dados atualizados e íntegros                  & RF19, RF20, RN10              & HU09              & CT13            \\
	Operar com segurança e desempenho                       & RNF01, RNF02, RNF08           & --                & CT16            \\
\end{longtable}
\vspace{-0.4cm}
\fonte{Elaborado pelo autor (2026).}

% posicionado após a subseção "Metodologia" (ou dentro de
% "Resultados preliminares: solução inicial").
%
% Pacotes necessários no preâmbulo do relatório:
%   \usepackage{longtable}
%   \usepackage{tabularx}
%   \usepackage{booktabs}
%   \usepackage{array}
%   \usepackage{amsmath}
% Referências novas: ver requisitos_geohabita.bib
% =====================================================================

\counterwithin{section}{chapter}
\renewcommand{\thesection}{\thechapter.\arabic{section}}
\counterwithin{table}{chapter}
\renewcommand{\thetable}{\thechapter.\arabic{table}}
\counterwithin{equation}{chapter}
\renewcommand{\theequation}{\thechapter.\arabic{equation}}

\newcolumntype{L}[1]{>{\raggedright\arraybackslash}p{#1}}
\newcolumntype{C}[1]{>{\centering\arraybackslash}p{#1}}


\section{Especificação de requisitos}

A especificação de requisitos do GeoHabita consolida, em termos verificáveis, aquilo que a plataforma deve oferecer ao cidadão e as condições de qualidade sob as quais deve operar. Em relação à concepção inicial, o escopo foi reorientado: a segurança deixa de ser tratada prioritariamente a partir da estatística criminal e passa a ser lida também pela infraestrutura pública de uso do espaço, isto é, pelos equipamentos e lugares que atraem pessoas às ruas em horários variados e sustentam o que \citeonline[cap.~1]{jacobs2011} denomina vigilância natural dos ``olhos na rua''. Praças, parques, equipamentos culturais e esportivos, feiras livres, pontos de transporte coletivo, unidades de saúde e escolas passam a compor um indicador próprio de vitalidade do espaço público, apresentado lado a lado com os registros oficiais da SSP-SP \cite{ssp2024,geosampa2024}.

A estrutura adotada segue a norma ISO/IEC/IEEE 29148 \cite{iso29148} e a classificação proposta por \citeonline{sommerville2018}, que distingue requisitos funcionais, requisitos não funcionais (de produto, organizacionais e externos) e requisitos de domínio. Os atributos de qualidade são organizados segundo o modelo da ISO/IEC 25010 \cite{iso25010}, e o desdobramento em histórias de usuário e critérios de aceitação acompanha as práticas ágeis descritas por \citeonline{cohn2004} e \citeonline{pressman2021}. Cada requisito recebe identificador único, prioridade segundo a técnica MoSCoW (\emph{Must}, \emph{Should}, \emph{Could}, \emph{Won't}) e rastreabilidade até os casos de teste que o verificam.

\subsection{Escopo do produto e atores}

O Produto Mínimo Viável (MVP) delimita-se ao município de São Paulo e aos seus 96 distritos administrativos, unidade territorial escolhida por dispor de polígonos oficiais no GeoSampa e de população recenseada pelo IBGE, o que permite o cálculo de taxas comparáveis entre regiões. A consulta parte de um CEP ou de um endereço informado pelo usuário, que é convertido em coordenadas geográficas e associado ao distrito correspondente. A partir desse ponto, o sistema produz um diagnóstico socioespacial composto por três blocos: a proximidade do endereço consultado aos equipamentos públicos de uso do espaço, medida pela distância até o equipamento mais próximo de cada tipo; o indicador de vitalidade do espaço público do distrito; e os indicadores oficiais de ocorrências criminais, sempre contextualizados.

Permanecem fora do escopo do MVP a intermediação ou divulgação de anúncios imobiliários, a coleta colaborativa de relatos de crimes, a predição de ocorrências por modelos de aprendizado de máquina, a extração automatizada não estruturada de páginas (\emph{web scraping}) e aplicativos nativos para dispositivos móveis. A exclusão dos relatos colaborativos decorre, em particular, do risco de reforço de estigmas territoriais e de reidentificação de vítimas já discutido na fundamentação teórica.

Identificam-se três atores principais, descritos no Quadro~\ref{qua:atores}.

\begin{table}[htb]
	\caption{Atores do sistema}
	\label{qua:atores}
	\centering
	\small
	\begin{tabularx}{\textwidth}{L{3.6cm}X}
		\toprule
		\textbf{Ator}                   & \textbf{Descrição e interesse}                                                                                                                                                            \\
		\midrule
		Munícipe consulente             & Pessoa ou família em transição habitacional, ou morador interessado em conhecer o próprio território. Consulta endereços, compara regiões e não precisa de cadastro para o uso essencial. \\
		Usuário de tecnologia assistiva & Munícipe que acessa a plataforma por leitor de tela, navegação exclusiva por teclado ou ampliação. Exige equivalência textual integral para mapas e atualizações dinâmicas.               \\
		Administrador de dados          & Integrante da equipe responsável por executar as rotinas de carga (ETL), validar as bases importadas e publicar novas versões dos indicadores.                                            \\
		Sistemas externos               & Serviços de CEP e geocodificação, servidores WFS do GeoSampa, portais de dados abertos da SSP-SP e do IBGE. Impõem limites de uso, licenças e obrigações de atribuição.                   \\
		\bottomrule
	\end{tabularx}
	\fonte{Elaborado pelo autor (2026).}
\end{table}

\subsection{Modelo de indicadores}

A plataforma calcula três indicadores em duas escalas espaciais. Na escala do endereço, o Índice de Proximidade Local mede a distância do ponto consultado aos equipamentos públicos mais próximos. Na escala do distrito, o Índice de Vitalidade do Espaço Público mede a oferta de infraestrutura da região, e os indicadores oficiais de ocorrências descrevem os registros criminais. Por decisão de projeto, os indicadores de infraestrutura e os de ocorrências são exibidos de forma paralela e nunca fundidos em uma nota única de ``segurança'' ou ``periculosidade''. Essa separação atende à crítica de \citeonline{caldeira2000} à ``fala do crime'' e evita que a ferramenta reproduza os mapas de perigo que alimentam a segregação.

Os três indicadores compartilham a mesma classificação de equipamentos, agrupados em seis categorias: (i) lazer e convivência (praças, parques e áreas verdes); (ii) cultura e esporte (bibliotecas, centros culturais, CEUs, clubes-escola e equipamentos esportivos); (iii) mobilidade (pontos e terminais de ônibus, estações de metrô e trem, ciclovias); (iv) comércio de rua de uso público (feiras livres e mercados municipais); (v) serviços essenciais (unidades de saúde e escolas públicas); e (vi) conectividade e permanência (pontos de Wi-Fi público e, quando disponível, iluminação pública). Cada categoria reúne um ou mais tipos de equipamento, como unidade básica de saúde, escola pública, praça, parque, ponto de ônibus ou estação de metrô e trem.

O \textbf{Índice de Proximidade Local} (IPL) é o indicador específico do endereço consultado. Para cada tipo de equipamento $j$, o sistema calcula a distância $d_j$, em metros, entre o ponto obtido pelo CEP e o equipamento mais próximo daquele tipo, por meio de busca de vizinho mais próximo no banco geoespacial (operador de distância indexado do PostGIS sobre o tipo \emph{geography}). A distância é convertida em uma nota de 0 a 100 por uma função linear decrescente, definida por dois limiares próprios de cada tipo: a distância ideal $a_j$, até a qual a nota é máxima, e a distância máxima $b_j$, a partir da qual a nota é nula (Equação~\ref{eq:nota}):

\begin{equation}
	\label{eq:nota}
	s_j =
	\begin{cases}
		100,                                    & \text{se } d_j \leq a_j    \\[4pt]
		100 \cdot \dfrac{b_j - d_j}{b_j - a_j}, & \text{se } a_j < d_j < b_j \\[8pt]
		0,                                      & \text{se } d_j \geq b_j
	\end{cases}
\end{equation}

A nota da categoria $k$, $S_k$, é a média das notas dos tipos que a compõem, e o IPL do endereço $c$ é a média ponderada das categorias, com os mesmos pesos $w_k$ ajustáveis pelo usuário (Equação~\ref{eq:ipl}):

\begin{equation}
	\label{eq:ipl}
	\text{IPL}_c = \frac{\sum_{k=1}^{6} w_k \, S_k}{\sum_{k=1}^{6} w_k}
\end{equation}

Os limiares iniciais, apresentados no Quadro~\ref{qua:limiares}, tomam como referência a distância de 400 metros, percorrida em cerca de cinco minutos de caminhada, e são tratados como parâmetros de configuração, a serem revistos com base nos testes com os moradores. Além da nota, o sistema informa para cada tipo a distância em linha reta, o tempo estimado de caminhada e a quantidade de equipamentos a até 500 metros e a até 1 quilômetro.

\begin{table}[htb]
	\caption{Limiares iniciais de distância por tipo de equipamento}
	\label{qua:limiares}
	\centering
	\small
	\begin{tabularx}{\textwidth}{L{3.6cm}XC{2.2cm}C{2.2cm}}
		\toprule
		\textbf{Categoria}  & \textbf{Tipos de equipamento}                               & \textbf{Ideal $a_j$ (m)} & \textbf{Máxima $b_j$ (m)} \\
		\midrule
		Lazer e convivência & Praça; parque ou área verde                                 & 400                      & 1\,200                    \\
		Cultura e esporte   & Biblioteca ou centro cultural; CEU ou equipamento esportivo & 800                      & 2\,000                    \\
		Mobilidade          & Ponto de ônibus                                             & 300                      & 800                       \\
		Mobilidade          & Estação de metrô ou trem; ciclovia                          & 800                      & 2\,000                    \\
		Comércio de rua     & Feira livre ou mercado municipal                            & 800                      & 1\,500                    \\
		Serviços essenciais & Unidade básica de saúde; escola pública                     & 800                      & 2\,000                    \\
		Conectividade       & Ponto de Wi-Fi público                                      & 400                      & 1\,200                    \\
		\bottomrule
	\end{tabularx}
	\fonte{Elaborado pelo autor (2026).}
\end{table}

O \textbf{Índice de Vitalidade do Espaço Público} (IVEP) complementa o IPL na escala do distrito, indicando se a região como um todo oferece infraestrutura que atrai uso cotidiano do espaço. Para cada categoria $k$ e distrito $d$, calcula-se a densidade de equipamentos por quilômetro quadrado e por dez mil habitantes, convertida em posição percentil entre os 96 distritos ($P_{k,d}$, de 0 a 100). O índice é a média ponderada dos percentis, conforme a Equação~\ref{eq:ivep}:

\begin{equation}
	\label{eq:ivep}
	\text{IVEP}_d = \frac{\sum_{k=1}^{6} w_k \, P_{k,d}}{\sum_{k=1}^{6} w_k}
\end{equation}

\noindent em que $w_k$ é o peso da categoria $k$, igual a 1 por padrão e ajustável pelo usuário entre 0 e 5, de modo que cada pessoa possa refletir suas prioridades de rotina.

O terceiro indicador compõe os \textbf{indicadores oficiais de ocorrências}, estimados com base nos microdados georreferenciados provenientes de boletins de ocorrência divulgados pela SSP-SP \cite{ssp2025}, correlacionados aos distritos mediante operação espacial do tipo ponto-em-polígono. Em conformidade com as orientações do manual de interpretação elaborado pela própria secretaria \cite{ssp2005}, que ressalta a subnotificação e advertência quanto a comparações entre unidades de dimensões distintas, os valores são expressos como taxa anualizada por cem mil habitantes, discriminados por natureza da ocorrência (roubo, furto, roubo de celular, roubo de veículo e crimes contra a vida), calculados em janela móvel trienal, conforme a Equação~\ref{eq:taxa}:

\begin{equation}
\label{eq:taxa}
T_{n,d} = \frac{O_{n,d}}{\text{Pop}_d \times a} \times 100\,000
\end{equation}

\noindent em que $O_{n,d}$ denota o quantitativo de ocorrências da natureza $n$ localizadas no distrito $d$ ao longo do período considerado, $\text{Pop}_d$ representa a população do distrito conforme dados do Censo Demográfico 2022 \cite{ibge2022} e $a$ corresponde ao número de anos abrangidos pela janela. Cada taxa é disponibilizada acompanhada da média municipal, da tendência observada no período e de notificação permanente sobre a questão da subnotificação.

\subsection{Fontes de dados e integrações}

O GeoHabita consome exclusivamente fontes abertas ou gratuitas, organizadas em dois fluxos. As consultas de endereço ocorrem em tempo real e passam por camada de cache; as bases estatísticas e cartográficas são importadas por rotinas periódicas de extração, transformação e carga (ETL) e versionadas com data de referência. O Quadro~\ref{qua:fontes} sintetiza as integrações previstas.

\begin{longtable}{L{3.1cm}L{4.3cm}L{2.3cm}L{4.4cm}}
	\caption{Fontes de dados e integrações}
	\label{qua:fontes}                                                                                                                                                                                                                                \\
	\toprule
	\textbf{Fonte}                        & \textbf{Uso no sistema}                                                                              & \textbf{Acesso}                    & \textbf{Restrições}                                           \\
	\midrule
	\endfirsthead
	\toprule
	\textbf{Fonte}                        & \textbf{Uso no sistema}                                                                              & \textbf{Acesso}                    & \textbf{Restrições}                                           \\
	\midrule
	\endhead
	\bottomrule
	\endfoot
	ViaCEP                                & Validação do CEP e obtenção de logradouro, bairro e código IBGE                                      & API REST, tempo real               & Não retorna coordenadas; bloqueio em uso massivo              \\
	BrasilAPI (CEP v2) e AwesomeAPI       & Contingência do ViaCEP e obtenção de latitude e longitude                                            & API REST, tempo real               & Coordenadas aproximadas e eventualmente ausentes              \\
	Nominatim (OpenStreetMap)             & Geocodificação de endereço completo quando o CEP não traz coordenadas                                & API REST, tempo real               & Até 1 requisição por segundo; cache e atribuição obrigatórios \\
	GeoSampa                              & Polígonos dos 96 distritos e camadas de equipamentos públicos, áreas verdes, transporte e mobilidade & WFS, carga periódica               & Sistema de referência SIRGAS 2000/UTM 23S; exige reprojeção   \\
	Portal de dados abertos da Prefeitura & Feiras livres, pontos de Wi-Fi público e demais cadastros municipais complementares                  & Arquivos CSV/JSON, carga periódica & Coordenadas a validar; atualização irregular                  \\
	SSP-SP                                & Microdados de boletins de ocorrência com coordenadas                                                 & Planilhas anuais, carga periódica  & Colunas variam entre anos; coordenadas nulas; subnotificação  \\
	IBGE (Censo 2022)                     & População por distrito para cálculo de taxas e densidades                                            & Arquivos e API, carga periódica    & Atualização decenal                                           \\
	OpenStreetMap e Leaflet               & Base cartográfica do mapa interativo                                                                 & Biblioteca JavaScript              & Atribuição obrigatória (licença ODbL)                         \\
\end{longtable}
\vspace{-0.4cm}
\fonte{Elaborado pelo autor (2026).}

\subsection{Requisitos funcionais}

Os requisitos funcionais descrevem os comportamentos observáveis do sistema \cite{sommerville2018}. Foram identificados 25 requisitos, agrupados em cinco módulos, dos quais 14 são classificados como essenciais (\emph{Must}) e compõem o MVP (Quadro~\ref{qua:rf}).

\begin{longtable}{C{1.2cm}L{9.2cm}C{1.5cm}C{2.2cm}}
	\caption{Requisitos funcionais}
	\label{qua:rf}                                                                                                                                                                                                                                \\
	\toprule
	\textbf{ID} & \textbf{Descrição}                                                                                                                                                                          & \textbf{Prior.} & \textbf{Módulo} \\
	\midrule
	\endfirsthead
	\toprule
	\textbf{ID} & \textbf{Descrição}                                                                                                                                                                          & \textbf{Prior.} & \textbf{Módulo} \\
	\midrule
	\endhead
	\bottomrule
	\endfoot
	RF01        & O sistema deve permitir a consulta por CEP, validando o formato de oito dígitos antes de qualquer chamada externa.                                                                          & Must            & Consulta        \\
	RF02        & O sistema deve preencher automaticamente logradouro, bairro e município a partir do CEP informado.                                                                                          & Must            & Consulta        \\
	RF03        & O sistema deve permitir a consulta por endereço textual quando o usuário não souber o CEP.                                                                                                  & Should          & Consulta        \\
	RF04        & O sistema deve converter o endereço em coordenadas geográficas, utilizando serviço de contingência em caso de falha do serviço principal.                                                   & Must            & Consulta        \\
	RF05        & O sistema deve identificar o distrito administrativo correspondente às coordenadas por operação de ponto-em-polígono.                                                                       & Must            & Consulta        \\
	RF06        & O sistema deve informar ao usuário quando o endereço estiver fora do município de São Paulo.                                                                                                & Must            & Consulta        \\
	RF07        & O sistema deve listar os equipamentos públicos de uso do espaço em raio de 500 metros e de 1 quilômetro do ponto consultado, agrupados pelas seis categorias do IVEP, com nome e distância. & Must            & Diagnóstico     \\
	RF23        & O sistema deve calcular, para o endereço consultado, a distância em linha reta até o equipamento mais próximo de cada tipo, com o tempo estimado de caminhada.                              & Must            & Diagnóstico     \\
	RF24        & O sistema deve calcular e exibir o Índice de Proximidade Local do endereço e a nota de cada categoria, destacando os tipos de equipamento ausentes no raio máximo.                          & Must            & Diagnóstico     \\
	RF25        & O sistema deve calcular a distância e o tempo pela rede de calçadas e vias até o equipamento mais próximo de cada tipo, quando o usuário solicitar.                                         & Could           & Diagnóstico     \\
	RF08        & O sistema deve exibir o IVEP do distrito, o percentil de cada categoria e a comparação com a mediana do município.                                                                          & Must            & Diagnóstico     \\
	RF09        & O sistema deve permitir que o usuário ajuste os pesos das categorias do IPL e do IVEP e recalcule os índices sem recarregar a página.                                                       & Should          & Diagnóstico     \\
	RF10        & O sistema deve exibir as taxas de ocorrências por cem mil habitantes do distrito, por natureza, com média municipal e tendência dos últimos três anos.                                      & Must            & Diagnóstico     \\
	RF11        & O sistema deve apresentar, junto às taxas criminais, nota explicativa sobre subnotificação e sobre a metodologia de cálculo.                                                                & Must            & Diagnóstico     \\
	RF12        & O sistema deve exibir a fonte e a data de referência de cada indicador apresentado.                                                                                                         & Must            & Diagnóstico     \\
	RF13        & O sistema deve exibir mapa interativo com o ponto consultado, o limite do distrito e os equipamentos do entorno, filtráveis por categoria.                                                  & Must            & Visualização    \\
	RF14        & O sistema deve oferecer, para todo conteúdo cartográfico, tabela textual equivalente navegável por teclado e leitor de tela.                                                                & Must            & Visualização    \\
	RF15        & O sistema deve anunciar por região viva (\emph{live region}) WAI-ARIA a conclusão de cada consulta assíncrona e eventuais erros.                                                            & Should          & Visualização    \\
	RF16        & O sistema deve permitir a comparação lado a lado de até três endereços ou distritos.                                                                                                        & Should          & Comparação      \\
	RF17        & O sistema deve permitir compartilhar um diagnóstico por link, sem armazenar dados pessoais do consulente.                                                                                   & Could           & Comparação      \\
	RF18        & O sistema deve permitir exportar o diagnóstico em formato PDF acessível ou CSV.                                                                                                             & Could           & Comparação      \\
	RF19        & O sistema deve permitir ao administrador executar as rotinas de carga das bases externas e registrar o resultado de cada execução.                                                          & Should          & Administração   \\
	RF20        & O sistema deve validar as bases importadas, descartando registros sem coordenadas ou fora dos limites do município e reportando a quantidade descartada.                                    & Should          & Administração   \\
	RF21        & O sistema deve manter o histórico de versões dos indicadores com data de referência de cada carga.                                                                                          & Could           & Administração   \\
	RF22        & O sistema deve disponibilizar uma API pública de leitura dos indicadores agregados por distrito, documentada em OpenAPI.                                                                    & Could           & Administração   \\
\end{longtable}
\vspace{-0.4cm}
\fonte{Elaborado pelo autor (2026).}

\subsection{Requisitos não funcionais}

Os requisitos não funcionais estabelecem restrições sobre os serviços oferecidos e sobre o processo de desenvolvimento \cite{sommerville2018}. Para garantir sua verificabilidade, cada requisito de produto está associado a uma métrica objetiva, conforme recomendam \citeonline{pressman2021}. Os requisitos de produto estão classificados segundo as características de qualidade da ISO/IEC 25010 (Quadro~\ref{qua:rnf}).

\begin{longtable}{C{1.3cm}L{2.6cm}L{6.1cm}L{4.1cm}}
	\caption{Requisitos não funcionais de produto}
	\label{qua:rnf}                                                                                                                                                                                                                                                                                                                                             \\
	\toprule
	\textbf{ID} & \textbf{Característica}              & \textbf{Requisito}                                                                                                                                                                                    & \textbf{Métrica de verificação}                                                                \\
	\midrule
	\endfirsthead
	\toprule
	\textbf{ID} & \textbf{Característica}              & \textbf{Requisito}                                                                                                                                                                                    & \textbf{Métrica de verificação}                                                                \\
	\midrule
	\endhead
	\bottomrule
	\endfoot
	RNF01       & Eficiência de desempenho             & O diagnóstico de um CEP já armazenado em cache deve ser exibido em até 2 segundos.                                                                                                                    & Percentil 95 $\leq$ 2 s com 30 usuários simultâneos                                            \\
	RNF02       & Eficiência de desempenho             & O diagnóstico de um CEP inédito deve ser exibido em até 6 segundos.                                                                                                                                   & Percentil 95 $\leq$ 6 s em teste de carga                                                      \\
	RNF03       & Usabilidade (acessibilidade)         & A interface deve atender às diretrizes WCAG 2.1 no nível AA e empregar marcações WAI-ARIA em todos os componentes dinâmicos.                                                                          & Zero erros críticos em Lighthouse e axe; contraste $\geq$ 4,5:1                                \\
	RNF04       & Usabilidade (operabilidade)          & Todas as funções devem ser operáveis exclusivamente por teclado, com foco visível e ordem lógica.                                                                                                     & Roteiro de teste manual com NVDA e Orca                                                        \\
	RNF05       & Usabilidade (inteligibilidade)       & Os textos explicativos devem usar linguagem simples, sem jargão técnico ou geográfico sem explicação.                                                                                                 & Avaliação em teste com usuários; SUS $\geq$ 70                                                 \\
	RNF06       & Confiabilidade (tolerância a falhas) & A falha de uma fonte externa não deve impedir a exibição dos demais blocos do diagnóstico.                                                                                                            & Teste com serviço externo simulado fora do ar                                                  \\
	RNF07       & Confiabilidade (disponibilidade)     & A aplicação deve manter disponibilidade mensal mínima de 98\%.                                                                                                                                        & Monitoramento externo de disponibilidade                                                       \\
	RNF08       & Segurança                            & Todo o tráfego deve ocorrer por HTTPS, e as rotas públicas devem limitar-se a 60 requisições por minuto por endereço IP.                                                                              & Teste automatizado; resposta HTTP 429 acima do limite                                          \\
	RNF09       & Segurança                            & A aplicação não deve apresentar vulnerabilidades de alta severidade do OWASP Top 10.                                                                                                                  & Varredura com OWASP ZAP e auditoria de dependências                                            \\
	RNF10       & Manutenibilidade                     & O código deve adotar arquitetura em camadas (apresentação, aplicação, domínio e infraestrutura), com adaptadores isolando cada serviço externo.                                                       & Revisão em \emph{pull request}; ausência de chamadas externas fora da camada de infraestrutura \\
	RNF11       & Manutenibilidade (testabilidade)     & As regras de negócio devem ter cobertura mínima de 70\% por testes unitários.                                                                                                                         & Relatório de cobertura na integração contínua                                                  \\
	RNF12       & Compatibilidade e portabilidade      & A interface deve funcionar nas duas versões mais recentes dos navegadores Chrome, Firefox, Edge e Safari, em telas a partir de 360 pixels de largura, e a aplicação deve ser empacotada em contêiner. & Roteiro multiplataforma; implantação em nuvem com um único comando                             \\
\end{longtable}
\vspace{-0.4cm}
\fonte{Elaborado pelo autor (2026).}

Os requisitos organizacionais e externos, que não se referem diretamente ao comportamento do produto, estão reunidos no Quadro~\ref{qua:rnf-org}.

\begin{table}[htb]
	\caption{Requisitos não funcionais organizacionais e externos}
	\label{qua:rnf-org}
	\centering
	\small
	\begin{tabularx}{\textwidth}{C{1.3cm}C{2.4cm}X}
		\toprule
		\textbf{ID} & \textbf{Tipo}      & \textbf{Requisito}                                                                                                                                                                                                         \\
		\midrule
		RNF13       & Organizacional     & O desenvolvimento deve seguir práticas do Scrum, com \emph{sprints} de duas semanas, controle de versão em Git e integração contínua executando os testes a cada envio ao repositório.                                     \\
		RNF14       & Organizacional     & A pilha tecnológica deve compreender Python com micro-framework web no serviço de retaguarda, JavaScript no cliente, PostgreSQL com extensão PostGIS na persistência e hospedagem em provedor de nuvem com plano gratuito. \\
		RNF15       & Externo (legal)    & O tratamento de dados deve observar a Lei Geral de Proteção de Dados \cite{lgpd2018}: a consulta não exige cadastro e nenhum CEP consultado é associado a pessoa identificável.                                            \\
		RNF16       & Externo (legal)    & A plataforma deve atender ao art. 63 da Lei Brasileira de Inclusão \cite{lbi2015}, que torna obrigatória a acessibilidade em sítios da internet.                                                                           \\
		RNF17       & Externo (licenças) & Todas as fontes devem ser atribuídas conforme suas licenças, e o código deve ser distribuído sob a GNU AGPL v3 \cite{agplv3}.                                                                                              \\
		RNF18       & Externo (ético)    & Os indicadores não devem empregar rótulos estigmatizantes, rankings de ``bairros perigosos'' nem escalas cromáticas de alarme para dados criminais.                                                                        \\
		\bottomrule
	\end{tabularx}
	\fonte{Elaborado pelo autor (2026).}
\end{table}

\subsection{Regras de negócio e requisitos de domínio}

As regras de negócio condensam o conhecimento do domínio urbano e estatístico que o sistema deve respeitar independentemente da interface. Elas são implementadas na camada de domínio e testadas de forma isolada (Quadro~\ref{qua:rn}).

\begin{longtable}{C{1.2cm}L{10.6cm}L{2.6cm}}
	\caption{Regras de negócio}
	\label{qua:rn}                                                                                                                                                                                                                                                                                               \\
	\toprule
	\textbf{ID} & \textbf{Regra}                                                                                                                                                                                                                                                        & \textbf{Relacionada a} \\
	\midrule
	\endfirsthead
	\toprule
	\textbf{ID} & \textbf{Regra}                                                                                                                                                                                                                                                        & \textbf{Relacionada a} \\
	\midrule
	\endhead
	\bottomrule
	\endfoot
	RN01        & Um CEP válido possui exatamente oito dígitos numéricos; hífens, pontos e espaços são removidos antes da validação.                                                                                                                                                    & RF01                   \\
	RN02        & Um CEP inexistente gera a mensagem ``CEP não encontrado'' e libera a consulta por endereço textual.                                                                                                                                                                   & RF02, RF03             \\
	RN03        & O serviço de contingência é acionado quando o serviço principal retorna erro ou excede 3 segundos de resposta.                                                                                                                                                        & RF04, RNF06            \\
	RN04        & Resultados de CEP e geocodificação permanecem em cache por 30 dias; indicadores agregados, até a próxima carga.                                                                                                                                                       & RF04, RNF01            \\
	RN05        & A coordenada obtida pelo CEP é tratada como aproximada; a interface informa que ela identifica o distrito, mas não a posição exata do imóvel.                                                                                                                         & RF05, RF07             \\
	RN06        & O IVEP é calculado por percentis entre os 96 distritos e varia de 0 a 100; os pesos padrão são iguais a 1 e aceitam valores inteiros de 0 a 5.                                                                                                                        & RF08, RF09             \\
	RN07        & Os indicadores de vitalidade e os indicadores de ocorrências nunca são combinados em uma pontuação única.                                                                                                                                                             & RF08, RF10, RNF18      \\
	RN08        & As taxas de ocorrências usam janela de três anos completos e são expressas por cem mil habitantes; o ano corrente incompleto não entra no cálculo.                                                                                                                    & RF10                   \\
	RN09        & Naturezas com menos de 30 ocorrências no distrito durante a janela são exibidas como ``número insuficiente para comparação'', sem taxa nem tendência.                                                                                                                 & RF10, RF11             \\
	RN10        & Registros da SSP-SP com coordenadas nulas, iguais a zero ou fora do município são descartados na carga e contabilizados no relatório de validação.                                                                                                                    & RF20                   \\
	RN11        & Indicadores com data de referência superior a 24 meses recebem o aviso ``dado desatualizado''.                                                                                                                                                                        & RF12                   \\
	RN12        & A distância até o equipamento mais próximo é a distância geodésica em linha reta, em metros, entre o ponto do CEP e o ponto do equipamento; para áreas como parques, usa-se o ponto mais próximo do contorno. A interface identifica a medida como ``em linha reta''. & RF23                   \\
	RN13        & A nota de cada tipo segue a Equação~\ref{eq:nota}, com os limiares do Quadro~\ref{qua:limiares}; os limiares são parâmetros de configuração alteráveis pelo administrador sem mudança de código.                                                                      & RF24                   \\
	RN14        & O tempo estimado de caminhada é calculado à velocidade de 80 metros por minuto (4,8 km/h) e arredondado para o minuto inteiro superior.                                                                                                                               & RF23                   \\
	RN15        & Tipo sem nenhum equipamento a menos da distância máxima recebe nota zero, e a interface informa ``nenhum equipamento deste tipo a menos de $b_j$ metros''.                                                                                                            & RF24                   \\
	RN16        & O IPL é calculado sempre sobre a coordenada do endereço consultado, e nunca sobre o centroide do distrito ou do bairro.                                                                                                                                               & RF24                   \\
\end{longtable}
\vspace{-0.4cm}
\fonte{Elaborado pelo autor (2026).}

A regra RN09 transpõe para o sistema a advertência do manual de interpretação da SSP-SP sobre comparações apoiadas em poucos casos, e a regra RN07 materializa a opção, discutida na fundamentação, de não reduzir a qualidade de um território a um único número.

\subsection{Casos de uso}

Os requisitos funcionais foram agrupados em oito casos de uso principais, que orientam o diagrama de casos de uso e o desenho das telas:

\begin{alineas}
	\item \textbf{UC01 -- Consultar endereço}: o munícipe informa CEP ou endereço, e o sistema valida, geocodifica e identifica o distrito (RF01 a RF06);
	\item \textbf{UC02 -- Visualizar diagnóstico socioespacial}: o sistema apresenta equipamentos do entorno, IVEP e indicadores de ocorrências, com fontes e datas (RF07, RF08, RF10 a RF12);
	\item \textbf{UC02a -- Avaliar proximidade do endereço}: o sistema calcula a distância até o equipamento mais próximo de cada tipo e o IPL do endereço (RF23 a RF25);
	\item \textbf{UC03 -- Personalizar prioridades}: o munícipe ajusta os pesos das categorias e acompanha o recálculo do IVEP (RF09);
	\item \textbf{UC04 -- Explorar mapa ou tabela equivalente}: o munícipe navega pelo mapa ou pela tabela textual e filtra categorias (RF13 a RF15);
	\item \textbf{UC05 -- Comparar localidades}: o munícipe compara até três endereços ou distritos (RF16);
	\item \textbf{UC06 -- Compartilhar ou exportar diagnóstico}: o munícipe gera link ou arquivo do diagnóstico (RF17, RF18);
	\item \textbf{UC07 -- Atualizar bases de dados}: o administrador executa as cargas, valida os registros e publica nova versão dos indicadores (RF19 a RF21).
\end{alineas}

\subsection{Histórias de usuário e critérios de aceitação}

Para o planejamento das \emph{sprints}, os casos de uso foram desdobrados em histórias de usuário no formato ``Como \emph{papel}, quero \emph{ação}, para \emph{benefício}'' \cite{cohn2004}, estimadas em pontos de complexidade relativa (Quadro~\ref{qua:hu}).

\begin{longtable}{C{1.2cm}L{9.0cm}C{2.1cm}C{1.3cm}}
	\caption{Histórias de usuário}
	\label{qua:hu}                                                                                                                                                                                                                             \\
	\toprule
	\textbf{ID} & \textbf{História}                                                                                                                                                                    & \textbf{Requisitos} & \textbf{Pontos} \\
	\midrule
	\endfirsthead
	\toprule
	\textbf{ID} & \textbf{História}                                                                                                                                                                    & \textbf{Requisitos} & \textbf{Pontos} \\
	\midrule
	\endhead
	\bottomrule
	\endfoot
	HU01        & Como munícipe, quero digitar um CEP e ver o endereço reconhecido, para iniciar a consulta sem preencher formulários longos.                                                          & RF01, RF02          & 3               \\
	HU02        & Como munícipe, quero saber em que distrito fica o endereço, para entender a região a que ele pertence.                                                                               & RF04--RF06          & 5               \\
	HU03        & Como munícipe, quero ver praças, parques, feiras, transporte e serviços públicos a até 1 km, para saber se a vizinhança tem vida nas ruas.                                           & RF07, RF13          & 8               \\
	HU04        & Como munícipe, quero ver o índice de vitalidade do distrito em comparação com a cidade, para avaliar a região de forma objetiva.                                                     & RF08                & 5               \\
	HU05        & Como munícipe, quero ver as taxas oficiais de ocorrências com explicação sobre seus limites, para decidir com informação confiável e sem alarmismo.                                  & RF10--RF12          & 8               \\
	HU06        & Como usuário de leitor de tela, quero acessar todo o conteúdo do mapa em forma de tabela e ser avisado quando a consulta terminar, para usar a plataforma com autonomia.             & RF14, RF15          & 5               \\
	HU07        & Como munícipe, quero dar mais peso ao transporte do que ao lazer, para que o índice reflita minha rotina.                                                                            & RF09                & 5               \\
	HU08        & Como munícipe, quero comparar três endereços lado a lado, para escolher entre as opções finais.                                                                                      & RF16                & 8               \\
	HU09        & Como administrador, quero carregar e validar novas bases da SSP-SP e do GeoSampa, para manter os indicadores atualizados.                                                            & RF19, RF20          & 8               \\
	HU10        & Como munícipe, quero saber a que distância do endereço ficam a UBS, a escola, a praça, o ponto de ônibus e a estação mais próximos, para avaliar se consigo fazer minha rotina a pé. & RF23, RF24          & 8               \\
\end{longtable}
\vspace{-0.4cm}
\fonte{Elaborado pelo autor (2026).}

Os critérios de aceitação são redigidos na sintaxe Gherkin (Dado, Quando, Então), o que permite convertê-los diretamente em testes automatizados de aceitação. A título de exemplo, a história HU01 é aceita quando todos os cenários a seguir forem atendidos:

\begin{alineas}
	\item \textbf{CEP válido}: \emph{dado} que o munícipe está na página inicial, \emph{quando} informa ``01310-100'', \emph{então} o sistema exibe ``Avenida Paulista, Bela Vista, São Paulo -- SP'';
	\item \textbf{Formato inválido}: \emph{quando} informa ``1234'', \emph{então} o sistema exibe ``O CEP deve ter 8 dígitos'' e nenhuma API externa é chamada;
	\item \textbf{CEP inexistente}: \emph{quando} informa ``99999-999'', \emph{então} o sistema exibe ``CEP não encontrado'' e oferece a consulta por endereço;
	\item \textbf{Serviço principal indisponível}: \emph{dado} que o ViaCEP não responde em 3 segundos, \emph{quando} o munícipe informa um CEP válido, \emph{então} o endereço é obtido pelo serviço de contingência.
\end{alineas}

Considera-se uma história pronta para desenvolvimento (\emph{Definition of Ready}) quando possui critérios de aceitação escritos, estimativa acordada e dependências externas identificadas. Considera-se concluída (\emph{Definition of Done}) quando o código foi revisado, os testes unitários e de aceitação passam na integração contínua, a verificação de acessibilidade não aponta erros críticos e o incremento está implantado no ambiente em nuvem.

\subsection{Priorização e planejamento incremental}

As dez histórias totalizam 63 pontos, distribuídos em quatro \emph{sprints} de duas semanas, precedidas de uma \emph{sprint} de preparação do ambiente (Quadro~\ref{qua:sprints}). O MVP encerra-se ao fim da terceira \emph{sprint}, quando o fluxo completo de consulta e diagnóstico estará disponível para validação com os moradores.

\begin{table}[htb]
	\caption{Planejamento das \emph{sprints}}
	\label{qua:sprints}
	\centering
	\small
	\begin{tabularx}{\textwidth}{C{1.4cm}XC{3.2cm}C{1.4cm}}
		\toprule
		\textbf{Sprint} & \textbf{Meta}                                                                                  & \textbf{Histórias} & \textbf{Pontos} \\
		\midrule
		0               & Repositório, integração contínua, banco com PostGIS e primeira carga dos distritos do GeoSampa & --                 & --              \\
		1               & Consultar um CEP e localizar o distrito                                                        & HU01, HU02, HU09   & 16              \\
		2               & Exibir a infraestrutura do entorno e a distância ao equipamento mais próximo de cada tipo      & HU03, HU10         & 16              \\
		3               & Calcular IVEP e indicadores de ocorrências, com mapa e tabela acessível (MVP)                  & HU04, HU05, HU06   & 18              \\
		4               & Personalização de pesos, comparação e ajustes após testes com usuários                         & HU07, HU08         & 13              \\
		\bottomrule
	\end{tabularx}
	\fonte{Elaborado pelo autor (2026).}
\end{table}

\subsection{Plano de testes}

A estratégia de verificação segue a pirâmide de testes: grande volume de testes unitários sobre as regras de negócio, testes de integração sobre os adaptadores de dados e o banco geoespacial, e um conjunto reduzido de testes ponta a ponta sobre os fluxos críticos \cite{pressman2021}. Nos testes automatizados, os serviços externos são substituídos por simulações (\emph{mocks}), o que torna os resultados reprodutíveis; testes de contrato agendados verificam, separadamente, se o formato das respostas reais das APIs permanece inalterado. O Quadro~\ref{qua:niveis} apresenta os níveis de teste previstos.

\begin{longtable}{L{2.6cm}L{5.0cm}L{3.5cm}L{3.0cm}}
	\caption{Níveis e tipos de teste}
	\label{qua:niveis}                                                                                                                                                                       \\
	\toprule
	\textbf{Nível} & \textbf{Objeto de verificação}                                  & \textbf{Ferramentas}                           & \textbf{Critério de saída}                           \\
	\midrule
	\endfirsthead
	\toprule
	\textbf{Nível} & \textbf{Objeto de verificação}                                  & \textbf{Ferramentas}                           & \textbf{Critério de saída}                           \\
	\midrule
	\endhead
	\bottomrule
	\endfoot
	Unitário       & Regras RN01 a RN16, cálculo do IPL, do IVEP e das taxas         & pytest                                         & 100\% aprovados; cobertura $\geq$ 70\%               \\
	Integração     & Adaptadores de CEP e geocodificação, consultas espaciais, cache & pytest com banco PostGIS em contêiner          & 100\% aprovados                                      \\
	Contrato       & Formato das respostas de ViaCEP, BrasilAPI e Nominatim          & Execução diária agendada                       & Alerta em caso de divergência                        \\
	Aceitação      & Cenários Gherkin das histórias de usuário                       & pytest-bdd ou Behave                           & Todas as histórias \emph{Must} aprovadas             \\
	Ponta a ponta  & Fluxo consulta, diagnóstico, mapa e tabela                      & Playwright                                     & Aprovado em Chrome e Firefox                         \\
	Desempenho     & RNF01 e RNF02                                                   & k6                                             & Metas de percentil 95 atendidas                      \\
	Segurança      & RNF08 e RNF09                                                   & OWASP ZAP; pip-audit                           & Nenhum alerta alto                                   \\
	Acessibilidade & RNF03 e RNF04                                                   & Lighthouse, axe, NVDA, Orca                    & Zero erros críticos                                  \\
	Usabilidade    & RNF05 e aderência às necessidades levantadas                    & Teste moderado com moradores; questionário SUS & Tarefas concluídas por $\geq$ 80\% dos participantes \\
\end{longtable}
\vspace{-0.4cm}
\fonte{Elaborado pelo autor (2026).}

Os casos de teste derivam diretamente dos requisitos e das regras de negócio. O Quadro~\ref{qua:ct} relaciona os principais.

\begin{longtable}{C{1.1cm}C{1.8cm}L{4.6cm}L{4.4cm}C{1.8cm}}
	\caption{Casos de teste principais}
	\label{qua:ct}                                                                                                                                                                                \\
	\toprule
	\textbf{ID} & \textbf{Requisito} & \textbf{Entrada ou condição}                                 & \textbf{Resultado esperado}                                                & \textbf{Tipo}  \\
	\midrule
	\endfirsthead
	\toprule
	\textbf{ID} & \textbf{Requisito} & \textbf{Entrada ou condição}                                 & \textbf{Resultado esperado}                                                & \textbf{Tipo}  \\
	\midrule
	\endhead
	\bottomrule
	\endfoot
	CT01        & RF01, RN01         & ``1234'', ``abcdefgh'' e ``123456789''                       & Mensagem de formato inválido e nenhuma chamada externa                     & Unitário       \\
	CT02        & RF02               & CEP 01310-100                                                & Logradouro, bairro e município preenchidos corretamente                    & Integração     \\
	CT03        & RN02               & CEP 99999-999                                                & ``CEP não encontrado'' e consulta textual liberada                         & Integração     \\
	CT04        & RN03, RNF06        & Serviço principal simulado com atraso de 5 s                 & Resposta obtida pelo serviço de contingência                               & Integração     \\
	CT05        & RF05               & Coordenadas conhecidas no interior de três distritos         & Distrito correto em todos os casos                                         & Integração     \\
	CT06        & RF06               & CEP de município vizinho                                     & Aviso de endereço fora de São Paulo                                        & Integração     \\
	CT07        & RF07               & Ponto com equipamentos cadastrados a 300 m, 800 m e 1,5 km   & Os dois primeiros listados com distância; o terceiro ausente               & Integração     \\
	CT08        & RF08, RN06         & Base fictícia com percentis conhecidos e pesos iguais        & IVEP igual à média aritmética dos percentis                                & Unitário       \\
	CT09        & RF09               & Peso de mobilidade igual a 5 e demais iguais a 1             & IVEP recalculado conforme a Equação~\ref{eq:ivep}, sem recarregar a página & Ponta a ponta  \\
	CT10        & RN08               & Distrito com 300 ocorrências em 3 anos e 100\,000 habitantes & Taxa de 100 por cem mil habitantes ao ano                                  & Unitário       \\
	CT11        & RN09               & Natureza com 12 ocorrências na janela                        & ``Número insuficiente para comparação'', sem taxa                          & Unitário       \\
	CT12        & RN07               & Resposta da API interna do diagnóstico                       & Ausência de qualquer campo que combine os dois conjuntos de indicadores    & Unitário       \\
	CT13        & RN10, RF20         & Planilha com 5 registros de coordenadas nulas ou zeradas     & Registros descartados e contabilizados no relatório de carga               & Integração     \\
	CT14        & RF14, RNF04        & Navegação apenas por teclado com leitor de tela              & Todos os equipamentos do mapa acessíveis pela tabela equivalente           & Acessibilidade \\
	CT15        & RF15               & Consulta concluída e consulta com erro                       & Ambas anunciadas pela região viva                                          & Acessibilidade \\
	CT16        & RNF08              & 61 requisições em um minuto a partir do mesmo IP             & A 61ª requisição recebe HTTP 429                                           & Segurança      \\
	CT17        & RF23, RN12         & Ponto com duas UBS cadastradas a 250 m e a 600 m             & UBS mais próxima informada a 250 m, ``em linha reta''                      & Integração     \\
	CT18        & RN13               & Ponto de ônibus a 550 m ($a_j = 300$, $b_j = 800$)           & Nota do tipo igual a 50                                                    & Unitário       \\
	CT19        & RN14               & Distância de 650 m                                           & Tempo estimado de 9 minutos                                                & Unitário       \\
	CT20        & RN15               & Nenhuma estação a menos de 2\,000 m                          & Nota zero e mensagem de ausência                                           & Unitário       \\
	CT21        & RF24               & Notas de categoria 100, 50, 0, 80, 60 e 40, pesos iguais a 1 & IPL igual a 55                                                             & Unitário       \\
	CT22        & RN16               & Dois CEPs distintos no mesmo distrito                        & IPL diferentes e mesmo IVEP                                                & Integração     \\
\end{longtable}
\vspace{-0.4cm}
\fonte{Elaborado pelo autor (2026).}

Os defeitos encontrados são registrados como \emph{issues} no repositório, classificados por severidade; defeitos críticos e altos impedem o encerramento da \emph{sprint}.

\subsection{Matriz de rastreabilidade}

A matriz de rastreabilidade (Quadro~\ref{qua:rastreab}) liga as necessidades identificadas junto à comunidade aos requisitos que as atendem e aos testes que comprovam seu cumprimento, permitindo avaliar o impacto de qualquer alteração de escopo.

\begin{longtable}{L{4.8cm}L{4.2cm}C{2.0cm}C{2.6cm}}
	\caption{Matriz de rastreabilidade}
	\label{qua:rastreab}                                                                                                          \\
	\toprule
	\textbf{Necessidade}                                    & \textbf{Requisitos e regras}  & \textbf{História} & \textbf{Testes} \\
	\midrule
	\endfirsthead
	\toprule
	\textbf{Necessidade}                                    & \textbf{Requisitos e regras}  & \textbf{História} & \textbf{Testes} \\
	\midrule
	\endhead
	\bottomrule
	\endfoot
	Consultar uma região de forma simples                   & RF01--RF06, RN01--RN05        & HU01, HU02        & CT01--CT06      \\
	Saber a que distância ficam os serviços do dia a dia    & RF23--RF25, RN12--RN16        & HU10              & CT17--CT22      \\
	Conhecer a infraestrutura pública do entorno            & RF07, RF13                    & HU03              & CT07            \\
	Avaliar a vitalidade do espaço público                  & RF08, RF09, RN06              & HU04, HU07        & CT08, CT09      \\
	Dispor de dados criminais auditáveis e contextualizados & RF10--RF12, RN07--RN09, RNF18 & HU05              & CT10--CT12      \\
	Usar a plataforma com tecnologia assistiva              & RF14, RF15, RNF03, RNF04      & HU06              & CT14, CT15      \\
	Manter os dados atualizados e íntegros                  & RF19, RF20, RN10              & HU09              & CT13            \\
	Operar com segurança e desempenho                       & RNF01, RNF02, RNF08           & --                & CT16            \\
\end{longtable}
\vspace{-0.4cm}
\fonte{Elaborado pelo autor (2026).}

% posicionado após a subseção "Metodologia" (ou dentro de
% "Resultados preliminares: solução inicial").
%
% Pacotes necessários no preâmbulo do relatório:
%   \usepackage{longtable}
%   \usepackage{tabularx}
%   \usepackage{booktabs}
%   \usepackage{array}
%   \usepackage{amsmath}
% Referências novas: ver requisitos_geohabita.bib
% =====================================================================

\counterwithin{section}{chapter}
\renewcommand{\thesection}{\thechapter.\arabic{section}}
\counterwithin{table}{chapter}
\renewcommand{\thetable}{\thechapter.\arabic{table}}
\counterwithin{equation}{chapter}
\renewcommand{\theequation}{\thechapter.\arabic{equation}}

\newcolumntype{L}[1]{>{\raggedright\arraybackslash}p{#1}}
\newcolumntype{C}[1]{>{\centering\arraybackslash}p{#1}}


\section{Especificação de requisitos}

A especificação de requisitos do GeoHabita consolida, em termos verificáveis, aquilo que a plataforma deve oferecer ao cidadão e as condições de qualidade sob as quais deve operar. Em relação à concepção inicial, o escopo foi reorientado: a segurança deixa de ser tratada prioritariamente a partir da estatística criminal e passa a ser lida também pela infraestrutura pública de uso do espaço, isto é, pelos equipamentos e lugares que atraem pessoas às ruas em horários variados e sustentam o que \citeonline[cap.~1]{jacobs2011} denomina vigilância natural dos ``olhos na rua''. Praças, parques, equipamentos culturais e esportivos, feiras livres, pontos de transporte coletivo, unidades de saúde e escolas passam a compor um indicador próprio de vitalidade do espaço público, apresentado lado a lado com os registros oficiais da SSP-SP \cite{ssp2024,geosampa2024}.

A estrutura adotada segue a norma ISO/IEC/IEEE 29148 \cite{iso29148} e a classificação proposta por \citeonline{sommerville2018}, que distingue requisitos funcionais, requisitos não funcionais (de produto, organizacionais e externos) e requisitos de domínio. Os atributos de qualidade são organizados segundo o modelo da ISO/IEC 25010 \cite{iso25010}, e o desdobramento em histórias de usuário e critérios de aceitação acompanha as práticas ágeis descritas por \citeonline{cohn2004} e \citeonline{pressman2021}. Cada requisito recebe identificador único, prioridade segundo a técnica MoSCoW (\emph{Must}, \emph{Should}, \emph{Could}, \emph{Won't}) e rastreabilidade até os casos de teste que o verificam.

\subsection{Escopo do produto e atores}

O Produto Mínimo Viável (MVP) delimita-se ao município de São Paulo e aos seus 96 distritos administrativos, unidade territorial escolhida por dispor de polígonos oficiais no GeoSampa e de população recenseada pelo IBGE, o que permite o cálculo de taxas comparáveis entre regiões. A consulta parte de um CEP ou de um endereço informado pelo usuário, que é convertido em coordenadas geográficas e associado ao distrito correspondente. A partir desse ponto, o sistema produz um diagnóstico socioespacial composto por três blocos: a proximidade do endereço consultado aos equipamentos públicos de uso do espaço, medida pela distância até o equipamento mais próximo de cada tipo; o indicador de vitalidade do espaço público do distrito; e os indicadores oficiais de ocorrências criminais, sempre contextualizados.

Permanecem fora do escopo do MVP a intermediação ou divulgação de anúncios imobiliários, a coleta colaborativa de relatos de crimes, a predição de ocorrências por modelos de aprendizado de máquina, a extração automatizada não estruturada de páginas (\emph{web scraping}) e aplicativos nativos para dispositivos móveis. A exclusão dos relatos colaborativos decorre, em particular, do risco de reforço de estigmas territoriais e de reidentificação de vítimas já discutido na fundamentação teórica.

Identificam-se três atores principais, descritos no Quadro~\ref{qua:atores}.

\begin{table}[htb]
	\caption{Atores do sistema}
	\label{qua:atores}
	\centering
	\small
	\begin{tabularx}{\textwidth}{L{3.6cm}X}
		\toprule
		\textbf{Ator}                   & \textbf{Descrição e interesse}                                                                                                                                                            \\
		\midrule
		Munícipe consulente             & Pessoa ou família em transição habitacional, ou morador interessado em conhecer o próprio território. Consulta endereços, compara regiões e não precisa de cadastro para o uso essencial. \\
		Usuário de tecnologia assistiva & Munícipe que acessa a plataforma por leitor de tela, navegação exclusiva por teclado ou ampliação. Exige equivalência textual integral para mapas e atualizações dinâmicas.               \\
		Administrador de dados          & Integrante da equipe responsável por executar as rotinas de carga (ETL), validar as bases importadas e publicar novas versões dos indicadores.                                            \\
		Sistemas externos               & Serviços de CEP e geocodificação, servidores WFS do GeoSampa, portais de dados abertos da SSP-SP e do IBGE. Impõem limites de uso, licenças e obrigações de atribuição.                   \\
		\bottomrule
	\end{tabularx}
	\fonte{Elaborado pelo autor (2026).}
\end{table}

\subsection{Modelo de indicadores}

A plataforma calcula três indicadores em duas escalas espaciais. Na escala do endereço, o Índice de Proximidade Local mede a distância do ponto consultado aos equipamentos públicos mais próximos. Na escala do distrito, o Índice de Vitalidade do Espaço Público mede a oferta de infraestrutura da região, e os indicadores oficiais de ocorrências descrevem os registros criminais. Por decisão de projeto, os indicadores de infraestrutura e os de ocorrências são exibidos de forma paralela e nunca fundidos em uma nota única de ``segurança'' ou ``periculosidade''. Essa separação atende à crítica de \citeonline{caldeira2000} à ``fala do crime'' e evita que a ferramenta reproduza os mapas de perigo que alimentam a segregação.

Os três indicadores compartilham a mesma classificação de equipamentos, agrupados em seis categorias: (i) lazer e convivência (praças, parques e áreas verdes); (ii) cultura e esporte (bibliotecas, centros culturais, CEUs, clubes-escola e equipamentos esportivos); (iii) mobilidade (pontos e terminais de ônibus, estações de metrô e trem, ciclovias); (iv) comércio de rua de uso público (feiras livres e mercados municipais); (v) serviços essenciais (unidades de saúde e escolas públicas); e (vi) conectividade e permanência (pontos de Wi-Fi público e, quando disponível, iluminação pública). Cada categoria reúne um ou mais tipos de equipamento, como unidade básica de saúde, escola pública, praça, parque, ponto de ônibus ou estação de metrô e trem.

O \textbf{Índice de Proximidade Local} (IPL) é o indicador específico do endereço consultado. Para cada tipo de equipamento $j$, o sistema calcula a distância $d_j$, em metros, entre o ponto obtido pelo CEP e o equipamento mais próximo daquele tipo, por meio de busca de vizinho mais próximo no banco geoespacial (operador de distância indexado do PostGIS sobre o tipo \emph{geography}). A distância é convertida em uma nota de 0 a 100 por uma função linear decrescente, definida por dois limiares próprios de cada tipo: a distância ideal $a_j$, até a qual a nota é máxima, e a distância máxima $b_j$, a partir da qual a nota é nula (Equação~\ref{eq:nota}):

\begin{equation}
	\label{eq:nota}
	s_j =
	\begin{cases}
		100,                                    & \text{se } d_j \leq a_j    \\[4pt]
		100 \cdot \dfrac{b_j - d_j}{b_j - a_j}, & \text{se } a_j < d_j < b_j \\[8pt]
		0,                                      & \text{se } d_j \geq b_j
	\end{cases}
\end{equation}

A nota da categoria $k$, $S_k$, é a média das notas dos tipos que a compõem, e o IPL do endereço $c$ é a média ponderada das categorias, com os mesmos pesos $w_k$ ajustáveis pelo usuário (Equação~\ref{eq:ipl}):

\begin{equation}
	\label{eq:ipl}
	\text{IPL}_c = \frac{\sum_{k=1}^{6} w_k \, S_k}{\sum_{k=1}^{6} w_k}
\end{equation}

Os limiares iniciais, apresentados no Quadro~\ref{qua:limiares}, tomam como referência a distância de 400 metros, percorrida em cerca de cinco minutos de caminhada, e são tratados como parâmetros de configuração, a serem revistos com base nos testes com os moradores. Além da nota, o sistema informa para cada tipo a distância em linha reta, o tempo estimado de caminhada e a quantidade de equipamentos a até 500 metros e a até 1 quilômetro.

\begin{table}[htb]
	\caption{Limiares iniciais de distância por tipo de equipamento}
	\label{qua:limiares}
	\centering
	\small
	\begin{tabularx}{\textwidth}{L{3.6cm}XC{2.2cm}C{2.2cm}}
		\toprule
		\textbf{Categoria}  & \textbf{Tipos de equipamento}                               & \textbf{Ideal $a_j$ (m)} & \textbf{Máxima $b_j$ (m)} \\
		\midrule
		Lazer e convivência & Praça; parque ou área verde                                 & 400                      & 1\,200                    \\
		Cultura e esporte   & Biblioteca ou centro cultural; CEU ou equipamento esportivo & 800                      & 2\,000                    \\
		Mobilidade          & Ponto de ônibus                                             & 300                      & 800                       \\
		Mobilidade          & Estação de metrô ou trem; ciclovia                          & 800                      & 2\,000                    \\
		Comércio de rua     & Feira livre ou mercado municipal                            & 800                      & 1\,500                    \\
		Serviços essenciais & Unidade básica de saúde; escola pública                     & 800                      & 2\,000                    \\
		Conectividade       & Ponto de Wi-Fi público                                      & 400                      & 1\,200                    \\
		\bottomrule
	\end{tabularx}
	\fonte{Elaborado pelo autor (2026).}
\end{table}

O \textbf{Índice de Vitalidade do Espaço Público} (IVEP) complementa o IPL na escala do distrito, indicando se a região como um todo oferece infraestrutura que atrai uso cotidiano do espaço. Para cada categoria $k$ e distrito $d$, calcula-se a densidade de equipamentos por quilômetro quadrado e por dez mil habitantes, convertida em posição percentil entre os 96 distritos ($P_{k,d}$, de 0 a 100). O índice é a média ponderada dos percentis, conforme a Equação~\ref{eq:ivep}:

\begin{equation}
	\label{eq:ivep}
	\text{IVEP}_d = \frac{\sum_{k=1}^{6} w_k \, P_{k,d}}{\sum_{k=1}^{6} w_k}
\end{equation}

\noindent em que $w_k$ é o peso da categoria $k$, igual a 1 por padrão e ajustável pelo usuário entre 0 e 5, de modo que cada pessoa possa refletir suas prioridades de rotina.

O terceiro indicador compõe os \textbf{indicadores oficiais de ocorrências}, estimados com base nos microdados georreferenciados provenientes de boletins de ocorrência divulgados pela SSP-SP \cite{ssp2025}, correlacionados aos distritos mediante operação espacial do tipo ponto-em-polígono. Em conformidade com as orientações do manual de interpretação elaborado pela própria secretaria \cite{ssp2005}, que ressalta a subnotificação e advertência quanto a comparações entre unidades de dimensões distintas, os valores são expressos como taxa anualizada por cem mil habitantes, discriminados por natureza da ocorrência (roubo, furto, roubo de celular, roubo de veículo e crimes contra a vida), calculados em janela móvel trienal, conforme a Equação~\ref{eq:taxa}:

\begin{equation}
\label{eq:taxa}
T_{n,d} = \frac{O_{n,d}}{\text{Pop}_d \times a} \times 100\,000
\end{equation}

\noindent em que $O_{n,d}$ denota o quantitativo de ocorrências da natureza $n$ localizadas no distrito $d$ ao longo do período considerado, $\text{Pop}_d$ representa a população do distrito conforme dados do Censo Demográfico 2022 \cite{ibge2022} e $a$ corresponde ao número de anos abrangidos pela janela. Cada taxa é disponibilizada acompanhada da média municipal, da tendência observada no período e de notificação permanente sobre a questão da subnotificação.

\subsection{Fontes de dados e integrações}

O GeoHabita consome exclusivamente fontes abertas ou gratuitas, organizadas em dois fluxos. As consultas de endereço ocorrem em tempo real e passam por camada de cache; as bases estatísticas e cartográficas são importadas por rotinas periódicas de extração, transformação e carga (ETL) e versionadas com data de referência. O Quadro~\ref{qua:fontes} sintetiza as integrações previstas.

\begin{longtable}{L{3.1cm}L{4.3cm}L{2.3cm}L{4.4cm}}
	\caption{Fontes de dados e integrações}
	\label{qua:fontes}                                                                                                                                                                                                                                \\
	\toprule
	\textbf{Fonte}                        & \textbf{Uso no sistema}                                                                              & \textbf{Acesso}                    & \textbf{Restrições}                                           \\
	\midrule
	\endfirsthead
	\toprule
	\textbf{Fonte}                        & \textbf{Uso no sistema}                                                                              & \textbf{Acesso}                    & \textbf{Restrições}                                           \\
	\midrule
	\endhead
	\bottomrule
	\endfoot
	ViaCEP                                & Validação do CEP e obtenção de logradouro, bairro e código IBGE                                      & API REST, tempo real               & Não retorna coordenadas; bloqueio em uso massivo              \\
	BrasilAPI (CEP v2) e AwesomeAPI       & Contingência do ViaCEP e obtenção de latitude e longitude                                            & API REST, tempo real               & Coordenadas aproximadas e eventualmente ausentes              \\
	Nominatim (OpenStreetMap)             & Geocodificação de endereço completo quando o CEP não traz coordenadas                                & API REST, tempo real               & Até 1 requisição por segundo; cache e atribuição obrigatórios \\
	GeoSampa                              & Polígonos dos 96 distritos e camadas de equipamentos públicos, áreas verdes, transporte e mobilidade & WFS, carga periódica               & Sistema de referência SIRGAS 2000/UTM 23S; exige reprojeção   \\
	Portal de dados abertos da Prefeitura & Feiras livres, pontos de Wi-Fi público e demais cadastros municipais complementares                  & Arquivos CSV/JSON, carga periódica & Coordenadas a validar; atualização irregular                  \\
	SSP-SP                                & Microdados de boletins de ocorrência com coordenadas                                                 & Planilhas anuais, carga periódica  & Colunas variam entre anos; coordenadas nulas; subnotificação  \\
	IBGE (Censo 2022)                     & População por distrito para cálculo de taxas e densidades                                            & Arquivos e API, carga periódica    & Atualização decenal                                           \\
	OpenStreetMap e Leaflet               & Base cartográfica do mapa interativo                                                                 & Biblioteca JavaScript              & Atribuição obrigatória (licença ODbL)                         \\
\end{longtable}
\vspace{-0.4cm}
\fonte{Elaborado pelo autor (2026).}

\subsection{Requisitos funcionais}

Os requisitos funcionais descrevem os comportamentos observáveis do sistema \cite{sommerville2018}. Foram identificados 25 requisitos, agrupados em cinco módulos, dos quais 14 são classificados como essenciais (\emph{Must}) e compõem o MVP (Quadro~\ref{qua:rf}).

\begin{longtable}{C{1.2cm}L{9.2cm}C{1.5cm}C{2.2cm}}
	\caption{Requisitos funcionais}
	\label{qua:rf}                                                                                                                                                                                                                                \\
	\toprule
	\textbf{ID} & \textbf{Descrição}                                                                                                                                                                          & \textbf{Prior.} & \textbf{Módulo} \\
	\midrule
	\endfirsthead
	\toprule
	\textbf{ID} & \textbf{Descrição}                                                                                                                                                                          & \textbf{Prior.} & \textbf{Módulo} \\
	\midrule
	\endhead
	\bottomrule
	\endfoot
	RF01        & O sistema deve permitir a consulta por CEP, validando o formato de oito dígitos antes de qualquer chamada externa.                                                                          & Must            & Consulta        \\
	RF02        & O sistema deve preencher automaticamente logradouro, bairro e município a partir do CEP informado.                                                                                          & Must            & Consulta        \\
	RF03        & O sistema deve permitir a consulta por endereço textual quando o usuário não souber o CEP.                                                                                                  & Should          & Consulta        \\
	RF04        & O sistema deve converter o endereço em coordenadas geográficas, utilizando serviço de contingência em caso de falha do serviço principal.                                                   & Must            & Consulta        \\
	RF05        & O sistema deve identificar o distrito administrativo correspondente às coordenadas por operação de ponto-em-polígono.                                                                       & Must            & Consulta        \\
	RF06        & O sistema deve informar ao usuário quando o endereço estiver fora do município de São Paulo.                                                                                                & Must            & Consulta        \\
	RF07        & O sistema deve listar os equipamentos públicos de uso do espaço em raio de 500 metros e de 1 quilômetro do ponto consultado, agrupados pelas seis categorias do IVEP, com nome e distância. & Must            & Diagnóstico     \\
	RF23        & O sistema deve calcular, para o endereço consultado, a distância em linha reta até o equipamento mais próximo de cada tipo, com o tempo estimado de caminhada.                              & Must            & Diagnóstico     \\
	RF24        & O sistema deve calcular e exibir o Índice de Proximidade Local do endereço e a nota de cada categoria, destacando os tipos de equipamento ausentes no raio máximo.                          & Must            & Diagnóstico     \\
	RF25        & O sistema deve calcular a distância e o tempo pela rede de calçadas e vias até o equipamento mais próximo de cada tipo, quando o usuário solicitar.                                         & Could           & Diagnóstico     \\
	RF08        & O sistema deve exibir o IVEP do distrito, o percentil de cada categoria e a comparação com a mediana do município.                                                                          & Must            & Diagnóstico     \\
	RF09        & O sistema deve permitir que o usuário ajuste os pesos das categorias do IPL e do IVEP e recalcule os índices sem recarregar a página.                                                       & Should          & Diagnóstico     \\
	RF10        & O sistema deve exibir as taxas de ocorrências por cem mil habitantes do distrito, por natureza, com média municipal e tendência dos últimos três anos.                                      & Must            & Diagnóstico     \\
	RF11        & O sistema deve apresentar, junto às taxas criminais, nota explicativa sobre subnotificação e sobre a metodologia de cálculo.                                                                & Must            & Diagnóstico     \\
	RF12        & O sistema deve exibir a fonte e a data de referência de cada indicador apresentado.                                                                                                         & Must            & Diagnóstico     \\
	RF13        & O sistema deve exibir mapa interativo com o ponto consultado, o limite do distrito e os equipamentos do entorno, filtráveis por categoria.                                                  & Must            & Visualização    \\
	RF14        & O sistema deve oferecer, para todo conteúdo cartográfico, tabela textual equivalente navegável por teclado e leitor de tela.                                                                & Must            & Visualização    \\
	RF15        & O sistema deve anunciar por região viva (\emph{live region}) WAI-ARIA a conclusão de cada consulta assíncrona e eventuais erros.                                                            & Should          & Visualização    \\
	RF16        & O sistema deve permitir a comparação lado a lado de até três endereços ou distritos.                                                                                                        & Should          & Comparação      \\
	RF17        & O sistema deve permitir compartilhar um diagnóstico por link, sem armazenar dados pessoais do consulente.                                                                                   & Could           & Comparação      \\
	RF18        & O sistema deve permitir exportar o diagnóstico em formato PDF acessível ou CSV.                                                                                                             & Could           & Comparação      \\
	RF19        & O sistema deve permitir ao administrador executar as rotinas de carga das bases externas e registrar o resultado de cada execução.                                                          & Should          & Administração   \\
	RF20        & O sistema deve validar as bases importadas, descartando registros sem coordenadas ou fora dos limites do município e reportando a quantidade descartada.                                    & Should          & Administração   \\
	RF21        & O sistema deve manter o histórico de versões dos indicadores com data de referência de cada carga.                                                                                          & Could           & Administração   \\
	RF22        & O sistema deve disponibilizar uma API pública de leitura dos indicadores agregados por distrito, documentada em OpenAPI.                                                                    & Could           & Administração   \\
\end{longtable}
\vspace{-0.4cm}
\fonte{Elaborado pelo autor (2026).}

\subsection{Requisitos não funcionais}

Os requisitos não funcionais estabelecem restrições sobre os serviços oferecidos e sobre o processo de desenvolvimento \cite{sommerville2018}. Para garantir sua verificabilidade, cada requisito de produto está associado a uma métrica objetiva, conforme recomendam \citeonline{pressman2021}. Os requisitos de produto estão classificados segundo as características de qualidade da ISO/IEC 25010 (Quadro~\ref{qua:rnf}).

\begin{longtable}{C{1.3cm}L{2.6cm}L{6.1cm}L{4.1cm}}
	\caption{Requisitos não funcionais de produto}
	\label{qua:rnf}                                                                                                                                                                                                                                                                                                                                             \\
	\toprule
	\textbf{ID} & \textbf{Característica}              & \textbf{Requisito}                                                                                                                                                                                    & \textbf{Métrica de verificação}                                                                \\
	\midrule
	\endfirsthead
	\toprule
	\textbf{ID} & \textbf{Característica}              & \textbf{Requisito}                                                                                                                                                                                    & \textbf{Métrica de verificação}                                                                \\
	\midrule
	\endhead
	\bottomrule
	\endfoot
	RNF01       & Eficiência de desempenho             & O diagnóstico de um CEP já armazenado em cache deve ser exibido em até 2 segundos.                                                                                                                    & Percentil 95 $\leq$ 2 s com 30 usuários simultâneos                                            \\
	RNF02       & Eficiência de desempenho             & O diagnóstico de um CEP inédito deve ser exibido em até 6 segundos.                                                                                                                                   & Percentil 95 $\leq$ 6 s em teste de carga                                                      \\
	RNF03       & Usabilidade (acessibilidade)         & A interface deve atender às diretrizes WCAG 2.1 no nível AA e empregar marcações WAI-ARIA em todos os componentes dinâmicos.                                                                          & Zero erros críticos em Lighthouse e axe; contraste $\geq$ 4,5:1                                \\
	RNF04       & Usabilidade (operabilidade)          & Todas as funções devem ser operáveis exclusivamente por teclado, com foco visível e ordem lógica.                                                                                                     & Roteiro de teste manual com NVDA e Orca                                                        \\
	RNF05       & Usabilidade (inteligibilidade)       & Os textos explicativos devem usar linguagem simples, sem jargão técnico ou geográfico sem explicação.                                                                                                 & Avaliação em teste com usuários; SUS $\geq$ 70                                                 \\
	RNF06       & Confiabilidade (tolerância a falhas) & A falha de uma fonte externa não deve impedir a exibição dos demais blocos do diagnóstico.                                                                                                            & Teste com serviço externo simulado fora do ar                                                  \\
	RNF07       & Confiabilidade (disponibilidade)     & A aplicação deve manter disponibilidade mensal mínima de 98\%.                                                                                                                                        & Monitoramento externo de disponibilidade                                                       \\
	RNF08       & Segurança                            & Todo o tráfego deve ocorrer por HTTPS, e as rotas públicas devem limitar-se a 60 requisições por minuto por endereço IP.                                                                              & Teste automatizado; resposta HTTP 429 acima do limite                                          \\
	RNF09       & Segurança                            & A aplicação não deve apresentar vulnerabilidades de alta severidade do OWASP Top 10.                                                                                                                  & Varredura com OWASP ZAP e auditoria de dependências                                            \\
	RNF10       & Manutenibilidade                     & O código deve adotar arquitetura em camadas (apresentação, aplicação, domínio e infraestrutura), com adaptadores isolando cada serviço externo.                                                       & Revisão em \emph{pull request}; ausência de chamadas externas fora da camada de infraestrutura \\
	RNF11       & Manutenibilidade (testabilidade)     & As regras de negócio devem ter cobertura mínima de 70\% por testes unitários.                                                                                                                         & Relatório de cobertura na integração contínua                                                  \\
	RNF12       & Compatibilidade e portabilidade      & A interface deve funcionar nas duas versões mais recentes dos navegadores Chrome, Firefox, Edge e Safari, em telas a partir de 360 pixels de largura, e a aplicação deve ser empacotada em contêiner. & Roteiro multiplataforma; implantação em nuvem com um único comando                             \\
\end{longtable}
\vspace{-0.4cm}
\fonte{Elaborado pelo autor (2026).}

Os requisitos organizacionais e externos, que não se referem diretamente ao comportamento do produto, estão reunidos no Quadro~\ref{qua:rnf-org}.

\begin{table}[htb]
	\caption{Requisitos não funcionais organizacionais e externos}
	\label{qua:rnf-org}
	\centering
	\small
	\begin{tabularx}{\textwidth}{C{1.3cm}C{2.4cm}X}
		\toprule
		\textbf{ID} & \textbf{Tipo}      & \textbf{Requisito}                                                                                                                                                                                                         \\
		\midrule
		RNF13       & Organizacional     & O desenvolvimento deve seguir práticas do Scrum, com \emph{sprints} de duas semanas, controle de versão em Git e integração contínua executando os testes a cada envio ao repositório.                                     \\
		RNF14       & Organizacional     & A pilha tecnológica deve compreender Python com micro-framework web no serviço de retaguarda, JavaScript no cliente, PostgreSQL com extensão PostGIS na persistência e hospedagem em provedor de nuvem com plano gratuito. \\
		RNF15       & Externo (legal)    & O tratamento de dados deve observar a Lei Geral de Proteção de Dados \cite{lgpd2018}: a consulta não exige cadastro e nenhum CEP consultado é associado a pessoa identificável.                                            \\
		RNF16       & Externo (legal)    & A plataforma deve atender ao art. 63 da Lei Brasileira de Inclusão \cite{lbi2015}, que torna obrigatória a acessibilidade em sítios da internet.                                                                           \\
		RNF17       & Externo (licenças) & Todas as fontes devem ser atribuídas conforme suas licenças, e o código deve ser distribuído sob a GNU AGPL v3 \cite{agplv3}.                                                                                              \\
		RNF18       & Externo (ético)    & Os indicadores não devem empregar rótulos estigmatizantes, rankings de ``bairros perigosos'' nem escalas cromáticas de alarme para dados criminais.                                                                        \\
		\bottomrule
	\end{tabularx}
	\fonte{Elaborado pelo autor (2026).}
\end{table}

\subsection{Regras de negócio e requisitos de domínio}

As regras de negócio condensam o conhecimento do domínio urbano e estatístico que o sistema deve respeitar independentemente da interface. Elas são implementadas na camada de domínio e testadas de forma isolada (Quadro~\ref{qua:rn}).

\begin{longtable}{C{1.2cm}L{10.6cm}L{2.6cm}}
	\caption{Regras de negócio}
	\label{qua:rn}                                                                                                                                                                                                                                                                                               \\
	\toprule
	\textbf{ID} & \textbf{Regra}                                                                                                                                                                                                                                                        & \textbf{Relacionada a} \\
	\midrule
	\endfirsthead
	\toprule
	\textbf{ID} & \textbf{Regra}                                                                                                                                                                                                                                                        & \textbf{Relacionada a} \\
	\midrule
	\endhead
	\bottomrule
	\endfoot
	RN01        & Um CEP válido possui exatamente oito dígitos numéricos; hífens, pontos e espaços são removidos antes da validação.                                                                                                                                                    & RF01                   \\
	RN02        & Um CEP inexistente gera a mensagem ``CEP não encontrado'' e libera a consulta por endereço textual.                                                                                                                                                                   & RF02, RF03             \\
	RN03        & O serviço de contingência é acionado quando o serviço principal retorna erro ou excede 3 segundos de resposta.                                                                                                                                                        & RF04, RNF06            \\
	RN04        & Resultados de CEP e geocodificação permanecem em cache por 30 dias; indicadores agregados, até a próxima carga.                                                                                                                                                       & RF04, RNF01            \\
	RN05        & A coordenada obtida pelo CEP é tratada como aproximada; a interface informa que ela identifica o distrito, mas não a posição exata do imóvel.                                                                                                                         & RF05, RF07             \\
	RN06        & O IVEP é calculado por percentis entre os 96 distritos e varia de 0 a 100; os pesos padrão são iguais a 1 e aceitam valores inteiros de 0 a 5.                                                                                                                        & RF08, RF09             \\
	RN07        & Os indicadores de vitalidade e os indicadores de ocorrências nunca são combinados em uma pontuação única.                                                                                                                                                             & RF08, RF10, RNF18      \\
	RN08        & As taxas de ocorrências usam janela de três anos completos e são expressas por cem mil habitantes; o ano corrente incompleto não entra no cálculo.                                                                                                                    & RF10                   \\
	RN09        & Naturezas com menos de 30 ocorrências no distrito durante a janela são exibidas como ``número insuficiente para comparação'', sem taxa nem tendência.                                                                                                                 & RF10, RF11             \\
	RN10        & Registros da SSP-SP com coordenadas nulas, iguais a zero ou fora do município são descartados na carga e contabilizados no relatório de validação.                                                                                                                    & RF20                   \\
	RN11        & Indicadores com data de referência superior a 24 meses recebem o aviso ``dado desatualizado''.                                                                                                                                                                        & RF12                   \\
	RN12        & A distância até o equipamento mais próximo é a distância geodésica em linha reta, em metros, entre o ponto do CEP e o ponto do equipamento; para áreas como parques, usa-se o ponto mais próximo do contorno. A interface identifica a medida como ``em linha reta''. & RF23                   \\
	RN13        & A nota de cada tipo segue a Equação~\ref{eq:nota}, com os limiares do Quadro~\ref{qua:limiares}; os limiares são parâmetros de configuração alteráveis pelo administrador sem mudança de código.                                                                      & RF24                   \\
	RN14        & O tempo estimado de caminhada é calculado à velocidade de 80 metros por minuto (4,8 km/h) e arredondado para o minuto inteiro superior.                                                                                                                               & RF23                   \\
	RN15        & Tipo sem nenhum equipamento a menos da distância máxima recebe nota zero, e a interface informa ``nenhum equipamento deste tipo a menos de $b_j$ metros''.                                                                                                            & RF24                   \\
	RN16        & O IPL é calculado sempre sobre a coordenada do endereço consultado, e nunca sobre o centroide do distrito ou do bairro.                                                                                                                                               & RF24                   \\
\end{longtable}
\vspace{-0.4cm}
\fonte{Elaborado pelo autor (2026).}

A regra RN09 transpõe para o sistema a advertência do manual de interpretação da SSP-SP sobre comparações apoiadas em poucos casos, e a regra RN07 materializa a opção, discutida na fundamentação, de não reduzir a qualidade de um território a um único número.

\subsection{Casos de uso}

Os requisitos funcionais foram agrupados em oito casos de uso principais, que orientam o diagrama de casos de uso e o desenho das telas:

\begin{alineas}
	\item \textbf{UC01 -- Consultar endereço}: o munícipe informa CEP ou endereço, e o sistema valida, geocodifica e identifica o distrito (RF01 a RF06);
	\item \textbf{UC02 -- Visualizar diagnóstico socioespacial}: o sistema apresenta equipamentos do entorno, IVEP e indicadores de ocorrências, com fontes e datas (RF07, RF08, RF10 a RF12);
	\item \textbf{UC02a -- Avaliar proximidade do endereço}: o sistema calcula a distância até o equipamento mais próximo de cada tipo e o IPL do endereço (RF23 a RF25);
	\item \textbf{UC03 -- Personalizar prioridades}: o munícipe ajusta os pesos das categorias e acompanha o recálculo do IVEP (RF09);
	\item \textbf{UC04 -- Explorar mapa ou tabela equivalente}: o munícipe navega pelo mapa ou pela tabela textual e filtra categorias (RF13 a RF15);
	\item \textbf{UC05 -- Comparar localidades}: o munícipe compara até três endereços ou distritos (RF16);
	\item \textbf{UC06 -- Compartilhar ou exportar diagnóstico}: o munícipe gera link ou arquivo do diagnóstico (RF17, RF18);
	\item \textbf{UC07 -- Atualizar bases de dados}: o administrador executa as cargas, valida os registros e publica nova versão dos indicadores (RF19 a RF21).
\end{alineas}

\subsection{Histórias de usuário e critérios de aceitação}

Para o planejamento das \emph{sprints}, os casos de uso foram desdobrados em histórias de usuário no formato ``Como \emph{papel}, quero \emph{ação}, para \emph{benefício}'' \cite{cohn2004}, estimadas em pontos de complexidade relativa (Quadro~\ref{qua:hu}).

\begin{longtable}{C{1.2cm}L{9.0cm}C{2.1cm}C{1.3cm}}
	\caption{Histórias de usuário}
	\label{qua:hu}                                                                                                                                                                                                                             \\
	\toprule
	\textbf{ID} & \textbf{História}                                                                                                                                                                    & \textbf{Requisitos} & \textbf{Pontos} \\
	\midrule
	\endfirsthead
	\toprule
	\textbf{ID} & \textbf{História}                                                                                                                                                                    & \textbf{Requisitos} & \textbf{Pontos} \\
	\midrule
	\endhead
	\bottomrule
	\endfoot
	HU01        & Como munícipe, quero digitar um CEP e ver o endereço reconhecido, para iniciar a consulta sem preencher formulários longos.                                                          & RF01, RF02          & 3               \\
	HU02        & Como munícipe, quero saber em que distrito fica o endereço, para entender a região a que ele pertence.                                                                               & RF04--RF06          & 5               \\
	HU03        & Como munícipe, quero ver praças, parques, feiras, transporte e serviços públicos a até 1 km, para saber se a vizinhança tem vida nas ruas.                                           & RF07, RF13          & 8               \\
	HU04        & Como munícipe, quero ver o índice de vitalidade do distrito em comparação com a cidade, para avaliar a região de forma objetiva.                                                     & RF08                & 5               \\
	HU05        & Como munícipe, quero ver as taxas oficiais de ocorrências com explicação sobre seus limites, para decidir com informação confiável e sem alarmismo.                                  & RF10--RF12          & 8               \\
	HU06        & Como usuário de leitor de tela, quero acessar todo o conteúdo do mapa em forma de tabela e ser avisado quando a consulta terminar, para usar a plataforma com autonomia.             & RF14, RF15          & 5               \\
	HU07        & Como munícipe, quero dar mais peso ao transporte do que ao lazer, para que o índice reflita minha rotina.                                                                            & RF09                & 5               \\
	HU08        & Como munícipe, quero comparar três endereços lado a lado, para escolher entre as opções finais.                                                                                      & RF16                & 8               \\
	HU09        & Como administrador, quero carregar e validar novas bases da SSP-SP e do GeoSampa, para manter os indicadores atualizados.                                                            & RF19, RF20          & 8               \\
	HU10        & Como munícipe, quero saber a que distância do endereço ficam a UBS, a escola, a praça, o ponto de ônibus e a estação mais próximos, para avaliar se consigo fazer minha rotina a pé. & RF23, RF24          & 8               \\
\end{longtable}
\vspace{-0.4cm}
\fonte{Elaborado pelo autor (2026).}

Os critérios de aceitação são redigidos na sintaxe Gherkin (Dado, Quando, Então), o que permite convertê-los diretamente em testes automatizados de aceitação. A título de exemplo, a história HU01 é aceita quando todos os cenários a seguir forem atendidos:

\begin{alineas}
	\item \textbf{CEP válido}: \emph{dado} que o munícipe está na página inicial, \emph{quando} informa ``01310-100'', \emph{então} o sistema exibe ``Avenida Paulista, Bela Vista, São Paulo -- SP'';
	\item \textbf{Formato inválido}: \emph{quando} informa ``1234'', \emph{então} o sistema exibe ``O CEP deve ter 8 dígitos'' e nenhuma API externa é chamada;
	\item \textbf{CEP inexistente}: \emph{quando} informa ``99999-999'', \emph{então} o sistema exibe ``CEP não encontrado'' e oferece a consulta por endereço;
	\item \textbf{Serviço principal indisponível}: \emph{dado} que o ViaCEP não responde em 3 segundos, \emph{quando} o munícipe informa um CEP válido, \emph{então} o endereço é obtido pelo serviço de contingência.
\end{alineas}

Considera-se uma história pronta para desenvolvimento (\emph{Definition of Ready}) quando possui critérios de aceitação escritos, estimativa acordada e dependências externas identificadas. Considera-se concluída (\emph{Definition of Done}) quando o código foi revisado, os testes unitários e de aceitação passam na integração contínua, a verificação de acessibilidade não aponta erros críticos e o incremento está implantado no ambiente em nuvem.

\subsection{Priorização e planejamento incremental}

As dez histórias totalizam 63 pontos, distribuídos em quatro \emph{sprints} de duas semanas, precedidas de uma \emph{sprint} de preparação do ambiente (Quadro~\ref{qua:sprints}). O MVP encerra-se ao fim da terceira \emph{sprint}, quando o fluxo completo de consulta e diagnóstico estará disponível para validação com os moradores.

\begin{table}[htb]
	\caption{Planejamento das \emph{sprints}}
	\label{qua:sprints}
	\centering
	\small
	\begin{tabularx}{\textwidth}{C{1.4cm}XC{3.2cm}C{1.4cm}}
		\toprule
		\textbf{Sprint} & \textbf{Meta}                                                                                  & \textbf{Histórias} & \textbf{Pontos} \\
		\midrule
		0               & Repositório, integração contínua, banco com PostGIS e primeira carga dos distritos do GeoSampa & --                 & --              \\
		1               & Consultar um CEP e localizar o distrito                                                        & HU01, HU02, HU09   & 16              \\
		2               & Exibir a infraestrutura do entorno e a distância ao equipamento mais próximo de cada tipo      & HU03, HU10         & 16              \\
		3               & Calcular IVEP e indicadores de ocorrências, com mapa e tabela acessível (MVP)                  & HU04, HU05, HU06   & 18              \\
		4               & Personalização de pesos, comparação e ajustes após testes com usuários                         & HU07, HU08         & 13              \\
		\bottomrule
	\end{tabularx}
	\fonte{Elaborado pelo autor (2026).}
\end{table}

\subsection{Plano de testes}

A estratégia de verificação segue a pirâmide de testes: grande volume de testes unitários sobre as regras de negócio, testes de integração sobre os adaptadores de dados e o banco geoespacial, e um conjunto reduzido de testes ponta a ponta sobre os fluxos críticos \cite{pressman2021}. Nos testes automatizados, os serviços externos são substituídos por simulações (\emph{mocks}), o que torna os resultados reprodutíveis; testes de contrato agendados verificam, separadamente, se o formato das respostas reais das APIs permanece inalterado. O Quadro~\ref{qua:niveis} apresenta os níveis de teste previstos.

\begin{longtable}{L{2.6cm}L{5.0cm}L{3.5cm}L{3.0cm}}
	\caption{Níveis e tipos de teste}
	\label{qua:niveis}                                                                                                                                                                       \\
	\toprule
	\textbf{Nível} & \textbf{Objeto de verificação}                                  & \textbf{Ferramentas}                           & \textbf{Critério de saída}                           \\
	\midrule
	\endfirsthead
	\toprule
	\textbf{Nível} & \textbf{Objeto de verificação}                                  & \textbf{Ferramentas}                           & \textbf{Critério de saída}                           \\
	\midrule
	\endhead
	\bottomrule
	\endfoot
	Unitário       & Regras RN01 a RN16, cálculo do IPL, do IVEP e das taxas         & pytest                                         & 100\% aprovados; cobertura $\geq$ 70\%               \\
	Integração     & Adaptadores de CEP e geocodificação, consultas espaciais, cache & pytest com banco PostGIS em contêiner          & 100\% aprovados                                      \\
	Contrato       & Formato das respostas de ViaCEP, BrasilAPI e Nominatim          & Execução diária agendada                       & Alerta em caso de divergência                        \\
	Aceitação      & Cenários Gherkin das histórias de usuário                       & pytest-bdd ou Behave                           & Todas as histórias \emph{Must} aprovadas             \\
	Ponta a ponta  & Fluxo consulta, diagnóstico, mapa e tabela                      & Playwright                                     & Aprovado em Chrome e Firefox                         \\
	Desempenho     & RNF01 e RNF02                                                   & k6                                             & Metas de percentil 95 atendidas                      \\
	Segurança      & RNF08 e RNF09                                                   & OWASP ZAP; pip-audit                           & Nenhum alerta alto                                   \\
	Acessibilidade & RNF03 e RNF04                                                   & Lighthouse, axe, NVDA, Orca                    & Zero erros críticos                                  \\
	Usabilidade    & RNF05 e aderência às necessidades levantadas                    & Teste moderado com moradores; questionário SUS & Tarefas concluídas por $\geq$ 80\% dos participantes \\
\end{longtable}
\vspace{-0.4cm}
\fonte{Elaborado pelo autor (2026).}

Os casos de teste derivam diretamente dos requisitos e das regras de negócio. O Quadro~\ref{qua:ct} relaciona os principais.

\begin{longtable}{C{1.1cm}C{1.8cm}L{4.6cm}L{4.4cm}C{1.8cm}}
	\caption{Casos de teste principais}
	\label{qua:ct}                                                                                                                                                                                \\
	\toprule
	\textbf{ID} & \textbf{Requisito} & \textbf{Entrada ou condição}                                 & \textbf{Resultado esperado}                                                & \textbf{Tipo}  \\
	\midrule
	\endfirsthead
	\toprule
	\textbf{ID} & \textbf{Requisito} & \textbf{Entrada ou condição}                                 & \textbf{Resultado esperado}                                                & \textbf{Tipo}  \\
	\midrule
	\endhead
	\bottomrule
	\endfoot
	CT01        & RF01, RN01         & ``1234'', ``abcdefgh'' e ``123456789''                       & Mensagem de formato inválido e nenhuma chamada externa                     & Unitário       \\
	CT02        & RF02               & CEP 01310-100                                                & Logradouro, bairro e município preenchidos corretamente                    & Integração     \\
	CT03        & RN02               & CEP 99999-999                                                & ``CEP não encontrado'' e consulta textual liberada                         & Integração     \\
	CT04        & RN03, RNF06        & Serviço principal simulado com atraso de 5 s                 & Resposta obtida pelo serviço de contingência                               & Integração     \\
	CT05        & RF05               & Coordenadas conhecidas no interior de três distritos         & Distrito correto em todos os casos                                         & Integração     \\
	CT06        & RF06               & CEP de município vizinho                                     & Aviso de endereço fora de São Paulo                                        & Integração     \\
	CT07        & RF07               & Ponto com equipamentos cadastrados a 300 m, 800 m e 1,5 km   & Os dois primeiros listados com distância; o terceiro ausente               & Integração     \\
	CT08        & RF08, RN06         & Base fictícia com percentis conhecidos e pesos iguais        & IVEP igual à média aritmética dos percentis                                & Unitário       \\
	CT09        & RF09               & Peso de mobilidade igual a 5 e demais iguais a 1             & IVEP recalculado conforme a Equação~\ref{eq:ivep}, sem recarregar a página & Ponta a ponta  \\
	CT10        & RN08               & Distrito com 300 ocorrências em 3 anos e 100\,000 habitantes & Taxa de 100 por cem mil habitantes ao ano                                  & Unitário       \\
	CT11        & RN09               & Natureza com 12 ocorrências na janela                        & ``Número insuficiente para comparação'', sem taxa                          & Unitário       \\
	CT12        & RN07               & Resposta da API interna do diagnóstico                       & Ausência de qualquer campo que combine os dois conjuntos de indicadores    & Unitário       \\
	CT13        & RN10, RF20         & Planilha com 5 registros de coordenadas nulas ou zeradas     & Registros descartados e contabilizados no relatório de carga               & Integração     \\
	CT14        & RF14, RNF04        & Navegação apenas por teclado com leitor de tela              & Todos os equipamentos do mapa acessíveis pela tabela equivalente           & Acessibilidade \\
	CT15        & RF15               & Consulta concluída e consulta com erro                       & Ambas anunciadas pela região viva                                          & Acessibilidade \\
	CT16        & RNF08              & 61 requisições em um minuto a partir do mesmo IP             & A 61ª requisição recebe HTTP 429                                           & Segurança      \\
	CT17        & RF23, RN12         & Ponto com duas UBS cadastradas a 250 m e a 600 m             & UBS mais próxima informada a 250 m, ``em linha reta''                      & Integração     \\
	CT18        & RN13               & Ponto de ônibus a 550 m ($a_j = 300$, $b_j = 800$)           & Nota do tipo igual a 50                                                    & Unitário       \\
	CT19        & RN14               & Distância de 650 m                                           & Tempo estimado de 9 minutos                                                & Unitário       \\
	CT20        & RN15               & Nenhuma estação a menos de 2\,000 m                          & Nota zero e mensagem de ausência                                           & Unitário       \\
	CT21        & RF24               & Notas de categoria 100, 50, 0, 80, 60 e 40, pesos iguais a 1 & IPL igual a 55                                                             & Unitário       \\
	CT22        & RN16               & Dois CEPs distintos no mesmo distrito                        & IPL diferentes e mesmo IVEP                                                & Integração     \\
\end{longtable}
\vspace{-0.4cm}
\fonte{Elaborado pelo autor (2026).}

Os defeitos encontrados são registrados como \emph{issues} no repositório, classificados por severidade; defeitos críticos e altos impedem o encerramento da \emph{sprint}.

\subsection{Matriz de rastreabilidade}

A matriz de rastreabilidade (Quadro~\ref{qua:rastreab}) liga as necessidades identificadas junto à comunidade aos requisitos que as atendem e aos testes que comprovam seu cumprimento, permitindo avaliar o impacto de qualquer alteração de escopo.

\begin{longtable}{L{4.8cm}L{4.2cm}C{2.0cm}C{2.6cm}}
	\caption{Matriz de rastreabilidade}
	\label{qua:rastreab}                                                                                                          \\
	\toprule
	\textbf{Necessidade}                                    & \textbf{Requisitos e regras}  & \textbf{História} & \textbf{Testes} \\
	\midrule
	\endfirsthead
	\toprule
	\textbf{Necessidade}                                    & \textbf{Requisitos e regras}  & \textbf{História} & \textbf{Testes} \\
	\midrule
	\endhead
	\bottomrule
	\endfoot
	Consultar uma região de forma simples                   & RF01--RF06, RN01--RN05        & HU01, HU02        & CT01--CT06      \\
	Saber a que distância ficam os serviços do dia a dia    & RF23--RF25, RN12--RN16        & HU10              & CT17--CT22      \\
	Conhecer a infraestrutura pública do entorno            & RF07, RF13                    & HU03              & CT07            \\
	Avaliar a vitalidade do espaço público                  & RF08, RF09, RN06              & HU04, HU07        & CT08, CT09      \\
	Dispor de dados criminais auditáveis e contextualizados & RF10--RF12, RN07--RN09, RNF18 & HU05              & CT10--CT12      \\
	Usar a plataforma com tecnologia assistiva              & RF14, RF15, RNF03, RNF04      & HU06              & CT14, CT15      \\
	Manter os dados atualizados e íntegros                  & RF19, RF20, RN10              & HU09              & CT13            \\
	Operar com segurança e desempenho                       & RNF01, RNF02, RNF08           & --                & CT16            \\
\end{longtable}
\vspace{-0.4cm}
\fonte{Elaborado pelo autor (2026).}


	\chapter{Coleta de Dados com a Comunidade}
	\section{Instrumento de Coleta de Dados e Compilação das Entrevistas de Imersão}
\label{apendice:questionario}

Apresenta-se a seguir a estrutura do roteiro semiestruturado aplicado presencialmente na fase de imersão do \emph{Design Thinking}, seguido da compilação e tabulação das respostas obtidas junto a dez cidadãos abordados em nós intermodais e praças da cidade de São Paulo.

% -------------------------------------------------------------
% PARTE I: ESTRUTURA DO FORMULÁRIO / QUESTIONÁRIO
% -------------------------------------------------------------
\subsection{Roteiro do Questionário de Imersão}
\begin{enumerate}[leftmargin=*]


    \item \textbf{[Pergunta-Filtro Eliminatória] Você esteve ou está ativamente buscando imóvel para alugar ou comprar na cidade de São Paulo nos últimos 12 meses?}
    \begin{itemize}[leftmargin=1.5em]
        \item[] (\quad) \textbf{Sim} $\rightarrow$ \emph{Prosseguir para a Questão 2.}
        \item[] (\quad) \textbf{Não} $\rightarrow$ \emph{Entrevista encerrada com agradecimento.}
    \end{itemize}

    \item \textbf{Perfil Geral e Escopo Habitacional:}
    \begin{itemize}[leftmargin=1.5em]
        \item Faixa etária e ocupação principal.
        \item Tipo de transição habitacional pretendida (aluguel individual, locação familiar, república estudantil ou compra de imóvel próprio).
        \item Regiões ou bairros de interesse preferencial na capital.
    \end{itemize}

    \item \textbf{Critérios Decisórios e Segurança Pública:}
    \begin{itemize}[leftmargin=1.5em]
        \item Além do valor financeiro do aluguel ou condomínio, quais fatores urbanos são mais decisivos na escolha do endereço?
        \item O que mais gera sensação de segurança no seu trajeto diário a pé: a presença de muros altos e portarias privadas ou ruas com fluxo constante de pedestres, comércio ativo e boa iluminação pública?
        \item Você já desistiu de visitar ou alugar um imóvel devido à ``fama'' de perigoso de um determinado bairro? De onde veio essa percepção?
    \end{itemize}

    \item \textbf{Transparência Pública e Acesso à Informação:}
    \begin{itemize}[leftmargin=1.5em]
        \item Você sabia que a Prefeitura Municipal de São Paulo mantém o portal GeoSampa, disponibilizando mapas abertos de equipamentos públicos (saúde, parques, zoneamento)?
        \item Você sabia que a Secretaria de Segurança Pública de São Paulo (SSP-SP) publica tabelas mensais com todas as ocorrências policiais registradas por Distrito Policial? Já tentou consultar essas bases ao escolher um imóvel?
    \end{itemize}

    \item \textbf{Avaliação da Proposta de Tecnologia Cívica (GeoHabita):}
    \begin{itemize}[leftmargin=1.5em]
        \item Em uma escala de 1 (nada útil) a 5 (essencial), qual o grau de utilidade de uma ferramenta pública e gratuita que informe o panorama de segurança e a distância a pé até UBSs, metrô e praças a partir da simples digitação do CEP?
        \item Você considera relevante que a ferramenta apresente dados em tabelas textuais com recursos de acessibilidade (modo de alto contraste, leitor de tela para pessoas com baixa visão)?
    \end{itemize}

    \item \textbf{Consentimento e Identificação do Participante:}
    \begin{itemize}[leftmargin=1.5em]
        \item[] (\quad) Desejo me identificar formalmente no relatório da pesquisa (Nome/E-mail).
        \item[] (\quad) Prefiro participar de forma \textbf{estritamente anônima}, ciente de que nenhum dado pessoal será exposto e que as respostas serão associadas exclusivamente a um código alfanumérico identificador (ex: Munícipe 1). Autorizo a guarda confidencial do meu contato exclusivamente para convocação das sessões de teste e feedback do protótipo.
    \end{itemize}
\end{enumerate}

\vspace{1em}

% -------------------------------------------------------------
% PARTE II: COMPILAÇÃO DAS 10 RESPOSTAS COLETADAS
% -------------------------------------------------------------
\subsection{Compilação Qualitativa das Respostas Coletadas}

Todas as pessoas selecionadas que responderam afirmativamente à pergunta-filtro (Questão 1) e optaram expressamente pela preservação do anonimato (Questão 6), sendo identificadas ordinalmente de Munícipe 1 a Munícipe 10.

\begin{table}[!htbp]
\centering
\small
\renewcommand{\arraystretch}{1.25}
\begin{tabular}{|c|p{3.2cm}|p{4.2cm}|c|c|c|}
\hline
\textbf{ID} & \textbf{Perfil do Entrevistado} & \textbf{Critérios Principais} & \textbf{GeoSampa?} & \textbf{SSP-SP?} & \textbf{Nota} \\
\hline
M1 & Estudante (21 anos), busca quarto próx. metrô & Transporte noturno, iluminação e comércio no caminho & Não & Não & 5 \\
\hline
M2 & Autônoma (35 anos), aluguel com 1 filho & Proximidade de creche/UBS, calçadas movimentadas & Não & Não & 5 \\
\hline
M3 & Dev. de software (28 anos), compra 1º imóvel & Acesso a eixos de mobilidade, praças para pets & Não & Não & 4 \\
\hline
M4 & Aposentada (64 anos), aluguel & Acessibilidade física, UBS próxima, iluminação & Não & Não & 3 \\
\hline
M5 & Atendente (24 anos), aluguel compartilhado & Pontos de ônibus seguros, fluxo de pedestres & Não & Não & 5 \\
\hline
M6 & Professor estadual (39 anos), aluguel & Corredor de ônibus, preço realista, escolas & Não & Não & 4 \\
\hline
M7 & Casal com bebê (31/29 anos), aluguel 2 dorm. & Parques públicos, calçadas largas, tranquilidade & Não & Não & 5 \\
\hline
M8 & Pós-graduando (26 anos, vindo do interior) & Segurança nos trajetos a pé, proximidade à USP & Não & Não & 5 \\
\hline
M9 & Comerciário (44 anos), compra imóvel popular & Linhas de trem/metrô, custo benefício do bairro & Não & Não & 4 \\
\hline
M10 & Designer (30 anos), usuária de bike/transporte & Ciclovias, vitalidade urbana, vida cultural & Não & Não & 5 \\
\hline
\end{tabular}
\caption{Quadro sintético dos entrevistados na fase de imersão empírica.}
\label{tab:resumo_entrevistados}
\end{table}

\vspace{0.5em}

\noindent\textbf{Síntese dos Relatos Coletados por Participante:}

\begin{description}[leftmargin=1.2em, itemsep=0.8em]
    \item[Munícipe 1 (Abordado no Terminal Pinheiros):] Jovem universitário buscando locação no Butantã ou Pinheiros. Relatou que sua maior aflição é o retorno das aulas após as 22h. Declarou preferir ruas com padarias e transeuntes a condomínios fechados em vias ermas: \emph{``O prédio pode ser blindado, mas o trajeto da estação até o portão é onde o assalto acontece''}. Desconhecia GeoSampa e os relatórios da SSP. Atribuiu nota 5 ao projeto.

    \item[Munícipe 2 (Abordada no Terminal Barra Funda):] Mãe solo buscando aluguel na Zona Norte/Oeste. Prioriza a proximidade com serviços de saúde e educação infantil. Declarou ter receio de bairros periféricos por causa do noticiário da televisão, mas admite não ter dados reais: \emph{``A imobiliária só fala maravilhas, a TV só fala de tragédia; a gente fica no escuro''}. Nunca ouviu falar de dados abertos do governo. Atribuiu nota 5.

    \item[Munícipe 3 (Abordado no Largo da Batata):] Profissional de tecnologia avaliando compra de apartamento na Zona Leste/Centro Expandido. Reclamou da opacidade dos anúncios de corretoras, que vendem segurança com base em câmeras e portarias remotas. Não conhecia os portais públicos municipais. Destacou a importância de haver métricas de áreas verdes no entorno. Atribuiu nota 4.

    \item[Munícipe 4 (Abordada no entorno da Praça da República):] Idosa que procura locação próxima a equipamentos públicos de saúde na região central. Declarou que não sabia da existência do GeoSampa e aprovou a inclusão de recursos de acessibilidade e leitor de tela em tabelas textuais. Atribuiu nota 3.

    \item[Munícipe 5 (Abordada no Terminal Barra Funda):] Trabalhadora do setor de serviços buscando moradia no eixo da Linha 3-Vermelha. Apontou que a volta para casa em pontos de ônibus escuros é seu principal fator de estresse diário. Atribuiu nota 5.

    \item[Munícipe 6 (Abordado no Terminal Pinheiros):] Professor da rede básica procurando imóvel na Zona Sul. Afirmou que o custo do aluguel o empurra para regiões que ele não conhece bem e que sofre com a falta de referências confiáveis sobre os bairros: \emph{``A gente fica refém de comentários vagos em grupos de redes sociais''}. Nunca consultou a SSP-SP por achar que os dados eram de uso restrito da polícia. Atribuiu nota 4.

    \item[Munícipe 7 (Abordados na Praça Roosevelt):] Casal à procura de imóvel de dois dormitórios para alugar na Bela Vista ou Cambuci. Enfatizaram a busca por calçadas caminháveis e praças arborizadas para passear com o carrinho do filho. Desconheciam qualquer canal de dados abertos da PMSP. Atribuíram nota 5.

    \item[Munícipe 8 (Abordado nas imediações do Butantã):] Estudante recém-chegado de outro estado buscando aluguel compartilhado. Declarou completa desorientação sobre quais distritos são tranquilos para pedestres: \emph{``Quem vem de fora não sabe nem o que é DP. Se eu puder colocar só o CEP do anúncio e entender o bairro, resolve metade das minhas dúvidas''}. Não conhecia as bases da SSP-SP ou GeoSampa. Atribuiu nota 5.

    \item[Munícipe 9 (Abordado na Estação Santo Amaro):] Comerciário buscando financiamento pelo programa habitacional na Zona Sul. Relatou que a segurança de transporte coletivo é prioritária sobre sistemas privados de alarme. Afirmou que as informações do governo parecem feitas para técnicos e burocratas, não para trabalhadores. Atribuiu nota 4.

    \item[Munícipe 10 (Abordada na Praça Benedito Calixto):] Moradora em busca de locação em bairros mistos (comércio e moradia integrados). Enfatizou que locais com bares e restaurantes na calçada transmitem muito mais segurança do que quarteirões com muros extensos e guaritas. Destacou a relevância de cruzar dados de mobilidade cicloviária e saúde com estatísticas de furto/roubo. Atribuiu nota 5.
\end{description}
	% \includepdf[
	% 	pages=-,
	% 	pagecommand={\thispagestyle{fancy}\fancyhf{}\fancyhead[R]{\thepage}\renewcommand{\headrulewidth}{0pt}},
	% ]{ERS_Prisma_Politico_MVP.pdf}

\end{apendicesenv}
% ---


% ----------------------------------------------------------
% Anexos
% ----------------------------------------------------------

% ---
% Inicia os anexos
% ---
% \begin{anexosenv}

% Imprime uma página indicando o início dos anexos
% \partanexos

% ---
% \chapter{Licença do Software AGPL}

% \includedoc{AGPL.tex}

% \end{anexosenv}

%---------------------------------------------------------------------
% INDICE REMISSIVO
%---------------------------------------------------------------------

% \phantompart

% \printindex

%---------------------------------------------------------------------
% Formulário de Identificação (opcional)
%---------------------------------------------------------------------
% \chapter*[Formulário de Identificação]{Formulário de Identificação}
% \addcontentsline{toc}{chapter}{Exemplo de Formulário de Identificação}
% \label{formulado-identificacao}

% Exemplo de Formulário de Identificação, compatível com o Anexo A (informativo)
% da ABNT NBR 10719:2015. Este formulário não é um anexo. Conforme definido na
% norma, ele é o último elemento pós-textual e opcional do relatório.

% \bigskip

% \begin{tabular}{|p{9cm}|p{5cm}|}
% \hline
% \multicolumn{2}{|c|}{\textbf{\large Dados do Relatório Técnico e/ou científico}}\\
% \hline
% \multirow{4}{10cm}[24pt]{Título e subtítulo}& Classificação de segurança\\
%                    & \\
%                    \cline{2-2}
%                    & No.\\
%                    & \\

% \hline
% Tipo de relatório & Data\\
% \hline
% Título do projeto/programa/plano & No.\\
% \hline
% \multicolumn{2}{|l|}{Autor(es)} \\
% \hline
% \multicolumn{2}{|l|}{Instituição executora e endereço completo} \\
% \hline
% \multicolumn{2}{|l|}{Instituição patrocinadora e endereço completo} \\
% \hline
% \multicolumn{2}{|l|}{Resumo}\\[3cm]
% \hline
% \multicolumn{2}{|l|}{Palavras-chave/descritores}\\
% \hline
% \multicolumn{2}{|l|}{
% Edição \hfill No. de páginas \hfill No. do volume \hfill Nº de classificação \phantom{XXXX}} \\
% \hline
% \multicolumn{2}{|l|}{
% ISSN \hfill \hfill Tiragem \hfill Preço \phantom{XXXXXXXX}} \\
% \hline
% \multicolumn{2}{|l|}{Distribuidor} \\
% \hline
% \multicolumn{2}{|l|}{Observações/notas}\\[3cm]
% \hline
% \end{tabular}

\end{document}
